% Énergies renouvelables et valorisation des déchets — SVTEEHB Tle C — Cours signé OnBuch+
% Programme officiel MINESEC. Compiler avec : tectonic N11-energies-renouvelables-valorisation-dechets.tex
\newcommand{\DOCMATIERE}{SVTEEHB}
\newcommand{\DOCNIVEAU}{Tle C}
\newcommand{\DOCMODULE}{Module IV : La Biotechnologie}
\newcommand{\DOCLECON}{N11}
\newcommand{\DOCTITRE}{Énergies renouvelables et valorisation des déchets}
% OnBuch+ — préambule « cours signés » ENRICHI (compilé avec tectonic / XeLaTeX)
% Système visuel « Pop » d'OnBuch+ : fond crème (jamais blanc), bordures encre
% épaisses, ombres « dures » décalées, titres Archivo Black en capitales.
% Version MATHÉMATIQUES : graphes (pgfplots/tikz), boîtes pédagogiques
% (définition, propriété, méthode, attention, le savais-tu, exercice résolu,
% exercice, TP, bilan), page de garde et signature OnBuch+ ancrée en bas.
%
% Macros attendues dans main.tex AVANT \input{preamble} :
%   \DOCMATIERE  \DOCNIVEAU  \DOCMODULE  \DOCLECON (numéro)  \DOCTITRE
\documentclass[11pt]{article}

\usepackage{fontspec}
\usepackage[french]{babel}
\usepackage[a4paper,margin=2cm,top=3.4cm,bottom=2.8cm,headheight=1.4cm,footskip=1.3cm]{geometry}
\usepackage{amsmath,amssymb,amsfonts}
\usepackage{mathtools} % \xmapsto, \coloneqq... souvent utilisés par les modèles
\usepackage{cancel} % simplifications barrées (bilans de forces)
% \abs{z} (module, valeur absolue) et \norm{u} (norme), utilisés par les modèles
\providecommand{\abs}{}\let\abs\relax\DeclarePairedDelimiter{\abs}{\lvert}{\rvert}
\providecommand{\norm}{}\let\norm\relax\DeclarePairedDelimiter{\norm}{\lVert}{\rVert}
\usepackage[table]{xcolor} % [table] : fournit aussi \rowcolors (alternance de lignes)
\usepackage{pagecolor}
\usepackage{tikz}
\usetikzlibrary{arrows.meta,positioning,calc,shapes.geometric,shapes.misc,shapes.arrows,decorations.pathmorphing,decorations.pathreplacing,decorations.markings,patterns,shadows,mindmap,trees,fit,backgrounds,matrix,intersections,angles,quotes}
\usepackage{pgfplots}
\usepackage[version=4]{mhchem} % équations nucléaires \ce{^{235}_{92}U}
\usepackage[european,siunitx]{circuitikz} % schémas électriques
\pgfplotsset{compat=1.18}
\usepgfplotslibrary{fillbetween,groupplots} % aires entre courbes (calcul intégral)
\usepackage{enumitem}
\usepackage{fancyhdr}
% [most] charge toutes les bibliothèques tcolorbox (listings, minted, raster,
% documentation...) et gonfle inutilement la mémoire TeX sur un document long
% (~300 boîtes) : on ne charge que ce qui est réellement utilisé.
\usepackage[skins,breakable]{tcolorbox}
\usepackage{graphicx}
\usepackage{colortbl}
\usepackage{booktabs}
\usepackage{tabularx}
\usepackage{multirow}
\usepackage{diagbox} % case d'en-tête diagonale des tableaux à double entrée
% Colonnes extensibles centrée / gauche / droite (C, L, R), souvent utilisées par les modèles
\newcolumntype{C}{>{\centering\arraybackslash}X}
\newcolumntype{L}{>{\raggedright\arraybackslash}X}
\newcolumntype{R}{>{\raggedleft\arraybackslash}X}
\usepackage{array}
\usepackage{multicol}
\usepackage{siunitx}
% parse-numbers=false : accepte aussi \SI{2,5 \times 10^{-3}}{...}
\sisetup{locale=FR,output-decimal-marker={,},group-minimum-digits=4,parse-numbers=false}
\DeclareSIUnit\ohm{\ensuremath{\symup{\Omega}}} % U+2126 absent de la police
\DeclareSIUnit\division{div} % réglages d'oscilloscope (ms/div, V/div)
\DeclareSIUnit\tr{tr} % tours (vitesse de rotation, ex. \SI{1500}{\tr\per\minute})
\DeclareSIUnit\molar{M} % concentration molaire (ex. \SI{3}{\molar}, \SI{1}{\micro\molar})
\DeclareSIUnit\an{an} % année (le modèle confond parfois avec \year, un compteur TeX) — ex. \SI{5}{\kilo\meter\per\an}
\DeclareSIUnit\FCFA{FCFA} % franc CFA (ex. \SI{500}{\FCFA\per\kilo\gram})
\DeclareSIUnit\degreeBrix{\textdegree Brix} % teneur en sucre d'un sirop/jus (réfractométrie)
\DeclareSIUnit\month{mois} % le modèle utilise parfois \month, un compteur TeX, comme unité de durée
\DeclareSIUnit\microliter{\micro\liter} % le modèle écrit parfois \microliter en un seul token
\DeclareSIUnit\million{millions} % dénombrement (ex. \SI{15}{\million\per\mL} spermatozoïdes)
\usepackage{qrcode}
% \degree : commande fréquemment utilisée par les modèles pour un angle ou une
% température en dehors de \SI{}{} (non fournie par les paquets chargés).
\providecommand{\degree}{^{\circ}}
% \qed (fin de démonstration), utilisé par les modèles sans amsthm
\providecommand{\qed}{\hfill$\square$}
% \barre : double barre des valeurs interdites dans un tableau de variations (tkz-tab)
\providecommand{\barre}{\ensuremath{\|}}
% environnement proof (amsthm non chargé) utilisé par les modèles
\ifdefined\proof\else\newenvironment{proof}[1][Démonstration]{\par\noindent\textit{#1.}\ }{\qed\par}\fi
\usepackage{lastpage}
\usepackage{hyperref}
\hypersetup{hidelinks}

% ============================================================
% Palette « Pop » OnBuch+ — mêmes hex que la classe P (plus_ui.dart)
% ============================================================
\definecolor{poporange}{HTML}{F59321}
\definecolor{popdark}{HTML}{E07A0C}
\definecolor{popcream}{HTML}{FFF6E8}
\definecolor{popink}{HTML}{1C1714}
\definecolor{poppurple}{HTML}{7A5AE0}
\definecolor{popgreen}{HTML}{2E9E6A}
\definecolor{poppink}{HTML}{E0457B}
\definecolor{popblue}{HTML}{2F7DE1}
\definecolor{popgold}{HTML}{C98A12}
\definecolor{popmuted}{HTML}{8A7F76}
\definecolor{popline}{HTML}{EADFD2}
\definecolor{popgray}{HTML}{8A7F76}
\colorlet{popgrey}{popgray}
% Alias tolérants (le modèle nomme parfois les couleurs différemment)
\colorlet{popwhite}{white}
\colorlet{popblack}{popink}
\colorlet{popred}{poppink}
\colorlet{popyellow}{popgold}
% Teintes claires (fonds de tableaux/encadrés) — le modèle nomme parfois
% les couleurs avec un suffixe L (ex. popblueL) sans les avoir définies.
\colorlet{poporangeL}{poporange!20}
\colorlet{popdarkL}{popdark!20}
\colorlet{poppurpleL}{poppurple!20}
\colorlet{popgreenL}{popgreen!20}
\colorlet{poppinkL}{poppink!20}
\colorlet{popblueL}{popblue!20}
\colorlet{popgoldL}{popgold!20}
\colorlet{popredL}{popred!20}
\colorlet{popgrayL}{popgray!20}
% Teintes claires pour fonds de tableaux / zones
\colorlet{popblueL}{popblue!10!white}
\colorlet{poporangeL}{poporange!14!white}
\colorlet{popgreenL}{popgreen!12!white}
\colorlet{poppurpleL}{poppurple!12!white}
\colorlet{poppinkL}{poppink!12!white}
\colorlet{popgoldL}{popgold!14!white}

\pagecolor{popcream}

% ============================================================
% Typographie
% ============================================================
\defaultfontfeatures{Ligatures=TeX}
\setmainfont{PlusJakartaSans-Medium.ttf}[Path=../../../fonts/,BoldFont=PlusJakartaSans-Bold.ttf]
% Maths en sans-serif (Fira Math) avec chiffres et lettres droites de Plus
% Jakarta, pour que les formules se fondent dans le texte.
\usepackage[math-style=ISO,bold-style=ISO]{unicode-math}
\setmathfont{FiraMath-Regular.otf}[Path=../../../fonts/]
\setmathfont{PlusJakartaSans-Medium.ttf}[Path=../../../fonts/,range={up/num,up/{Latin,latin},\mathup}]
\usepackage{icomma} % virgule décimale sans espace en mode math
% Caractères mathématiques Unicode absents des polices de texte (Plus Jakarta,
% Archivo Black) : rendus par leur équivalent mathématique, en texte comme en
% formule — sinon ils s'affichent en carré vide (ex. « Arithmétique dans ℤ »).
\usepackage{newunicodechar}
\newunicodechar{ℝ}{\ensuremath{\mathbb{R}}}
\newunicodechar{ℤ}{\ensuremath{\mathbb{Z}}}
\newunicodechar{ℂ}{\ensuremath{\mathbb{C}}}
\newunicodechar{ℕ}{\ensuremath{\mathbb{N}}}
\newunicodechar{ℚ}{\ensuremath{\mathbb{Q}}}
\newunicodechar{𝔻}{\ensuremath{\mathbb{D}}}
\newunicodechar{α}{\ensuremath{\alpha}}
\newunicodechar{β}{\ensuremath{\beta}}
\newunicodechar{γ}{\ensuremath{\gamma}}
\newunicodechar{δ}{\ensuremath{\delta}}
\newunicodechar{ε}{\ensuremath{\varepsilon}}
\newunicodechar{θ}{\ensuremath{\theta}}
\newunicodechar{λ}{\ensuremath{\lambda}}
\newunicodechar{μ}{\ensuremath{\mu}}
\newunicodechar{σ}{\ensuremath{\sigma}}
\newunicodechar{τ}{\ensuremath{\tau}}
\newunicodechar{φ}{\ensuremath{\varphi}}
\newunicodechar{ψ}{\ensuremath{\psi}}
\newunicodechar{ω}{\ensuremath{\omega}}
\newunicodechar{Σ}{\ensuremath{\Sigma}}
% majuscules produites par \MakeUppercase dans les titres (α-aminés -> Α)
\newunicodechar{Α}{\ensuremath{\alpha}}
\newunicodechar{Β}{\ensuremath{\beta}}
\newunicodechar{Γ}{\ensuremath{\Gamma}}
\newunicodechar{Δ}{\ensuremath{\Delta}}
\newunicodechar{Ε}{\ensuremath{\varepsilon}}
\newunicodechar{Θ}{\ensuremath{\Theta}}
\newunicodechar{Λ}{\ensuremath{\Lambda}}
\newunicodechar{Μ}{\ensuremath{\mu}}
\newunicodechar{Τ}{\ensuremath{\tau}}
\newunicodechar{Φ}{\ensuremath{\Phi}}
\newunicodechar{Ψ}{\ensuremath{\Psi}}
\newunicodechar{Ω}{\ensuremath{\Omega}}
\newunicodechar{↦}{\ensuremath{\mapsto}}
\newunicodechar{∈}{\ensuremath{\in}}
\newunicodechar{∉}{\ensuremath{\notin}}
\newunicodechar{∧}{\ensuremath{\wedge}}
\newunicodechar{⇔}{\ensuremath{\Leftrightarrow}}
\newunicodechar{⟺}{\ensuremath{\Longleftrightarrow}}
\newunicodechar{⇒}{\ensuremath{\Rightarrow}}
\newunicodechar{⟹}{\ensuremath{\Longrightarrow}}
\newunicodechar{∘}{\ensuremath{\circ}}
\newunicodechar{∪}{\ensuremath{\cup}}
\newunicodechar{∩}{\ensuremath{\cap}}
\newunicodechar{≡}{\ensuremath{\equiv}}
\newunicodechar{⊂}{\ensuremath{\subset}}
\newunicodechar{⊥}{\ensuremath{\perp}}
\newunicodechar{∃}{\ensuremath{\exists}}
\newunicodechar{∀}{\ensuremath{\forall}}
\newunicodechar{∛}{\ensuremath{\sqrt[3]{\,}}}
\newunicodechar{∜}{\ensuremath{\sqrt[4]{\,}}}
\newunicodechar{⃗}{\ensuremath{{}^{\to}}}
\newunicodechar{ⁿ}{\ifmmode^{n}\else\textsuperscript{\ensuremath{n}}\fi}
\newunicodechar{ˣ}{\ifmmode^{x}\else\textsuperscript{\ensuremath{x}}\fi}
\newunicodechar{ᵇ}{\ifmmode^{b}\else\textsuperscript{\ensuremath{b}}\fi}
\newunicodechar{ʳ}{\ifmmode^{r}\else\textsuperscript{\ensuremath{r}}\fi}
\newunicodechar{ᵘ}{\ifmmode^{u}\else\textsuperscript{\ensuremath{u}}\fi}
\newunicodechar{ʸ}{\ifmmode^{y}\else\textsuperscript{\ensuremath{y}}\fi}
\newunicodechar{ᵃ}{\ifmmode^{a}\else\textsuperscript{\ensuremath{a}}\fi}
\newunicodechar{ᵉ}{\ifmmode^{e}\else\textsuperscript{\ensuremath{e}}\fi}
\newunicodechar{ᵏ}{\ifmmode^{k}\else\textsuperscript{\ensuremath{k}}\fi}
\newunicodechar{ᵖ}{\ifmmode^{p}\else\textsuperscript{\ensuremath{p}}\fi}
\newunicodechar{ᴺ}{\ifmmode^{N}\else\textsuperscript{\ensuremath{N}}\fi}
\newunicodechar{ᶜ}{\ifmmode^{c}\else\textsuperscript{\ensuremath{c}}\fi}
\newunicodechar{ʰ}{\ifmmode^{h}\else\textsuperscript{\ensuremath{h}}\fi}
\newunicodechar{ᵠ}{\ifmmode^{q}\else\textsuperscript{\ensuremath{q}}\fi}
\newunicodechar{ₙ}{\ifmmode_{n}\else\textsubscript{\ensuremath{n}}\fi}
\newunicodechar{ₐ}{\ifmmode_{a}\else\textsubscript{\ensuremath{a}}\fi}
\newunicodechar{ᵢ}{\ifmmode_{i}\else\textsubscript{\ensuremath{i}}\fi}
\newunicodechar{ᵣ}{\ifmmode_{r}\else\textsubscript{\ensuremath{r}}\fi}
\newunicodechar{ₖ}{\ifmmode_{k}\else\textsubscript{\ensuremath{k}}\fi}
\newunicodechar{ᵦ}{\ifmmode_{\beta}\else\textsubscript{\ensuremath{\beta}}\fi}
\newunicodechar{ₓ}{\ifmmode_{x}\else\textsubscript{\ensuremath{x}}\fi}
\newfontfamily\popdisplay{ArchivoBlack-Regular.ttf}[Path=../../../fonts/]
\newfontfamily\poplogo{SpaceGrotesk-Bold.ttf}[Path=../../../fonts/]
\newfontfamily\popsemi{PlusJakartaSans-SemiBold.ttf}[Path=../../../fonts/]
\newfontfamily\popxbold{PlusJakartaSans-ExtraBold.ttf}[Path=../../../fonts/]
\newfontfamily\popxboldls{PlusJakartaSans-ExtraBold.ttf}[Path=../../../fonts/,LetterSpace=3.0]

\setlist[enumerate]{leftmargin=1.7em,itemsep=5pt,topsep=4pt}
\setlist[itemize]{leftmargin=1.7em,itemsep=4pt,topsep=4pt,label={\color{poporange}\textbullet}}
\setlength{\parindent}{0pt}
\setlength{\parskip}{5pt}
\renewcommand{\arraystretch}{1.25}

% pgfplots : style Pop par défaut
\pgfplotsset{
  every axis/.append style={axis line style={popink,line width=1pt},tick style={popink},
    grid style={popline,line width=0.5pt},label style={font=\small},tick label style={font=\footnotesize}},
  popcurve/.style={poporange,line width=1.8pt,no markers,smooth},
}
% Même style disponible dans un \draw TikZ simple (hors pgfplots)
\tikzset{popcurve/.style={poporange,line width=1.8pt}}
\tikzset{popgrid/.style={popline,very thin}}
\colorlet{popcurve}{poporange} % le nom du style est parfois utilisé comme couleur

% ============================================================
% Pill / squiggle
% ============================================================
\newcommand{\poppill}[3]{%
  \tikz[baseline=-0.55ex]\node[draw=popink,line width=1.2pt,rounded corners=2.6pt,%
    fill=#2,inner xsep=6.5pt,inner ysep=3pt]{%
    \popxboldls\fontsize{7}{7}\selectfont\color{#3}\MakeUppercase{#1}};%
}
\newcommand{\popsquiggle}[1]{%
  \tikz[baseline=-0.4ex]{
    \draw[#1,line width=2.6pt,line cap=round]
      plot [smooth] coordinates {(0,0.09) (0.34,0.01) (0.68,0.21) (1.02,0.01) (1.36,0.10)};
  }%
}
% Mot-clé surligné (marqueur orange sous le texte)
\newcommand{\cle}[1]{{\popxbold\color{popdark}#1}}

% ============================================================
% En-tête / pied
% ============================================================
\pagestyle{fancy}
\fancyhf{}
\renewcommand{\headrulewidth}{0pt}
\renewcommand{\footrulewidth}{0pt}
\lhead{\raisebox{-0.15ex}{\poplogo\fontsize{13}{13}\selectfont\color{popink}OnBuch\color{poporange}+}}
\rhead{\poppill{\DOCMATIERE\ \textbullet\ \DOCNIVEAU\ \textbullet\ Leçon \DOCLECON}{white}{popink}}
\lfoot{\popxbold\fontsize{7.6}{7.6}\selectfont\color{popmuted}\MakeUppercase{Cours signé OnBuch+}\ \textbullet\ {\popsemi\DOCTITRE}}
\rfoot{\tikz[baseline=-0.6ex]\node[rounded corners=8.5pt,fill=popink,minimum height=17pt,minimum width=17pt,inner xsep=5pt,inner sep=0]{\popxbold\fontsize{8}{8}\selectfont\color{popcream}\thepage\,/\,\pageref*{LastPage}};}

% ============================================================
% Titres
% ============================================================
\newcommand{\chaptitre}[2]{%
  \begin{tcolorbox}[enhanced,colback=poporange,colframe=popink,boxrule=2.6pt,arc=11pt,
      drop shadow=popink,left=15pt,right=15pt,top=12pt,bottom=11pt,before skip=3pt,after skip=18pt]
    \poppill{#2}{popink}{popcream}\par\vspace{9pt}
    {\raggedright\hyphenpenalty=10000\exhyphenpenalty=10000\popdisplay\fontsize{22}{24}\selectfont\color{popcream}\MakeUppercase{#1}\par}
  \end{tcolorbox}
}
% Section numérotée : pastille orange avec le numéro + titre Archivo
\newcounter{popsec}
\newcommand{\coursec}[1]{%
  \stepcounter{popsec}\setcounter{popsub}{0}%
  \par\vspace{16pt}\noindent
  \begin{minipage}{\linewidth}
  \tikz[baseline=-0.7ex]\node[draw=popink,line width=1.6pt,fill=poporange,rounded corners=4pt,minimum size=22pt,inner sep=0]{\popdisplay\fontsize{12}{12}\selectfont\color{popcream}\thepopsec};%
  \hspace{8pt}{\popdisplay\fontsize{15}{17}\selectfont\color{popink}\MakeUppercase{#1}}\par
  \vspace{4pt}\noindent{\color{poporange}\rule{36pt}{3.6pt}}
  \end{minipage}\par\nopagebreak\vspace{8pt}
}
\newcounter{popsub}
\newcommand{\courssub}[1]{%
  \stepcounter{popsub}%
  \par\vspace{9pt}\noindent{\popxbold\fontsize{11.5}{13}\selectfont\color{popink}{\color{poporange}\thepopsec.\thepopsub}\hspace{6pt}#1}\par\nopagebreak\vspace{4pt}
}

% ============================================================
% Boîtes pédagogiques — même recette : fond blanc, bordure épaisse colorée,
% ombre dure encre, badge plein. Titre optionnel [#1].
% ============================================================
\tcbset{popbase/.style={breakable,enhanced,colback=white,boxrule=2.3pt,arc=9pt,
  shadow={3pt}{-3pt}{0pt}{popink},left=14pt,right=14pt,top=11pt,bottom=11pt,before skip=11pt,after skip=15pt,
  fontupper=\popsemi\fontsize{10.6}{14.5}\selectfont\color{popink},
  fontlower=\popsemi\fontsize{10.6}{14.5}\selectfont\color{popink}}}
% Coche dessinée (le glyphe ✓ n'existe pas dans les polices du document)
\newcommand{\popcheck}[1]{\tikz[baseline=-0.5ex]\draw[#1,line width=1.6pt,line cap=round,line join=round](-0.13,0.0)--(-0.04,-0.09)--(0.13,0.1);}
\providecommand{\coche}{\popcheck{popgreen}}
\providecommand{\resultat}[1]{\par\smallskip\noindent\fcolorbox{popgreen}{popgreenL}{\parbox{\dimexpr\linewidth-2\fboxsep-2\fboxrule\relax}{\textbf{Résultat.} #1}}\par\smallskip}
\newcommand{\popboxhead}[3]{\poppill{#1}{#2}{white}\ifx\relax#3\relax\else\hspace{6pt}{\popxbold\fontsize{10.6}{12}\selectfont\color{popink}#3}\fi\par\vspace{7pt}}

\newtcolorbox{aretenir}[1][]{popbase,colframe=popgreen,before upper={\popboxhead{À retenir}{popgreen}{#1}}}
\newtcolorbox{exemplebox}[1][]{popbase,colframe=poporange,before upper={\popboxhead{Exemple}{poporange}{#1}}}
\newtcolorbox{definition}[1][]{popbase,colframe=popblue,before upper={\popboxhead{Définition}{popblue}{#1}}}
\newtcolorbox{propriete}[1][]{popbase,colframe=poppurple,before upper={\popboxhead{Propriété}{poppurple}{#1}}}
\newtcolorbox{methode}[1][]{popbase,colframe=popgold,before upper={\popboxhead{Méthode}{popgold}{#1}}}
\newtcolorbox{attention}[1][]{popbase,colframe=poppink,before upper={\popboxhead{Attention}{poppink}{#1}}}
\newtcolorbox{experience}[1][]{popbase,colframe=popblue,colback=popblueL,before upper={\popboxhead{Expérience}{popblue}{#1}}}
% « Le savais-tu ? » : carte violette pleine (contexte, culture, Cameroun)
\newtcolorbox{savaistu}[1][]{popbase,colback=poppurple,colframe=popink,
  fontupper=\popsemi\fontsize{10.4}{14.2}\selectfont\color{white},
  before upper={\poppill{Le savais-tu\,?}{white}{poppurple}\ifx\relax#1\relax\else\hspace{6pt}{\popxbold\color{white}#1}\fi\par\vspace{7pt}}}
% Exercice résolu : énoncé en haut, \tcblower puis la solution sur fond crème
\newcounter{popexo}\newcounter{popexr}
\newtcolorbox{exoresolu}[1][]{popbase,colframe=poporange,colbacklower=popcream,
  segmentation style={popink,line width=1.2pt,dashed},
  before upper={\stepcounter{popexr}\popboxhead{Exercice résolu \thepopexr}{poporange}{#1}},
  before lower={\poppill{Solution}{popink}{popcream}\par\vspace{6pt}}}
% Exercice (non résolu en place) — la correction est dans la partie Corrigés
\newtcolorbox{exercice}[1][]{popbase,colframe=popink,
  before upper={\stepcounter{popexo}\popboxhead{Exercice \thepopexo}{popink}{#1}}}
% Corrigé
\newtcolorbox{corrige}[1][]{popbase,colframe=popgreen,colback=popgreenL,
  before upper={\popboxhead{Corrigé}{popgreen}{#1}}}
% Objectifs / prérequis / bilan
\newtcolorbox{objectifs}{popbase,colframe=popink,colback=poporangeL,
  before upper={\popboxhead{Objectifs de la leçon}{popink}{}}}
\newtcolorbox{prerequis}{popbase,colframe=popink,colback=popblueL,
  before upper={\popboxhead{Prérequis}{popblue}{}}}
\newtcolorbox{bilan}{popbase,colframe=popink,colback=white,boxrule=2.8pt,
  before upper={\popboxhead{Fiche bilan}{poporange}{L'essentiel à connaître pour l'examen}}}
% Figure encadrée avec légende
\newtcolorbox{popfigure}[1][]{enhanced,breakable=false,colback=white,colframe=popline,boxrule=1.4pt,arc=8pt,
  left=10pt,right=10pt,top=10pt,bottom=8pt,before skip=10pt,after skip=12pt,halign=center,
  lower separated=false,halign lower=center,
  fontlower=\popsemi\fontsize{9}{11.5}\selectfont\color{popmuted},
  #1}
\newcommand{\legende}[1]{\par\vspace{6pt}{\popsemi\fontsize{9}{11.5}\selectfont\color{popmuted}#1}}

% ============================================================
% Signature OnBuch+
% ============================================================
\newcommand{\poplogoinline}[1]{{\poplogo\fontsize{#1}{#1}\selectfont\color{popink}OnBuch\color{poporange}+}}
\newcommand{\signatureonbuch}{%
  \begin{tcolorbox}[enhanced,colback=popink,colframe=popink,boxrule=2.2pt,arc=10pt,
      drop shadow=poporange,left=14pt,right=14pt,top=10pt,bottom=10pt,before skip=0pt,after skip=0pt]
    \tikz[baseline=-0.7ex]\node[circle,fill=popgreen,draw=popcream,line width=1.3pt,minimum size=22pt,inner sep=0]{\popcheck{white}};%
    \hspace{9pt}\begin{minipage}[c]{0.62\linewidth}
      {\popxbold\fontsize{12}{13}\selectfont\color{popcream}Signé\ \textbullet\ {\poplogo OnBuch\color{poporange}+}}\par\vspace{2pt}
      {\raggedright\popsemi\fontsize{8}{10.5}\selectfont\color{popline}\MakeUppercase{Équipe pédagogique OnBuch+}\\\MakeUppercase{Conforme au programme officiel MINESEC}\par}
    \end{minipage}\hfill
    {\poplogo\fontsize{22}{22}\selectfont\color{popcream}OnBuch\color{poporange}+}
  \end{tcolorbox}%
}

% ============================================================
% Page de garde
% #1 titre · #2 matière · #3 niveau · #4 module · #5 n° leçon · #6 durée ·
% #7 description / objectifs (texte ou itemize)
% ============================================================
\newcommand{\pagegarde}[7]{%
  \thispagestyle{empty}
  \newgeometry{a4paper,margin=1.7cm,top=1.6cm,bottom=1.4cm}%
  \begin{tikzpicture}[remember picture,overlay]
    \fill[poporange!18] ([xshift=-3.2cm,yshift=-3.6cm]current page.north east) circle (4.6cm);
    \fill[poppurple!14] ([xshift=2.4cm,yshift=7.6cm]current page.south west) circle (3.8cm);
  \end{tikzpicture}%
  {\poplogo\fontsize{30}{30}\selectfont\color{popink}OnBuch\color{poporange}+}\hfill
  \raisebox{0.6ex}{\poppill{Cours signé \textbullet\ Premium}{popink}{popcream}}\par
  \vspace{1pt}\popsquiggle{poppurple}\par
  \vspace{8pt}
  \begin{tcolorbox}[enhanced,colback=poporange,colframe=popink,boxrule=2.8pt,arc=13pt,
      drop shadow=popink,left=18pt,right=18pt,top=12pt,bottom=13pt,after skip=10pt]
    \poppill{#2 \textbullet\ #3}{popink}{popcream}\hspace{5pt}\poppill{#4}{white}{popink}\par\vspace{8pt}
    {\popdisplay\fontsize{14}{14}\selectfont\color{popink}LEÇON #5}\par\vspace{4pt}
    {\raggedright\hyphenpenalty=10000\exhyphenpenalty=10000\popdisplay\fontsize{25}{27}\selectfont\color{popcream}\MakeUppercase{#1}\par}
    \vspace{8pt}
    {\popxbold\fontsize{9.5}{9.5}\selectfont\color{popink}\MakeUppercase{Durée conseillée : #6}}
  \end{tcolorbox}
  \begin{tcolorbox}[enhanced,colback=white,colframe=popink,boxrule=2.2pt,arc=10pt,
      drop shadow=popink,left=16pt,right=16pt,top=10pt,bottom=10pt,after skip=8pt]
    \poppill{Dans cette leçon}{poppurple}{white}\par\vspace{5pt}
    \popsemi\fontsize{9.3}{12.6}\selectfont\color{popink}#7
  \end{tcolorbox}
  \noindent\begin{minipage}[t]{0.487\linewidth}
    \begin{tcolorbox}[enhanced,colback=popgreen,colframe=popink,boxrule=2.2pt,arc=10pt,
        drop shadow=popink,left=11pt,right=11pt,top=10pt,bottom=10pt,height=3.1cm,valign=center]
      \tcbox[colback=white,colframe=popink,boxrule=1.3pt,arc=5pt,boxsep=0pt,nobeforeafter,
        left=3pt,right=3pt,top=3pt,bottom=3pt,valign=center]{\qrcode[height=1.9cm]{https://play.google.com/store/apps/details?id=com.luvvix.onbuch&hl=fr}}%
      \hspace{8pt}\begin{minipage}[c]{0.5\linewidth}
        {\popdisplay\fontsize{11}{12.5}\selectfont\color{white}\MakeUppercase{Télécharge OnBuch}}\par\vspace{4pt}
        {\raggedright\popsemi\fontsize{8.4}{11}\selectfont\color{white}Gratuit \textbullet\ Google Play\\Tuteur IA, annales, concours, résultats\par}
      \end{minipage}
    \end{tcolorbox}
  \end{minipage}\hfill
  \begin{minipage}[t]{0.487\linewidth}
    \begin{tcolorbox}[enhanced,colback=poppurple,colframe=popink,boxrule=2.2pt,arc=10pt,
        drop shadow=popink,left=11pt,right=11pt,top=10pt,bottom=10pt,height=3.1cm,valign=center]
      \tikz[baseline=-0.7ex]\node[circle,fill=white,draw=popink,line width=1.3pt,minimum size=1.6cm,inner sep=0]{\popdisplay\fontsize{24}{24}\selectfont\color{poppurple}+};%
      \hspace{8pt}\begin{minipage}[c]{0.6\linewidth}
        {\popdisplay\fontsize{11}{12.5}\selectfont\color{white}\MakeUppercase{Passe à OnBuch+}}\par\vspace{4pt}
        {\raggedright\popsemi\fontsize{8.4}{11}\selectfont\color{white}Tous les cours signés, Tuteur IA illimité, fiches PDF — onglet OnBuch+ de l'app\par}
      \end{minipage}
    \end{tcolorbox}
  \end{minipage}
  \vspace{8pt}
  \popcontenu
  \vfill
  \signatureonbuch
  \restoregeometry
  \newpage
}

% Rangée « Ce cours contient » (page de garde)
\newcommand{\poptile}[3]{% #1 couleur #2 gros texte #3 légende
  \begin{tcolorbox}[enhanced,colback=#1,colframe=popink,boxrule=2pt,arc=9pt,shadow={2.5pt}{-2.5pt}{0pt}{popink},
      left=6pt,right=6pt,top=9pt,bottom=9pt,halign=center,valign=center,height=1.75cm,width=\linewidth,nobeforeafter]
    {\popdisplay\fontsize{12.5}{13}\selectfont\color{white}\MakeUppercase{#2}}\par\vspace{3pt}
    {\centering\popsemi\fontsize{7.8}{9.5}\selectfont\color{white}#3\par}
  \end{tcolorbox}}
\newcommand{\popcontenu}{%
  \par\noindent{\popxboldls\fontsize{7.5}{7.5}\selectfont\color{popmuted}CE COURS SIGNÉ CONTIENT}\par\vspace{7pt}
  \noindent\begin{minipage}[t]{0.235\linewidth}\poptile{popblue}{Cours}{complet, illustré, pas à pas}\end{minipage}\hfill
  \begin{minipage}[t]{0.235\linewidth}\poptile{poporange}{Méthodes}{et exercices résolus}\end{minipage}\hfill
  \begin{minipage}[t]{0.235\linewidth}\poptile{poppink}{Exercices}{type examen + corrigés}\end{minipage}\hfill
  \begin{minipage}[t]{0.235\linewidth}\poptile{popgreen}{Fiche bilan}{l'essentiel à retenir}\end{minipage}\par}

% Sommaire
\newtcolorbox{sommaire}{enhanced,colback=white,colframe=popink,boxrule=2.2pt,arc=10pt,
  drop shadow=popink,left=18pt,right=18pt,top=13pt,bottom=13pt,before skip=6pt,after skip=20pt,
  before upper={\poppill{Sommaire}{poppurple}{white}\par\vspace{9pt}\popsemi\fontsize{10.6}{14.5}\selectfont\color{popink}}}

% Signature de fin de cours — poussée tout en bas de la dernière page
\newcommand{\findecours}{%
  \par\vfill\nopagebreak
  \begin{center}\popsquiggle{poporange}\end{center}\vspace{4pt}
  \signatureonbuch
}
\AtBeginDocument{\def\llbracket{\mathopen{[\mkern-3.5mu[}}\def\rrbracket{\mathclose{]\mkern-3.5mu]}}} % intervalles d'entiers (glyphe ⟦ absent de la police)
% Glyphes absents de Fira Math : redéfinis à partir de glyphes disponibles
\AtBeginDocument{%
  \def\setminus{\mathbin{\backslash}}\let\smallsetminus\setminus
  \def\vdots{\mathord{\vbox{\baselineskip3.2pt\lineskiplimit0pt\kern4pt\hbox{.}\hbox{.}\hbox{.}}}}%
  \def\ddots{\mathinner{\mkern1mu\raise6.5pt\hbox{.}\mkern2mu\raise3.6pt\hbox{.}\mkern2mu\raise0.7pt\hbox{.}\mkern1mu}}%
  \def\triangle{\mathord{\scalebox{0.9}{$\symup{\Delta}$}}}%
  \def\nabla{\mathord{\raisebox{\depth}{\scalebox{1}[-1]{$\symup{\Delta}$}}}}%
  \def\checkmark{\popcheck{popdark}}%
  \def\star{\mathbin{\ast}}\def\bigstar{\mathord{\ast}}%
  \def\diamond{\mathbin{\scalebox{0.6}{$\Diamond$}}}%
  \def\bigcup{\mathop{\mathchoice{\raisebox{-0.3em}{\scalebox{1.7}{$\cup$}}}{\scalebox{1.25}{$\cup$}}{\cup}{\cup}}\displaylimits}%
  \def\bigcap{\mathop{\mathchoice{\raisebox{-0.3em}{\scalebox{1.7}{$\cap$}}}{\scalebox{1.25}{$\cap$}}{\cap}{\cap}}\displaylimits}%
  \def\lesseqqgtr{\mathrel{\lessgtr}}\def\bigoplus{\mathop{\oplus}\displaylimits}%
}
\providecommand{\ssi}{si et seulement si} % abréviation fréquente
\AtBeginDocument{\providecommand{\wideparen}{\overparen}} % arc de cercle

%% FIGFIX-BEGIN v3 — correctifs globaux des figures (docs/onbuch-generation/tools/figfix)
\usepackage{adjustbox}\usepackage{etoolbox}
% toute figure TikZ/pgfplots plus large que la ligne est réduite à la largeur disponible
\BeforeBeginEnvironment{tikzpicture}{\begin{adjustbox}{max width=\linewidth}}
\AfterEndEnvironment{tikzpicture}{\end{adjustbox}}
% légendes pgfplots lisibles par-dessus les courbes
\pgfplotsset{every axis/.append style={legend style={fill=white,fill opacity=0.92,text opacity=1,draw=black!15,inner sep=3pt}}}
% étiquettes d'arêtes (syntaxe edge["..."]) avec fond blanc
\tikzset{every edge quotes/.append style={fill=white,inner sep=1.2pt}}
% bulles de cartes mentales : texte à la ligne dans la bulle, bulle un peu plus grande
\tikzset{every concept/.append style={text width=2.0cm,align=center,minimum size=2.3cm,inner sep=2pt}}
%% FIGFIX-END
\begin{document}

\pagegarde{Énergies renouvelables et valorisation des déchets}{SVTEEHB}{Tle C}{Module IV : La Biotechnologie}{N11}{6 h}{Cette leçon présente les techniques de production des biocarburants et les différentes voies de valorisation des déchets (papiers et plastiques) au Cameroun. Elle montre comment le recyclage et la transformation des déchets contribuent à la protection de l'environnement et au développement durable.
\vspace{4pt}
\begin{multicols}{2}\small
\begin{itemize}
\item À la fin de cette leçon, je sais définir les notions d'énergie renouvelable, de biocarburant et de valorisation des déchets
\item À la fin de cette leçon, je sais expliquer les techniques de production des biocarburants (bioéthanol, biodiesel, biogaz)
\item À la fin de cette leçon, je sais expliquer la technique de production de l'énergie à partir du papier
\item À la fin de cette leçon, je sais expliquer les techniques de recyclage du papier en serviettes jetables, en papier hygiénique et en papier-cadeau
\item À la fin de cette leçon, je sais expliquer la technique de recyclage des déchets plastiques en bouteilles recyclées
\item À la fin de cette leçon, je sais expliquer la technique de transformation des déchets plastiques en pavés
\end{itemize}\end{multicols}}

\begin{sommaire}
\begin{multicols}{2}
\begin{enumerate}
\item Généralités : énergies renouvelables et valorisation des déchets
\item Les techniques de production des biocarburants
\item La production d'énergie à partir du papier
\item Le recyclage du papier : serviettes jetables, papier hygiénique, papier-cadeau
\item Le recyclage des déchets plastiques en bouteilles recyclées
\item La transformation des déchets plastiques en pavés et les avantages de la valorisation
\item Sensibilisation de la communauté au recyclage
\item Activité d'intégration
\item Méthodes et exercices résolus
\item Exercices
\item Corrigés des exercices
\item Fiche bilan
\end{enumerate}
\end{multicols}
\end{sommaire}

\begin{objectifs}
\begin{itemize}
\item À la fin de cette leçon, je sais définir les notions d'énergie renouvelable, de biocarburant et de valorisation des déchets
\item À la fin de cette leçon, je sais expliquer les techniques de production des biocarburants (bioéthanol, biodiesel, biogaz)
\item À la fin de cette leçon, je sais expliquer la technique de production de l'énergie à partir du papier
\item À la fin de cette leçon, je sais expliquer les techniques de recyclage du papier en serviettes jetables, en papier hygiénique et en papier-cadeau
\item À la fin de cette leçon, je sais expliquer la technique de recyclage des déchets plastiques en bouteilles recyclées
\item À la fin de cette leçon, je sais expliquer la technique de transformation des déchets plastiques en pavés
\item À la fin de cette leçon, je sais citer les avantages de la valorisation des déchets
\item À la fin de cette leçon, je sais informer la communauté sur la nécessité de transformer et de recycler les déchets
\end{itemize}
\end{objectifs}

\begin{prerequis}
\begin{itemize}
\item Notion de fermentation et de respiration cellulaire (Première)
\item Notion de biotechnologie et d'utilisation des micro-organismes (Module IV, leçons précédentes)
\item Structure de la cellulose et des glucides (Première)
\item Notion de pollution de l'environnement et de développement durable (Seconde/Première)
\end{itemize}
\end{prerequis}

\newpage

\coursec{Généralités : énergies renouvelables et valorisation des déchets}

\textbf{Situation d'accroche.} À Yaoundé comme à Douala, les montagnes de déchets plastiques et de papiers s'accumulent dans les quartiers, tandis que le prix du carburant ne cesse d'augmenter. Pourtant, dans certains quartiers, des jeunes transforment ces mêmes déchets en pavés pour les cours d'écoles et en bouteilles réutilisables, et des unités pilotes produisent du biodiesel à partir d'huile de palme. Comment des déchets qui polluent nos rues peuvent-ils devenir une source d'énergie et de revenus ? Quelles techniques biologiques et industrielles permettent cette transformation ? C'est ce que cette leçon va nous apprendre.

\textbf{Problématique.} Toute société humaine produit des déchets et consomme de l'énergie. Lorsque l'énergie provient de ressources épuisables et que les déchets sont abandonnés dans la nature, il en résulte une double crise : raréfaction des ressources et pollution de l'environnement. La biotechnologie propose une réponse élégante : utiliser la matière vivante (la \cle{biomasse}) et les déchets eux-mêmes comme matières premières. Avant d'étudier les techniques de production des biocarburants et de recyclage, il faut d'abord définir rigoureusement les notions d'énergie renouvelable, de biocarburant et de valorisation des déchets, et comprendre la nature des déchets qui nous entourent.

\courssub{Les énergies renouvelables}

\textbf{Observation et intuition.} Allume un feu de bois dans la cour de la maison : le bois brûle et libère de la chaleur. L'année suivante, de nouveaux arbres ont poussé dans la savane ou la plantation. Compare maintenant avec le gasoil que l'on achète à la station-service : ce pétrole s'est formé il y a des millions d'années à partir d'organismes marins enfouis ; une fois brûlé dans le moteur d'un taxi, il est définitivement perdu, et aucun gisement ne se reformera à l'échelle d'une vie humaine. Cette différence de \emph{vitesse de renouvellement} est le cœur de la distinction entre énergies renouvelables et énergies fossiles.

\begin{definition}[Énergie renouvelable]
Une \cle{énergie renouvelable} est une énergie issue d'une source qui se \textbf{renouvelle naturellement à l'échelle du temps humain} (quelques années à quelques décennies), de sorte que son exploitation n'épuise pas la ressource. Les principales sources sont : le \textbf{soleil} (énergie solaire), la \textbf{biomasse} (bois, résidus agricoles, biocarburants), le \textbf{vent} (énergie éolienne) et l'\textbf{eau} (énergie hydraulique).
\end{definition}

\begin{definition}[Énergie fossile]
Une \cle{énergie fossile} est une énergie issue de la transformation, pendant des \textbf{millions d'années}, de débris d'organismes vivants enfouis dans le sous-sol : le \textbf{pétrole}, le \textbf{charbon} et le \textbf{gaz naturel}. Leur stock est \textbf{limité} et leur combustion libère massivement du dioxyde de carbone \ce{CO2}, principal \textbf{gaz à effet de serre} responsable du réchauffement climatique.
\end{definition}

\textbf{Raisonnement à partir des faits.} Le bilan carbone distingue fondamentalement les deux types d'énergie. Un arbre absorbe du \ce{CO2} atmosphérique par photosynthèse pendant sa croissance ; lorsqu'on brûle son bois (ou un biocarburant issu de plantes), le \ce{CO2} relâché est celui qui avait été prélevé peu de temps auparavant : le cycle est \emph{bouclé} à court terme. À l'inverse, la combustion du pétrole remet dans l'atmosphère un carbone séquestré depuis des millions d'années : c'est un \emph{ajout net} de \ce{CO2}. C'est pourquoi les énergies fossiles aggravent l'effet de serre, alors que la biomasse, gérée durablement, est proche de la neutralité carbone.

\begin{popfigure}
\begin{center}
\begin{tikzpicture}[font=\small, >=Stealth]
% --- Cycle fossile ---
\node[draw=popdark, fill=popdark!15, rounded corners, minimum width=2.6cm, minimum height=1cm, align=center] (gis) at (0,0) {Gisement fossile\\(pétrole, charbon)};
\node[draw=poporange, fill=poporangeL, rounded corners, minimum width=2.6cm, minimum height=1cm, align=center] (usine) at (4.2,0) {Combustion\\(moteur, usine)};
\node[draw=popink, fill=poppinkL, rounded corners, minimum width=2.6cm, minimum height=1cm, align=center] (atm1) at (8.4,0) {\ce{CO2} accumulé\\dans l'atmosphère};
\draw[->, very thick, popdark] (gis) -- (usine) node[midway, above]{extraction};
\draw[->, very thick, poporange] (usine) -- (atm1) node[midway, above]{émission nette};
\draw[->, dashed, popmuted] (atm1.south) .. controls (4.2,-1.6) .. (gis.south) node[midway, below, align=center]{retour impossible : formation en millions d'années};
\node[font=\bfseries, popdark] at (4.2,1.3) {Cycle OUVERT : ressource épuisable};

% --- Cycle renouvelable ---
\node[draw=popgreen, fill=popgreenL, rounded corners, minimum width=2.6cm, minimum height=1cm, align=center] (plante) at (0,-3.6) {Plante\\(biomasse)};
\node[draw=popblue, fill=popblueL, rounded corners, minimum width=2.6cm, minimum height=1cm, align=center] (carbu) at (4.2,-3.6) {Biocarburant\\ou bois énergie};
\node[draw=popink, fill=poppinkL, rounded corners, minimum width=2.6cm, minimum height=1cm, align=center] (atm2) at (8.4,-3.6) {\ce{CO2} atmosphérique};
\draw[->, very thick, popgreen] (plante) -- (carbu) node[midway, above]{récolte, transformation};
\draw[->, very thick, popblue] (carbu) -- (atm2) node[midway, above]{combustion};
\draw[->, very thick, popgreen] (atm2.south) .. controls (4.2,-5.4) .. (plante.south) node[midway, below, align=center]{photosynthèse : le \ce{CO2} est réabsorbé en quelques années};
\node[font=\bfseries, popgreen!60!black] at (4.2,-2.3) {Cycle BOUCLÉ : ressource renouvelable};
\end{tikzpicture}
\end{center}
\legende{Comparaison des cycles de vie d'une énergie fossile (en haut, cycle ouvert : le carbone libéré ne retourne pas au gisement) et d'une énergie renouvelable issue de la biomasse (en bas, cycle bouclé : la photosynthèse réabsorbe le \ce{CO2} émis par la combustion).}
\end{popfigure}

\begin{aretenir}[L'essentiel sur les énergies]
Une énergie est dite \textbf{renouvelable} lorsque sa source se reconstitue naturellement à l'échelle humaine (soleil, vent, eau, biomasse). Les énergies \textbf{fossiles} (pétrole, charbon, gaz) sont \textbf{épuisables} et leur combustion accroît l'effet de serre par émission nette de \ce{CO2}.
\end{aretenir}

\courssub{Les biocarburants : une énergie renouvelable issue du vivant}

\textbf{Observation.} Dans plusieurs localités du Cameroun, des générateurs et des véhicules expérimentaux fonctionnent avec un carburant obtenu à partir d'huile de palme ou d'huile de coton. Ce carburant n'est pas tiré du sous-sol : il est \emph{fabriqué} à partir de plantes cultivées. C'est l'idée générale des biocarburants.

\begin{definition}[Biocarburant]
Un \cle{biocarburant} est un \textbf{carburant produit à partir de la biomasse végétale ou animale}, c'est-à-dire à partir de matière organique d'origine récente (végétaux cultivés, résidus agricoles, huiles usagées, déchets organiques). Exemples : le \textbf{biodiesel} obtenu à partir d'huiles végétales (palme, soja, jatropha) et le \textbf{bioéthanol} obtenu par fermentation alcoolique de sucres ou d'amidon (canne à sucre, maïs, manioc).
\end{definition}

\textbf{Justification biologique.} Les plantes stockent l'énergie solaire sous forme de molécules organiques riches en carbone (glucides, lipides) grâce à la photosynthèse, dont le bilan simplifié est :
\[ \ce{6CO2 + 6H2O ->[lumiere] C6H12O6 + 6O2} \]
La biotechnologie consiste ensuite à \emph{convertir} cette matière organique en carburant liquide utilisable dans les moteurs : fermentation des sucres en éthanol par des levures, ou transformation chimique des huiles végétales en biodiesel. Ces techniques seront détaillées dans la section suivante.

\begin{attention}[Ne pas confondre « recyclable » et « renouvelable »]
C'est l'erreur la plus fréquente au Baccalauréat. \textbf{Renouvelable} qualifie une \emph{source d'énergie} qui se reconstitue naturellement à l'échelle humaine (soleil, biomasse). \textbf{Recyclable} qualifie un \emph{matériau} que l'on peut transformer pour le réutiliser (papier, plastique, métal). Un plastique recyclé n'est \textbf{pas} une énergie renouvelable : il est issu du pétrole, une ressource fossile, et son recyclage ne fait qu'économiser la matière première, sans créer d'énergie nouvelle. De même, le papier recyclé est un matériau valorisé, pas une source d'énergie renouvelable en soi.
\end{attention}

\courssub{La valorisation des déchets}

\textbf{Observation.} Dans une décharge sauvage de la périphérie de Douala, on trouve pêle-mêle des restes de nourriture, des cartons, des sachets plastiques et des boîtes de conserve. Un maraîcher voisin récupère pourtant les déchets organiques pour fabriquer du compost, et une coopérative collecte les plastiques pour les revendre à une unité de transformation. Ce qui était une nuisance devient une \emph{ressource} : c'est exactement la logique de la valorisation.

\begin{definition}[Valorisation des déchets]
La \cle{valorisation des déchets} est l'\textbf{ensemble des techniques permettant de transformer un déchet en produit utile ou en énergie}. On distingue :
\begin{itemize}
\item la \textbf{valorisation matière} (recyclage) : le déchet devient une nouvelle matière première (papier recyclé en papier hygiénique, plastique fondu en pavés) ;
\item la \textbf{valorisation organique} : les déchets organiques sont transformés en \textbf{compost} ou en \textbf{biogaz} (méthanisation par des micro-organismes) ;
\item la \textbf{valorisation énergétique} : le déchet est brûlé ou transformé pour produire de la chaleur ou de l'électricité (incinération avec récupération d'énergie, production de combustible à partir du papier).
\end{itemize}
\end{definition}

\textbf{Typologie des déchets de l'environnement de l'Homme.} Pour valoriser, il faut d'abord \emph{trier}, car chaque catégorie de déchet exige une technique différente. On distingue classiquement :

\begin{center}
\begin{tabularx}{\linewidth}{|l|X|X|}\hline
\rowcolor{popblueL} \textbf{Type de déchet} & \textbf{Exemples} & \textbf{Voie de valorisation} \\ \hline
Déchets organiques & épluchures, restes de repas, fumier, herbes coupées & compost, biogaz (méthanisation) \\ \hline
Papiers et cartons & cahiers usagés, journaux, emballages en carton & recyclage en papier hygiénique, serviettes, papier-cadeau ; combustible \\ \hline
Plastiques & sachets, bouteilles, bidons & recyclage en bouteilles, transformation en pavés \\ \hline
Métaux & boîtes de conserve, ferraille, aluminium & fonte et refabrication d'objets métalliques \\ \hline
\end{tabularx}
\end{center}

\begin{popfigure}
\begin{center}
\begin{tikzpicture}[font=\small]
% Diagramme circulaire : parts cumulées
% Organiques 50%, Papiers 15%, Plastiques 15%, Métaux 8%, Autres 12%
\def\r{2.6}
\fill[popgreenL]  (0,0) -- (90:\r) arc (90:-90:\r) -- cycle;          % 50%
\fill[popblueL]   (0,0) -- (-90:\r) arc (-90:-144:\r) -- cycle;       % 15%
\fill[poporangeL] (0,0) -- (-144:\r) arc (-144:-198:\r) -- cycle;     % 15%
\fill[poppurpleL] (0,0) -- (-198:\r) arc (-198:-226.8:\r) -- cycle;   % 8%
\fill[popgoldL]   (0,0) -- (-226.8:\r) arc (-226.8:-270:\r) -- cycle; % 12%
\draw[popdark, thick] (0,0) circle (\r);
% Étiquettes
\node at (0:1.5) {\textbf{50 \,\%}};
\node[align=center] at (0:3.6) {Déchets\\organiques};
\node at (-117:1.6) {\textbf{15 \,\%}};
\node[align=center] at (-117:3.7) {Papiers\\cartons};
\node at (-171:1.7) {\textbf{15 \,\%}};
\node[align=center] at (-171:3.7) {Plastiques};
\node at (-212:1.9) {\textbf{8 \,\%}};
\node[align=center] at (-212:3.9) {Métaux};
\node at (-248:1.8) {\textbf{12 \,\%}};
\node[align=center] at (-248:3.8) {Autres\\(verre, textile)};
% Liaisons
\draw[popmuted] (0:2.6) -- (0:3.0);
\draw[popmuted] (-117:2.6) -- (-117:3.1);
\draw[popmuted] (-171:2.6) -- (-171:3.1);
\draw[popmuted] (-212:2.6) -- (-212:3.2);
\draw[popmuted] (-248:2.6) -- (-248:3.1);
\end{tikzpicture}
\end{center}
\legende{Répartition approximative des types de déchets ménagers produits dans une grande ville camerounaise (ordres de grandeur). Les déchets organiques dominent largement (environ la moitié), ce qui explique le potentiel considérable du compostage et de la méthanisation ; les papiers et plastiques, bien que minoritaires en masse, posent les problèmes de pollution les plus visibles.}
\end{popfigure}

\textbf{Lecture du document.} Ce diagramme montre que près de la moitié des déchets ménagers sont \emph{organiques} : ils pourraient nourrir les sols (compost) ou produire du gaz (biogaz) au lieu d'encombrer les rues. Les papiers et plastiques représentent chacun environ 15 \,\% de la masse, mais leur persistance dans la nature en fait les polluants les plus visibles et les plus durables.

\begin{exemplebox}[Durées de dégradation : un contraste frappant]
Tous les déchets ne se dégradent pas à la même vitesse dans la nature :
\begin{itemize}
\item une \textbf{feuille de papier} abandonnée au sol est décomposée par les micro-organismes du sol en \textbf{quelques mois} (2 à 6 mois environ) ;
\item un \textbf{sac plastique} (sachet) persiste pendant \textbf{plusieurs centaines d'années} (de 100 à 450 ans selon les estimations), se fragmentant en microplastiques qui contaminent les sols et les cours d'eau.
\end{itemize}
Cette différence explique pourquoi les plastiques s'accumulent dans l'environnement et pourquoi leur \emph{recyclage} est une priorité, alors que le papier peut aussi être valorisé comme combustible ou composté.
\end{exemplebox}

\courssub{La gestion des déchets au Cameroun : une problématique nationale}

\textbf{Document et situation concrète.} Au Cameroun, la collecte des ordures ménagères dans les grandes villes est assurée principalement par la société \textbf{HYSACAM} (Hygiène et Salubrité du Cameroun). Malgré son action, une part importante des déchets échappe à la collecte : faute de bennes suffisantes ou de desserte de certains quartiers, les ordures sont abandonnées dans des \textbf{décharges sauvages}, le long des routes, dans les ravins et les cours d'eau. Conséquences directes : insalubrité, prolifération des moustiques et des mouches (vecteurs du paludisme et des maladies diarrhéiques), obstruction des caniveaux qui aggrave les inondations en saison des pluies, et pollution des nappes phréatiques.

\begin{savaistu}[Des millions de tonnes de déchets, une faible part valorisée]
Les villes camerounaises produisent chaque année des \textbf{millions de tonnes} de déchets ménagers : Douala et Yaoundé génèrent à elles seules plusieurs milliers de tonnes \emph{par jour}. Or seule une \textbf{faible fraction} de ces déchets est réellement valorisée (recyclée, compostée ou transformée en énergie) ; l'essentiel finit en décharge. Pourtant, cette « poubelle » est une mine : le plastique peut devenir pavés et bouteilles, le papier peut devenir papier hygiénique ou combustible, et la matière organique peut devenir compost et biogaz. Le recyclage est aussi un gisement d'emplois pour les jeunes, comme le montrent les coopératives de transformation de déchets plastiques en pavés à Yaoundé et à Douala.
\end{savaistu}

\textbf{Interprétation et conclusion de la démarche.} L'observation du terrain (décharges sauvages, accumulation des plastiques) et l'analyse des données (composition des déchets, durées de dégradation) conduisent à une conclusion claire : le déchet n'est pas une fatalité mais une \emph{matière première mal placée}. La valorisation des déchets répond simultanément à trois problèmes : elle \textbf{assainit} l'environnement, elle \textbf{économise} les ressources naturelles et les énergies fossiles, et elle \textbf{crée} des revenus. Les sections suivantes de cette leçon détailleront les techniques concrètes : production des biocarburants, production d'énergie à partir du papier, recyclage du papier et des plastiques.

\begin{exoresolu}[Distinguer les notions dans une situation concrète]
\textbf{Énoncé.} Une coopérative de jeunes de Bafoussam réalise trois activités : (1) elle produit du biodiesel à partir d'huile de palme ; (2) elle transforme les sachets plastiques usagés en pavés ; (3) elle composte les déchets de marché pour les maraîchers. Pour chaque activité, indique s'il s'agit d'une production d'énergie renouvelable ou d'une valorisation de déchets, en justifiant.
\tcblower
\textbf{Solution détaillée.}
\begin{enumerate}
\item \textbf{Production de biodiesel à partir d'huile de palme} : il s'agit de la production d'une \textbf{énergie renouvelable}. Le biodiesel est un \textbf{biocarburant}, c'est-à-dire un carburant produit à partir de la biomasse végétale ; le palmier à huile se recultive à l'échelle humaine, donc la source se renouvelle.
\item \textbf{Transformation des sachets plastiques en pavés} : il s'agit d'une \textbf{valorisation matière (recyclage)} des déchets plastiques. Attention : ce n'est \emph{pas} une énergie renouvelable, car le plastique est issu du pétrole, une ressource fossile ; on transforme simplement un déchet en produit utile.
\item \textbf{Compostage des déchets de marché} : il s'agit d'une \textbf{valorisation organique} des déchets : des micro-organismes décomposent la matière organique en compost, un engrais naturel utile à l'agriculture.
\end{enumerate}
\textbf{Conclusion.} Les trois activités relèvent de la biotechnologie appliquée à l'environnement, mais seule la première produit une énergie renouvelable ; les deux autres sont des formes de valorisation des déchets.
\end{exoresolu}

\begin{attention}[Pièges de rédaction à éviter au Baccalauréat]
\begin{itemize}
\item Ne jamais écrire qu'une énergie renouvelable « ne pollue pas du tout » : la combustion de la biomasse émet du \ce{CO2} ; on dit qu'elle est \emph{proche de la neutralité carbone} si la ressource est renouvelée.
\item Ne pas confondre \textbf{biomasse} (matière organique utilisée comme source d'énergie) et \textbf{biodiversité} (variété des espèces vivantes).
\item Ne pas affirmer que « recycler supprime les déchets » : le recyclage \emph{transforme} le déchet ; il existe toujours des résidus ultimes.
\item Toujours accompagner les durées de dégradation de leur unité (mois, années) et rester sur des ordres de grandeur prudents.
\end{itemize}
\end{attention}

\begin{aretenir}[Bilan de la section]
\begin{itemize}
\item Une \textbf{énergie renouvelable} se reconstitue naturellement à l'échelle humaine (soleil, vent, eau, biomasse) ; les \textbf{énergies fossiles} (pétrole, charbon, gaz) sont épuisables et aggravent l'effet de serre.
\item Un \textbf{biocarburant} est un carburant produit à partir de la biomasse végétale ou animale (biodiesel, bioéthanol).
\item La \textbf{valorisation des déchets} transforme un déchet en produit utile (recyclage, compost) ou en énergie (biogaz, combustible).
\item Les déchets se classent en \textbf{organiques}, \textbf{papiers}, \textbf{plastiques} et \textbf{métaux} ; chaque catégorie appelle une technique de valorisation spécifique.
\item Au Cameroun, l'essentiel des déchets n'est pas valorisé : le recyclage est à la fois un enjeu sanitaire, écologique et économique.
\end{itemize}
\end{aretenir}

\coursec{Les techniques de production des biocarburants}

Dans la section précédente, nous avons défini les énergies renouvelables et montré que les déchets, loin d'être une simple nuisance, constituent une véritable ressource. Parmi les énergies renouvelables, les \cle{biocarburants} occupent une place particulière : ce sont des carburants liquides ou gazeux produits à partir de la \cle{biomasse} végétale ou de déchets organiques, grâce à des réactions chimiques ou à l'activité de micro-organismes. Contrairement au pétrole, dont les réserves s'épuisent et dont la combustion libère un carbone séquestré depuis des millions d'années, les biocarburants s'inscrivent dans le cycle court du carbone. Nous allons étudier trois techniques de production, toutes réalisables ou déjà réalisées au Cameroun : le \cle{bioéthanol}, le \cle{biodiesel} et le \cle{biogaz}.

\courssub{Le bioéthanol : production par fermentation alcoolique}

\textbf{Observation et problème.}\par
Au village, lorsqu'on laisse reposer du jus de canne à sucre ou du jus de palme dans un récipient fermé pendant quelques jours, le liquide devient pétillant et dégage une odeur d'alcool : c'est le vin de palme bien connu. D'où viennent l'alcool et les bulles de gaz ? Peut-on exploiter ce phénomène pour fabriquer un carburant ?

\begin{experience}[Mise en évidence de la fermentation alcoolique]
\textbf{Matériel :} une solution de glucose, de la levure de boulanger (\emph{Saccharomyces cerevisiae}), un ballon, un tube à dégagement, de l'eau de chaux.\par
\textbf{Protocole :} on introduit dans le ballon la solution de glucose additionnée de levure ; on bouche le ballon avec un tube à dégagement plongeant dans l'eau de chaux ; on place l'ensemble à l'étuve à \SI{30}{\celsius}.\par
\textbf{Observations :} après quelques heures, des bulles de gaz se dégagent et l'eau de chaux se trouble (présence de \ce{CO2}). Après distillation du liquide, on recueille un liquide inflammable : l'éthanol. En l'absence de levure (témoin), rien ne se produit.\par
\textbf{Interprétation :} la levure, micro-organisme vivant, transforme le glucose en éthanol et en dioxyde de carbone, en l'absence de dioxygène. C'est une réaction biologique catalysée par les enzymes de la levure.
\end{experience}

\begin{definition}[Fermentation alcoolique]
La \cle{fermentation alcoolique} est la transformation, en anaérobiose (absence de dioxygène), des sucres en éthanol et en dioxyde de carbone, réalisée par des levures telles que \emph{Saccharomyces cerevisiae}. Son équation bilan est :
\[ \ce{C6H12O6 -> 2 C2H5OH + 2 CO2} \]
Cette équation, vue en Première, est admise au programme de Terminale.
\end{definition}

\textbf{Les étapes de la production industrielle du bioéthanol.}\par
À partir de plantes riches en sucres ou en amidon — canne à sucre, maïs, manioc, tous cultivés au Cameroun — la production suit quatre étapes :

\begin{enumerate}
\item \textbf{Broyage :} la matière première (tiges de canne, grains de maïs, tubercules de manioc) est broyée pour libérer les sucres ou l'amidon et augmenter la surface de contact.
\item \textbf{Hydrolyse (saccharification) :} pour les plantes à amidon (maïs, manioc), l'amidon, polymère de glucose, est hydrolysé en glucose par des enzymes (amylases). La canne à sucre, riche en saccharose directement utilisable par les levures, n'a pas besoin de cette étape.
\item \textbf{Fermentation :} des levures sont ajoutées au moût sucré en cuve close et chauffée vers \SI{30}{\celsius}. En quelques jours, le glucose est converti en éthanol (\SI{10}{\percent} à \SI{15}{\percent} en volume) et en \ce{CO2}.
\item \textbf{Distillation :} le mélange est chauffé ; l'éthanol, plus volatil que l'eau (température d'ébullition \SI{78}{\celsius} contre \SI{100}{\celsius}), s'évapore en premier et est recueilli par condensation. On obtient un éthanol concentré (\SI{95}{\percent} et plus) utilisable comme carburant, seul ou mélangé à l'essence.
\end{enumerate}

\begin{popfigure}
\begin{center}
\begin{tikzpicture}[
box/.style={draw=popblue, fill=popblueL, rounded corners=3pt, minimum height=1.1cm, align=center, font=\small, text width=2.6cm},
fle/.style={-{Stealth[length=3mm]}, thick, popdark},
lab/.style={font=\footnotesize\itshape, popmuted, midway, above}]
\node[box, align=center] (mp) at (0,0){Matière première\\ (canne, maïs, manioc)};
\node[box] (br) at (3.6,0) {Broyage};
\node[box, align=center] (hy) at (7.2,0){Hydrolyse\\ (saccharification)};
\node[box, align=center] (fe) at (3.6,-2.2){Fermentation\\ (levures, anaérobie)};
\node[box] (di) at (7.2,-2.2) {Distillation};
\node[box, fill=popgreenL, draw=popgreen, align=center] (ca) at (10.8,-2.2){Bioéthanol\\ (carburant)};
\draw[fle] (mp) -- (br);
\draw[fle] (br) -- (hy) node[lab]{glucose};
\draw[fle] (hy.south) |- (fe.east);
\draw[fle] (fe) -- (di) node[lab]{moût alcoolisé};
\draw[fle] (di) -- (ca);
\node[font=\footnotesize, poporange] at (3.6,-3.1) {\ce{C6H12O6 -> 2 C2H5OH + 2 CO2}};
\end{tikzpicture}
\end{center}
\legende{Chaîne de production du bioéthanol : de la plante sucrière ou amylacée au carburant. La fermentation est assurée par la levure \emph{Saccharomyces cerevisiae} en anaérobiose.}
\end{popfigure}

\begin{popfigure}
\begin{center}
\begin{tikzpicture}
\begin{axis}[
width=0.85\linewidth, height=6cm,
xlabel={Temps (jours)}, ylabel={Quantité (unités arbitraires)},
xmin=0, xmax=6, ymin=0, ymax=11,
xtick={0,1,2,3,4,5,6},
legend style={at={(0.03,0.97)}, anchor=north west, font=\footnotesize},
grid=major, grid style={popline!40}]
\addplot[popcurve, domain=0:6, samples=100, popblue] {9*(1-exp(-0.8*x))};
\addlegendentry{Éthanol produit}
\addplot[popcurve, domain=0:6, samples=100, poporange] {9*(1-exp(-0.8*x))};
\addlegendentry{\ce{CO2} dégagé}
\addplot[popcurve, domain=0:6, samples=100, popgreen, dashed] {10*exp(-0.8*x)};
\addlegendentry{Glucose restant}
\end{axis}
\end{tikzpicture}
\end{center}
\legende{Évolution des quantités de glucose, d'éthanol et de \ce{CO2} au cours d'une fermentation alcoolique. L'éthanol et le \ce{CO2} apparaissent en quantités équimolaires (courbes confondues), conformément à l'équation bilan, pendant que le glucose s'épuise. Le ralentissement final s'explique par la disparition du substrat et la toxicité de l'éthanol accumulé pour les levures.}
\end{popfigure}

\begin{savaistu}[Le bioéthanol au Brésil]
Au Brésil, premier producteur mondial de bioéthanol, une grande partie des voitures roule à l'éthanol issu de la canne à sucre, pur ou mélangé à l'essence (jusqu'à \SI{27}{\percent} d'éthanol dans l'essence ordinaire). Ce programme, lancé dans les années 1970 après les chocs pétroliers, montre qu'un pays tropical peut réduire fortement sa dépendance au pétrole importé. Le Cameroun, avec sa canne à sucre (SOSUCAM à Nkoteng) et son manioc, dispose d'un potentiel comparable.
\end{savaistu}

\courssub{Le biodiesel : production par transestérification des huiles végétales}

\textbf{Observation.}\par
L'huile de palme, abondante au Cameroun (SOCAPALM, plantations du littoral), est un corps gras : elle ne brûle pas facilement dans un moteur car elle est trop visqueuse. Comment la transformer en carburant utilisable ?

\begin{definition}[Transestérification]
La \cle{transestérification} est la réaction chimique entre un corps gras (triglycéride d'une huile végétale : palme, soja, jatropha) et un alcool (méthanol ou éthanol), en présence d'un catalyseur (soude ou potasse). Elle produit des \cle{esters méthyliques} — le biodiesel — et du glycérol (glycérine), coproduit valorisable en savonnerie.
\end{definition}

\textbf{Principe de la technique.}\par
Un triglycéride est formé d'une molécule de glycérol liée à trois acides gras. L'alcool, aidé du catalyseur, « décroche » les acides gras du glycérol et forme avec eux des esters, molécules moins visqueuses, qui brûlent correctement dans un moteur diesel. Les étapes pratiques sont :

\begin{enumerate}
\item \textbf{Extraction et filtration} de l'huile végétale (palme, soja, jatropha — plante rustique poussant sur sols pauvres, très étudiée pour la production de biodiesel en zone tropicale).
\item \textbf{Réaction de transestérification :} mélange de l'huile avec l'alcool et le catalyseur, sous agitation et chauffage modéré (environ \SI{60}{\celsius}).
\item \textbf{Décantation :} le mélange se sépare en deux couches — le biodiesel (phase supérieure) et le glycérol (phase inférieure, plus dense).
\item \textbf{Purification :} lavage et séchage du biodiesel avant utilisation.
\end{enumerate}

\courssub{Le biogaz : production par méthanisation des déchets organiques}

\textbf{Observation.}\par
Dans les marécages et les rizières, des bulles de gaz remontent à la surface et s'enflamment si on les approche d'une flamme : c'est le « gaz des marais ». Ce gaz peut-il être produit volontairement à partir de nos déchets ?

\begin{definition}[Méthanisation]
La \cle{méthanisation} (ou digestion anaérobie) est la dégradation de la matière organique (déchets de cuisine, fumier, lisier, boues, résidus de récolte) par des micro-organismes, en l'absence de dioxygène. Elle est achevée par des \cle{bactéries méthanogènes} (archées) qui produisent un mélange gazeux appelé \cle{biogaz}, composé essentiellement de méthane \ce{CH4} (\SI{50}{\percent} à \SI{70}{\percent}) et de dioxyde de carbone \ce{CO2} (\SI{30}{\percent} à \SI{50}{\percent}). Le résidu solide et liquide, le \cle{digestat}, est un excellent engrais organique.
\end{definition}

\textbf{Fonctionnement d'un digesteur domestique.}\par
Les déchets organiques sont introduits dans une cuve étanche, sans air, maintenue à température tiède. Les bactéries s'y développent et dégradent la matière en plusieurs semaines. Le biogaz s'accumule dans la partie supérieure et est conduit par un tuyau vers une cuisinière ou une lampe ; le digestat est vidangé périodiquement pour fertiliser les champs.

\begin{popfigure}
\begin{center}
\begin{tikzpicture}[scale=1]
% Cuve
\draw[fill=popblueL, draw=popblue, thick] (0,0) ellipse (3.2 and 1.6);
\draw[fill=popcream, draw=popblue, thick] (0,0.35) ellipse (2.6 and 1.1);
% Entrée
\draw[fill=poporangeL, draw=poporange, thick] (-4.6,1.6) rectangle (-3.6,0.4);
\draw[-{Stealth[length=3mm]}, thick, poporange] (-3.6,0.8) -- (-2.4,0.3);
\node[font=\footnotesize, align=center, popdark] at (-4.1,2.1) {Entrée des\\ déchets organiques};
% Sortie gaz
\draw[fill=popgreenL, draw=popgreen, thick] (0.5,2.1) rectangle (1.5,1.2);
\draw[-{Stealth[length=3mm]}, thick, popgreen] (1.5,1.7) -- (3.2,1.7);
\node[font=\footnotesize, align=center, popdark] at (4.4,1.7) {Biogaz\\ (\ce{CH4} + \ce{CO2})};
\draw[-{Stealth[length=3mm]}, thick, popgreen] (0,1.1) -- (0,1.55);
% Sortie digestat
\draw[-{Stealth[length=3mm]}, thick, poppurple] (2.4,-0.9) -- (4,-1.8);
\node[font=\footnotesize, align=center, popdark] at (4.6,-2.2) {Digestat\\ (engrais)};
% Labels internes
\node[font=\small\itshape, popdark] at (0,0.35) {Cuve anaérobie};
\node[font=\footnotesize, popmuted] at (0,-0.55) {bactéries méthanogènes (sans \ce{O2})};
\end{tikzpicture}
\end{center}
\legende{Schéma d'un digesteur à biogaz domestique. Les déchets organiques entrent dans la cuve étanche (anaérobie) où les bactéries méthanogènes les dégradent ; le biogaz est recueilli en partie haute et le digestat, résidu fertilisant, est évacué en partie basse.}
\end{popfigure}

\begin{exemplebox}[Un rendement chiffré]
Une tonne de déchets organiques (fumier, résidus de cuisine, tiges de cultures) peut produire environ \SI{80}{\cubic\metre} à \SI{100}{\cubic\metre} de biogaz, soit l'équivalent énergétique d'environ \SI{50}{\kilo\gram} de butane. Une ferme disposant de quelques porcs ou bovins peut ainsi couvrir l'essentiel de ses besoins de cuisson, tout en réduisant les odeurs et les déchets.
\end{exemplebox}

\courssub{Intérêt énergétique : un bilan carbone quasi neutre}

Pourquoi dit-on que les biocarburants sont « propres » pour le climat ? Conduisons le raisonnement au lieu de l'affirmer.

\begin{enumerate}
\item La plante (canne, palme, maïs) fabrique sa matière organique par photosynthèse, en \textbf{captant} le \ce{CO2} de l'atmosphère : \ce{6 CO2 + 6 H2O -> C6H12O6 + 6 O2}.
\item La fermentation transforme ce glucose : \ce{C6H12O6 -> 2 C2H5OH + 2 CO2}. Une partie du carbone repart dans l'air, l'autre se retrouve dans l'éthanol.
\item La combustion de l'éthanol dans le moteur libère ce carbone restant : \ce{C2H5OH + 3 O2 -> 2 CO2 + 3 H2O}.
\end{enumerate}

Au total, le \ce{CO2} libéré par la combustion correspond (au premier ordre) au \ce{CO2} que la plante avait fixé quelques mois plus tôt : le carbone circule en \textbf{cycle court} entre l'atmosphère et la biomasse. À l'inverse, brûler du pétrole rejette un carbone séquestré depuis des millions d'années, qui s'\textbf{ajoute} au stock atmosphérique. Le bilan carbone des biocarburants est donc dit \emph{quasi neutre} — « quasi » car la mécanisation agricole, les engrais et la distillation consomment encore des énergies fossiles.

\begin{attention}[Complément Bac (hors programme strict, mais utile) : les limites des biocarburants]
Les biocarburants ne sont pas une solution miracle. Cultiver de la canne, du maïs ou du palmier pour produire du carburant entre en \textbf{compétition avec les cultures vivrières} : surfaces agricoles, eau et engrais détournés de l'alimentation peuvent faire augmenter le prix des denrées. Dans un pays comme le Cameroun où la sécurité alimentaire reste un enjeu, il faut privilégier les voies qui n'entrent pas en concurrence avec la nourriture : biogaz à partir des déchets, jatropha sur sols non cultivables, résidus de récolte.
\end{attention}

\begin{methode}[Expliquer une technique de production de biocarburant en 4 étapes]
Face à une question du type « Expliquez la technique de production du biogaz », structure ta réponse ainsi :
\begin{enumerate}
\item \textbf{Matière première :} identifier la biomasse utilisée (déchets organiques, huile végétale, plante sucrière).
\item \textbf{Micro-organisme ou réaction :} préciser l'agent de la transformation (bactéries méthanogènes, levures, réaction chimique de transestérification) et ses conditions (anaérobiose, catalyseur, température).
\item \textbf{Extraction / séparation :} indiquer comment le produit est recueilli (distillation, décantation, captage du gaz).
\item \textbf{Produit et utilisation :} nommer le biocarburant obtenu, sa composition et son usage (carburant, cuisson, engrais comme coproduit).
\end{enumerate}
\end{methode}

\begin{attention}[Erreur fréquente : biogaz et gaz naturel]
Ne confonds jamais le \cle{biogaz} et le \cle{gaz naturel}. Le biogaz est d'origine \textbf{biologique} : il est produit aujourd'hui par la fermentation de déchets organiques, donc renouvelable. Le gaz naturel est d'origine \textbf{fossile} : il s'est formé il y a des millions d'années à partir d'organismes morts enfouis, comme le pétrole, et il n'est pas renouvelable. Tous deux contiennent du méthane, mais leur origine — et donc leur bilan carbone — est radicalement différente.
\end{attention}

\begin{exoresolu}[Bilan d'une fermentation alcoolique]
\textbf{Énoncé.} Une unité artisanale fait fermenter \SI{360}{\gram} de glucose par des levures. (1) Écris l'équation bilan de la fermentation alcoolique. (2) Calcule la masse d'éthanol théoriquement obtenue. On donne : $M(\ce{C6H12O6}) = \SI{180}{\gram\per\mole}$ ; $M(\ce{C2H5OH}) = \SI{46}{\gram\per\mole}$.
\tcblower
\textbf{Solution.} (1) L'équation bilan est : \ce{C6H12O6 -> 2 C2H5OH + 2 CO2}.\par
(2) Quantité de glucose : $n = \dfrac{m}{M} = \dfrac{360}{180} = \SI{2}{\mole}$. D'après l'équation, 1 mole de glucose donne 2 moles d'éthanol, donc $n(\text{éthanol}) = 2 \times 2 = \SI{4}{\mole}$. Masse d'éthanol : $m = n \times M = 4 \times 46 = \SI{184}{\gram}$. La fermentation de \SI{360}{\gram} de glucose produit théoriquement \SI{184}{\gram} d'éthanol (et \SI{176}{\gram} de \ce{CO2}).
\end{exoresolu}

\begin{exercice}[Pour t'entraîner]
Une coopérative agricole de la région de l'Ouest collecte \SI{5}{\tonne} de fumier et de résidus végétaux par mois. (1) Estime le volume mensuel de biogaz produit. (2) Nomme le micro-organisme responsable et précise la condition indispensable du milieu. (3) Cite les deux coproduits de la méthanisation et leur usage.
\end{exercice}

\begin{corrige}[Exercice : production de biogaz d'une coopérative]
(1) À raison de \SI{80}{\cubic\metre} à \SI{100}{\cubic\metre} de biogaz par tonne, \SI{5}{\tonne} de déchets produisent environ \SI{400}{\cubic\metre} à \SI{500}{\cubic\metre} de biogaz par mois.\par
(2) Ce sont les \textbf{bactéries méthanogènes}, qui travaillent impérativement en \textbf{anaérobiose} (absence totale de dioxygène, d'où la cuve étanche).\par
(3) Le \textbf{biogaz} (\ce{CH4} + \ce{CO2}), utilisé pour la cuisson et l'éclairage, et le \textbf{digestat}, résidu utilisé comme engrais organique pour les champs.
\end{corrige}

\begin{aretenir}[L'essentiel de la section]
\begin{itemize}
\item Le \cle{bioéthanol} s'obtient par \textbf{fermentation alcoolique} des sucres (canne, maïs, manioc) par la levure \emph{Saccharomyces cerevisiae}, selon \ce{C6H12O6 -> 2 C2H5OH + 2 CO2} ; étapes : broyage, hydrolyse, fermentation, distillation.
\item Le \cle{biodiesel} s'obtient par \textbf{transestérification} des huiles végétales (palme, soja, jatropha) avec un alcool et un catalyseur ; coproduit : le glycérol.
\item Le \cle{biogaz} (\ce{CH4} + \ce{CO2}) s'obtient par \textbf{méthanisation} des déchets organiques par des bactéries méthanogènes en anaérobiose ; coproduit : le digestat, engrais.
\item Le bilan carbone des biocarburants est \textbf{quasi neutre} : le \ce{CO2} libéré a été préalablement fixé par la plante (cycle court du carbone).
\item Limite majeure : la compétition possible avec les cultures vivrières.
\end{itemize}
\end{aretenir}

\coursec{La production d'énergie à partir du papier}

Après avoir vu comment les végétaux (canne à sucre, maïs, huile de palme) peuvent être transformés en \cle{biocarburants} liquides, une question naturelle se pose : que faire des déchets papiers qui s'accumulent chaque jour dans nos écoles, nos administrations et nos marchés ? À Yaoundé comme à Douala, des tonnes de vieux cahiers, de cartons et de journaux finissent brûlés à l'air libre ou abandonnés dans les ruisseaux. Pourtant, ce papier est une véritable \textbf{source d'énergie cachée}. Cette section montre comment, par une démarche rigoureuse de collecte, de transformation et de combustion contrôlée, le papier peut être \cle{valorisé} en chaleur et en électricité.

\courssub{Le papier : une biomasse transformée riche en énergie chimique}

\textbf{Observation et questionnement.} Prends une feuille de papier bien sèche et approche-la d'une flamme : elle s'enflamme rapidement et dégage une chaleur sensible. D'où vient cette énergie ? Pourquoi le papier brûle-t-il si facilement alors qu'il n'est ni du bois ni du charbon ?

\begin{experience}[Mise en évidence de l'énergie stockée dans le papier]
\textbf{Matériel :} papier sec, papier humide, pinces métalliques, bougie, thermomètre, bécher d'eau, balance.

\textbf{Protocole :}
\begin{enumerate}
\item Peser \SI{1}{g} de papier sec et \SI{1}{g} de papier humidifié puis essoré.
\item Brûler chaque échantillon sous un bécher contenant \SI{50}{mL} d'eau à température initiale $\theta_i = \SI{25}{\celsius}$.
\item Mesurer la température finale $\theta_f$ de l'eau après combustion complète.
\end{enumerate}

\textbf{Observations :}
\begin{itemize}
\item Le papier sec brûle entièrement, avec une flamme vive ; la température de l'eau s'élève nettement ($\theta_f \approx \SI{40}{\celsius}$).
\item Le papier humide brûle mal, dégage beaucoup de fumée ; l'élévation de température est faible ($\theta_f \approx \SI{28}{\celsius}$).
\end{itemize}

\textbf{Interprétation :} la combustion du papier libère de la chaleur qui est transférée à l'eau. Cette chaleur provient de l'\cle{énergie chimique} stockée dans les molécules constituant le papier. L'eau contenue dans le papier humide absorbe une partie de cette énergie pour s'évaporer, ce qui réduit le rendement de la combustion.
\end{experience}

\begin{definition}[Nature chimique du papier]
Le \cle{papier} est un matériau constitué essentiellement de \cle{cellulose} (environ 70 à 90 \% de sa masse), un polysaccharide de formule brute $(\ce{C6H10O5})_n$ formé de longues chaînes de glucose, extrait des fibres végétales (bois, coton, bagasse de canne). La cellulose est une \textbf{matière organique combustible} : ses liaisons chimiques carbone-carbone et carbone-hydrogène renferment de l'énergie chimique libérable par combustion.
\end{definition}

\textbf{D'où vient cette énergie ?} Le papier est une \cle{biomasse transformée}. La cellulose a été fabriquée par les végétaux grâce à la \textbf{photosynthèse}, qui convertit l'énergie lumineuse du soleil en énergie chimique stockée dans le glucose puis dans la cellulose (savoir admis) :
\[ \ce{6CO2 + 6H2O ->[lumiere, chlorophylle] C6H12O6 + 6O2} \]
Brûler du papier, c'est donc en réalité \textbf{libérer une énergie solaire ancienne}, captée par l'arbre des années auparavant. C'est pourquoi l'énergie tirée du papier est classée parmi les \cle{énergies renouvelables} (à condition que les forêts soient replantées).

\begin{definition}[Pouvoir calorifique]
Le \cle{pouvoir calorifique} (PC) d'un combustible est la quantité de chaleur libérée par la combustion complète d'une unité de masse de ce combustible. Il s'exprime en joules par kilogramme (\SI{}{J/kg}) ou, plus couramment, en mégajoules par kilogramme (\SI{}{MJ/kg}). Plus le pouvoir calorifique est élevé, plus le combustible est « énergétique ».
\end{definition}

\begin{center}
\begin{tabularx}{\linewidth}{|l|X|X|}
\hline
\rowcolor{popblueL} \textbf{Combustible} & \textbf{Pouvoir calorifique (MJ/kg)} & \textbf{Remarque} \\
\hline
Papier sec & 13 à 17 & comparable au bois de chauffage \\
\hline
Bois sec & 15 à 18 & combustible traditionnel au Cameroun \\
\hline
Charbon de bois & 28 à 32 & obtenu par carbonisation du bois \\
\hline
\end{tabularx}
\end{center}

\begin{popfigure}
\begin{center}
\begin{tikzpicture}
\begin{axis}[
  ybar, bar width=0.9cm,
  width=0.85\linewidth, height=6cm,
  ylabel={Pouvoir calorifique (MJ/kg)},
  symbolic x coords={Papier sec, Bois sec, Charbon de bois},
  xtick=data,
  ymin=0, ymax=35,
  ymajorgrids=true, grid style={popline!40},
  nodes near coords,
  every node near coord/.append style={font=\small\bfseries, color=popdark},
]
\addplot[fill=popblue, draw=popdark] coordinates {(Papier sec,15) (Bois sec,16.5) (Charbon de bois,30)};
\end{axis}
\end{tikzpicture}
\end{center}
\legende{Comparaison des pouvoirs calorifiques moyens du papier sec, du bois sec et du charbon de bois. Le papier sec possède un pouvoir calorifique voisin de celui du bois : c'est un combustible réellement valorisable.}
\end{popfigure}

\begin{aretenir}[L'essentiel sur le papier combustible]
Le papier, fait de cellulose issue de la photosynthèse, est une biomasse transformée dont le pouvoir calorifique (environ \SI{15}{MJ/kg}) est comparable à celui du bois. Sa combustion libère l'énergie chimique stockée dans la cellulose sous forme de chaleur utilisable.
\end{aretenir}

\courssub{La technique de production d'énergie à partir du papier}

\textbf{Situation-problème.} Un lycée de Bafoussam produit chaque année plusieurs centaines de kilogrammes de vieux cahiers, de copies corrigées et de cartons. Comment transformer ce « problème » en ressource énergétique ? La réponse tient en une chaîne technique en quatre étapes.

\begin{methode}[Expliquer la production d'énergie à partir du papier en 4 étapes]
Pour rédiger une explication complète au Baccalauréat, décris toujours la chaîne dans l'ordre logique suivant :
\begin{enumerate}
\item \textbf{Collecte et tri :} rassembler les déchets papiers et éliminer les corps étrangers (agrafes, plastiques, papiers souillés ou gras) qui gêneraient la combustion ou pollueraient les fumées.
\item \textbf{Séchage :} exposer le papier au soleil ou le faire sécher à l'air jusqu'à obtention d'un papier bien sec, car l'humidité absorbe la chaleur de combustion et réduit le rendement.
\item \textbf{Compression :} déchiqueter le papier, éventuellement le tremper puis le compresser fortement (presse manuelle ou mécanique) pour obtenir des \cle{briquettes} ou des \cle{granulés} denses, faciles à stocker et à transporter, qui brûlent lentement et régulièrement.
\item \textbf{Combustion contrôlée avec récupération d'énergie :} brûler les briquettes dans un foyer ou une chaudière ; la chaleur dégagée chauffe directement (cuisson, chauffage) ou transforme de l'eau en \textbf{vapeur sous pression} qui entraîne une \textbf{turbine} couplée à un alternateur : on produit alors de l'\textbf{électricité} (incinération avec récupération d'énergie).
\end{enumerate}
\end{methode}

\begin{popfigure}
\begin{center}
\begin{tikzpicture}[
  node distance=0.55cm,
  box/.style={rectangle, rounded corners=3pt, draw=popdark, thick, fill=popblueL, minimum height=1.1cm, text width=2.6cm, align=center, font=\small},
  box2/.style={rectangle, rounded corners=3pt, draw=popdark, thick, fill=poporangeL, minimum height=1.1cm, text width=2.6cm, align=center, font=\small},
  box3/.style={rectangle, rounded corners=3pt, draw=popdark, thick, fill=popgreenL, minimum height=1.1cm, text width=2.6cm, align=center, font=\small},
  arr/.style={-{Stealth[length=3mm]}, very thick, popdark}
]
\node[box, align=center] (a){\textbf{1. Collecte et tri}\\ déchets papiers propres};
\node[box, right=of a, align=center] (b){\textbf{2. Séchage}\\ papier sec, non souillé};
\node[box2, right=of b, align=center] (c){\textbf{3. Compression}\\ briquettes / granulés};
\node[box2, right=of c, align=center] (d){\textbf{4. Combustion}\\ foyer ou chaudière};
\node[box3, below=0.8cm of d, align=center] (e){\textbf{Chaleur}\\ eau $\rightarrow$ vapeur sous pression};
\node[box3, left=of e, align=center] (f){\textbf{Turbine + alternateur}\\ production d'électricité};
\draw[arr] (a) -- (b);
\draw[arr] (b) -- (c);
\draw[arr] (c) -- (d);
\draw[arr] (d) -- (e);
\draw[arr] (e) -- (f);
\end{tikzpicture}
\end{center}
\legende{Chaîne de valorisation énergétique du papier : des déchets papiers à l'électricité. La chaleur de combustion peut aussi être utilisée directement (cuisson, chauffage, séchage agricole).}
\end{popfigure}

\textbf{Le principe physique de la centrale.} Dans la chaudière, la combustion des briquettes porte l'eau à haute température ; la vapeur produite, sous pression, pousse les pales d'une turbine qui tourne. La turbine entraîne un alternateur qui convertit l'énergie mécanique en énergie électrique. C'est exactement le principe des centrales thermiques classiques, sauf que le combustible est ici un \textbf{déchet valorisé} : on parle d'\cle{incinération avec récupération d'énergie} (ou « valorisation énergétique »).

\begin{attention}[Condition d'efficacité : papier sec et non souillé]
La combustion n'est efficace que si le papier est \textbf{sec} (l'eau résiduelle absorbe la chaleur pour s'évaporer et fait chuter le rendement) et \textbf{non souillé} (papiers gras, plastifiés ou mélangés à des matières organiques brûlent mal et dégagent des fumées toxiques). Le tri et le séchage ne sont donc pas des étapes facultatives : ce sont les conditions du succès de toute la chaîne.
\end{attention}

\begin{attention}[Erreur fréquente : brûler n'est pas valoriser !]
Brûler le papier \textbf{à l'air libre} (dans la cour, dans un bidon, au bord de la route) n'est \textbf{pas} de la valorisation : l'énergie est perdue dans l'atmosphère, la combustion incomplète dégage des fumées nocives et des gaz polluants, et l'incendie peut se propager. Il n'y a valorisation énergétique que si la combustion est \textbf{contrôlée} (foyer, chaudière fermée) et si la chaleur est \textbf{récupérée} et utilisée (chauffage, cuisson, vapeur, électricité). Au Baccalauréat, cette distinction est souvent attendue.
\end{attention}

\begin{exoresolu}[Bilan énergétique d'une briquette de papier]
\textbf{Énoncé.} Un centre de santé rural utilise des briquettes de papier compressé pour chauffer l'eau de stérilisation. Le pouvoir calorifique du papier sec est $PC = \SI{15}{MJ/kg}$. On brûle complètement $m = \SI{2}{kg}$ de briquettes.
\begin{enumerate}
\item Calculer la quantité de chaleur $Q$ libérée par cette combustion.
\item Sachant qu'il faut environ $\SI{0,42}{MJ}$ pour porter 1 litre d'eau de $\SI{25}{\celsius}$ à $\SI{100}{\celsius}$, combien de litres d'eau peut-on théoriquement porter à ébullition ?
\item En réalité, on ne porte que 50 litres à ébullition. Calculer le rendement énergétique du foyer.
\end{enumerate}
\tcblower
\textbf{Solution détaillée.}
\begin{enumerate}
\item La chaleur libérée est le produit de la masse par le pouvoir calorifique :
\[ Q = m \times PC = 2 \times 15 = \SI{30}{MJ}. \]
\item Le nombre de litres théorique est :
\[ N = \frac{Q}{\num{0,42}} = \frac{30}{\num{0,42}} \approx 71 \text{ litres}. \]
\item Le rendement est le rapport de l'énergie utile à l'énergie fournie. L'énergie réellement utilisée est $Q_u = 50 \times \num{0,42} = \SI{21}{MJ}$, d'où :
\[ \eta = \frac{Q_u}{Q} = \frac{21}{30} = \num{0,70} \text{ soit } 70\,\%. \]
Les 30 \% restants correspondent aux pertes (chaleur perdue par les parois, fumées chaudes, combustion imparfaite).
\end{enumerate}
\textbf{Conclusion :} même avec un foyer simple, la valorisation énergétique du papier permet de récupérer une fraction importante de l'énergie stockée dans les déchets.
\end{exoresolu}

\begin{savaistu}[Des briquettes de papier contre la déforestation]
Au Cameroun, la demande en bois de chauffe et en charbon de bois est l'une des premières causes de la déforestation des savanes et des lisières forestières. Les \textbf{briquettes de papier compressé}, fabriquées artisanalement à partir de vieux papiers trempés, malaxés et pressés puis séchés au soleil, peuvent remplacer partiellement le charbon de bois pour la cuisson domestique. Des associations et des jeunes entrepreneurs de Yaoundé et de Douala transforment ainsi les déchets de papeteries et d'imprimeries en combustible bon marché : on réduit à la fois la pollution urbaine, la pression sur les forêts et les dépenses énergétiques des ménages. Recycler le papier, c'est donc aussi protéger le patrimoine forestier national.
\end{savaistu}

\begin{exercice}[À toi de jouer]
Un quartier collecte chaque semaine \SI{500}{kg} de déchets papiers. Après tri et séchage, on obtient \SI{400}{kg} de briquettes de pouvoir calorifique $PC = \SI{14}{MJ/kg}$.
\begin{enumerate}
\item Calculer l'énergie totale disponible chaque semaine.
\item Comparer cette énergie à celle que fournirait la même masse de charbon de bois ($PC = \SI{30}{MJ/kg}$).
\item Citer deux conditions nécessaires pour que cette valorisation soit efficace et non polluante.
\end{enumerate}
\end{exercice}

\begin{corrige}[Exercice : briquettes du quartier]
\begin{enumerate}
\item $Q = m \times PC = 400 \times 14 = \SI{5600}{MJ}$, soit \SI{5,6}{GJ} par semaine.
\item Avec du charbon : $Q' = 400 \times 30 = \SI{12000}{MJ}$. Le papier fournit donc environ 47 \% de l'énergie du charbon à masse égale, mais il est gratuit (déchet récupéré) et épargne les forêts.
\item Il faut un papier \textbf{sec et non souillé} (tri et séchage rigoureux) et une \textbf{combustion contrôlée avec récupération de chaleur} (foyer ou chaudière), jamais un brûlage à l'air libre.
\end{enumerate}
\end{corrige}

\begin{aretenir}[Bilan de la section]
\begin{itemize}
\item Le papier est une \textbf{biomasse transformée} riche en cellulose ; son énergie chimique provient de la photosynthèse.
\item Son pouvoir calorifique ($\approx \SI{15}{MJ/kg}$) est proche de celui du bois.
\item La chaîne de valorisation comporte 4 étapes : \textbf{collecte-tri, séchage, compression en briquettes, combustion contrôlée} avec récupération de la chaleur (chauffage direct ou production de vapeur, de turbine et d'électricité).
\item La valorisation exige un papier \textbf{sec et propre} et une \textbf{récupération contrôlée de l'énergie} : brûler à l'air libre est une pollution, pas une valorisation.
\end{itemize}
\end{aretenir}

Cette énergie récupérée n'est toutefois qu'une première vie nouvelle pour le papier. Comme nous le verrons dans la section suivante, les fibres de cellulose elles-mêmes peuvent être régénérées pour fabriquer de nouveaux produits : serviettes jetables, papier hygiénique et papier-cadeau — c'est le \textbf{recyclage matière} du papier.

\coursec{Le recyclage du papier : serviettes jetables, papier hygiénique, papier-cadeau}

Dans la section précédente, nous avons vu qu'une partie des déchets papiers peut être valorisée énergétiquement, par combustion contrôlée ou par méthanisation. Mais brûler le papier, c'est détruire définitivement une matière première précieuse : la \cle{fibre de cellulose}. Or cette fibre, issue du bois, peut servir plusieurs fois avant d'être trop dégradée. Une deuxième voie de valorisation, souvent plus rentable et plus respectueuse des forêts, consiste donc à \textbf{récupérer les fibres des papiers usagés pour fabriquer de nouveaux produits papetiers} : c'est le \cle{recyclage du papier}. Dans les quartiers de Yaoundé, Douala ou Bafoussam, des associations de jeunes et des coopératives transforment déjà les vieux cahiers, journaux et cartons en serviettes jetables, en papier hygiénique et en papier-cadeau vendus sur les marchés. Comprendre cette technique, c'est à la fois maîtriser un savoir exigible au Baccalauréat et découvrir une piste concrète d'entrepreneuriat vert.

\courssub{Principe général du recyclage du papier}

\begin{definition}[Recyclage du papier]
Le \cle{recyclage du papier} est l'ensemble des opérations qui consistent à récupérer les \textbf{fibres de cellulose} contenues dans les papiers et cartons usagés, à les remettre en suspension dans l'eau (mise en pâte), puis à les reformer en nouvelles feuilles pour fabriquer de nouveaux produits : serviettes jetables, papier hygiénique, papier-cadeau, carton d'emballage, etc.
\end{definition}

L'intuition est simple : une feuille de papier n'est rien d'autre qu'un \textbf{feutre de fibres de cellulose} enchevêtrées et pressées, liées entre elles par des \textbf{liaisons hydrogène}. Tant que ces fibres restent assez longues et intactes, il suffit de les séparer à nouveau dans l'eau (rupture des liaisons hydrogène par l'eau, action mécanique), de les nettoyer, puis de les re-presser et sécher (reformation des liaisons hydrogène) pour obtenir une feuille neuve. Le recyclage est donc une \emph{re-fabrication} du papier à partir de sa propre matière première.

\begin{attention}[Tous les papiers ne se recyclent pas]
Erreur fréquente : croire que tout déchet « papier » peut entrer dans le circuit de recyclage. En réalité, certains papiers sont \textbf{impropres au recyclage} et doivent être éliminés dès le tri, sous peine de gâcher la qualité de tout le lot recyclé :
\begin{itemize}
\item les \textbf{papiers plastifiés ou couchés} (certains emballages, couvertures glacées, sachets alimentaires), dont le film plastique ou la couche de kaolin ne se séparent pas des fibres ;
\item les \textbf{papiers gras ou souillés} (papier ayant enveloppé de la viande, de l'huile, excréments, serviettes hygiéniques usagées), dont les graisses et saletés contaminent toute la pâte et empêchent le collage des fibres ;
\item les \textbf{papiers thermiques} (tickets de caisse, fax anciens) et les papiers carbone, chargés en produits chimiques (développeurs de couleur) ;
\item les papiers souillés par des produits toxiques (peintures, solvants, pesticides).
\end{itemize}
\end{attention}

\courssub{Les cinq étapes de la technique de recyclage du papier}

\textbf{Étape 1 : Collecte et tri des déchets papiers}\par
La première étape consiste à \textbf{rassembler} les papiers usagés (vieux cahiers, journaux, papiers de bureau, cartons d'emballage) puis à les \textbf{trier} soigneusement. Le tri a un double objectif :
\begin{itemize}
\item \textbf{éliminer les corps étrangers} : agrafes, trombones, spirales de cahiers, rubans adhésifs, films plastiques, ficelles, qui endommageraient le matériel (broyeurs, pompes) et tacheraient la pâte ;
\item \textbf{écarter les papiers non recyclables} (plastifiés, gras, souillés, thermiques) et séparer les papiers par catégories de qualité (papiers blancs de bureau d'un côté, journaux et cartons de l'autre), car la qualité de la matière première conditionne la qualité du produit final.
\end{itemize}

\textbf{Étape 2 : La mise en pâte (désagrégation)}\par
\begin{definition}[Mise en pâte]
La \cle{mise en pâte} est l'opération qui consiste à \textbf{tremper le papier trié et déchiré en morceaux dans l'eau}, puis à le \textbf{brasser énergiquement} (agitation mécanique dans un pulpeur industriel, mixage ou pilage artisanal) jusqu'à la \textbf{désagrégation complète} des feuilles. On obtient une suspension de fibres de cellulose individualisées dans l'eau, appelée \cle{pulpe} (ou pâte à papier secondaire).
\end{definition}

L'eau pénètre entre les fibres, rompt les liaisons hydrogène et ramollit la structure. Le brassage mécanique achève la séparation des fibres. Plus le brassage est long et énergique, plus les fibres sont désagrégées et fines, et plus la pâte est homogène : ce paramètre sera déterminant pour le choix du produit final (papier hygiénique très fin ou papier-cadeau plus épais).

\textbf{Étape 3 : Désencrage et blanchiment éventuel}\par
\begin{definition}[Désencrage]
Le \cle{désencrage} est l'opération qui consiste à \textbf{débarrasser la pulpe des encres d'impression}, des colles et des salissures. À l'échelle industrielle, on utilise principalement la \textbf{flotation} : on injecte de l'air dans la pulpe en présence de \textbf{tensioactifs} (savons, agents mouillants) ; les bulles d'air se chargent des particules d'encre hydrophobes et les remontent en surface sous forme d'une mousse que l'on écrème. En complément, des \textbf{lavages} sur tamis fins éliminent les plus fines particules. Un \textbf{blanchiment} chimique (peroxyde d'hydrogène, hydrosulfite de sodium) peut suivre lorsque l'on vise un produit blanc pur (papier hygiénique, serviettes de table).
\end{definition}

Cette étape est \textbf{modulée selon le produit visé} : elle est poussée (flotation + lavage + blanchiment) pour le papier hygiénique et les serviettes blanches, mais elle peut être réduite, voire supprimée, pour le papier-cadeau que l'on colorera de toute façon, ou pour le carton d'emballage. C'est une économie d'eau, d'énergie et de produits chimiques.

\textbf{Étape 4 : Égouttage, pressage et séchage}\par
La pâte lavée (consistance \SI{1}{\%} à \SI{3}{\%} de fibres) est \textbf{étalée en couche fine et régulière sur un tamis} (toile fine sans fin sur machine continue industrielle ; cadre à moustiquaire en artisanat). L'eau s'écoule par gravité à travers les mailles (\textbf{égouttage}) tandis que les fibres, retenues, s'enchevêtrent à nouveau en une nappe continue. On exerce ensuite une \textbf{pression} mécanique (rouleaux presseurs industriels ; planches lestées ou presse hydraulique en artisanat) pour chasser l'eau résiduelle (consistance \SI{40}{\%} à \SI{50}{\%}) et souder les fibres entre elles par reformations des liaisons hydrogène. La feuille humide est enfin mise à \textbf{sécher} (cylindres chauffés à la vapeur industriels ; air libre, soleil ou étuve en artisanat). Au séchage, la feuille devient solide et rigide.

\textbf{Étape 5 : Découpe et façonnage}\par
Les feuilles sèches (ou la bobine mère industrielle) sont enfin \textbf{découpées et façonnées} selon le produit final : pliage et découpe en carrés pour les serviettes, enroulage en rouleaux perforés pour le papier hygiénique, coloration en masse ou en surface et décoration pour le papier-cadeau.

\begin{methode}[Décrire la technique de recyclage du papier en 5 étapes ordonnées]
Pour rédiger correctement cette technique au Baccalauréat, procède toujours dans l'ordre chronologique et nomme chaque étape :
\begin{enumerate}
\item \textbf{Collecte et tri} : rassembler les papiers usagés, éliminer agrafes, plastiques et papiers souillés ou plastifiés ;
\item \textbf{Mise en pâte} : déchirer, tremper dans l'eau et brasser jusqu'à obtention d'une pulpe de fibres désagrégées ;
\item \textbf{Désencrage} (flotation et/ou lavage) \textbf{et blanchiment éventuel} : traiter la pulpe pour éliminer encres et colorer selon le produit visé ;
\item \textbf{Égouttage, pressage et séchage} : former la feuille sur tamis, presser pour chasser l'eau, puis sécher ;
\item \textbf{Découpe et façonnage} : couper, plier, colorer ou enrouler selon le produit final (serviette, papier hygiénique, papier-cadeau).
\end{enumerate}
Conclus toujours par l'intérêt de l'opération : économie de fibres vierges (arbres, eau, énergie) et réduction des déchets mis en décharge.
\end{methode}

\begin{popfigure}
\begin{center}
\begin{tikzpicture}[
  font=\small,
  etape/.style={rectangle, rounded corners=4pt, draw=popblue, fill=popblueL, text width=3.2cm, minimum height=1.6cm, align=center},
  fleche/.style={-{Stealth[length=3mm]}, very thick, poporange},
  prod/.style={rectangle, rounded corners=4pt, draw=popgreen, fill=popgreenL, text width=3.2cm, minimum height=1.6cm, align=center}
]
\node[etape, align=center] (e1){\textbf{1. Collecte \& Tri}\\ élimination corps étrangers,\\ papiers non recyclables};
\node[etape, right=0.8cm of e1, align=center] (e2){\textbf{2. Mise en pâte}\\ trempage + brassage intense\\ $\rightarrow$ pulpe fibreuse};
\node[etape, right=0.8cm of e2, align=center] (e3){\textbf{3. Désencrage / Blanchiment}\\ flotation, lavage, blanchiment\\ (selon produit visé)};
\node[etape, right=0.8cm of e3, align=center] (e4){\textbf{4. Formage / Pressage / Séchage}\\ égouttage sur tamis, pressage,\\ séchage $\rightarrow$ feuilles sèches};
\node[etape, right=0.8cm of e4, align=center] (e5){\textbf{5. Façonnage}\\ découpe, pliage, gaufrage,\\ enroulement, coloration};
\draw[fleche] (e1) -- (e2);
\draw[fleche] (e2) -- (e3);
\draw[fleche] (e3) -- (e4);
\draw[fleche] (e4) -- (e5);
\node[prod, below=1.2cm of e3, align=center] (p){\textbf{Produits finis}\\ Serviettes jetables\\ Papier hygiénique\\ Papier-cadeau\\ Carton d'emballage};
\draw[fleche, popgreen] (e5) -- ++(0,-0.6) -| (p);
\end{tikzpicture}
\end{center}
\legende{Organigramme des cinq étapes du recyclage du papier (Procédé industriel / semi-industriel). La qualité du tri (1) conditionne la pureté de la pulpe (2). Le désencrage (3) est adapté au degré de blancheur requis. Le formage sur tamis (4) reforme le feutre de fibres. Le façonnage (5) adapte la feuille à l'usage final.}
\end{popfigure}

\courssub{Fabrication artisanale en milieu scolaire : le cadre à tamis}

La bonne nouvelle est que cette technique ne demande \textbf{aucune machine industrielle complexe} : elle est réalisable dans un club scientifique de lycée ou un atelier de quartier avec du matériel très simple. C'est une activité de TP idéale pour illustrer le cycle de la matière.

\begin{experience}[Fabriquer une feuille de papier recyclé au lycée (TP)]
\textbf{Matériel} : vieux cahiers ou journaux triés (sans agrafes), une bassine d'eau (\SI{\approx 10}{L}), un mixeur plongeant ou un pilon en bois, un \textbf{cadre à tamis} (cadre de bois \SI{30x40}{cm} sur lequel est tendue et agrafée une moustiquaire fine ou une toile de nylon), une éponge plate, deux planches lisses (\SI{40x50}{cm}), des poids (livres, parpaings), éventuellement des colorants alimentaires, des pétales de fleurs séchées ou des fils de coton pour décorer.

\textbf{Protocole} :
\begin{enumerate}
\item \textbf{Trempage} : déchirer le papier en carrés de \SI{\approx 2}{cm} et le faire tremper \SI{\ge 2}{h} (idéalement une nuit) dans l'eau pour bien ramollir les fibres.
\item \textbf{Mixage} : mixer par petites quantités (\SI{100}{g} papier pour \SI{1}{L} eau) jusqu'à obtention d'une pulpe homogène, sans morceaux visibles. Verser la pulpe dans la bassine, compléter en eau pour obtenir une consistance « soupe légère » (\SI{\approx 1}{\%} fibres).
\item \textbf{Formage} : bien mélanger la bassine. Plonger le cadre à tamis \textbf{horizontalement} au fond de la bassine, le relever \textbf{lentement et à l'horizontale} : une couche régulière de fibres se dépose sur la toile. Laisser égoutter \SI{30}{s}.
\item \textbf{Pressage (couchage)} : poser le cadre face fibres sur une planche propre (ou un feutre). Éponger énergiquement le dos de la toile pour aspirer l'eau. Retourner l'ensemble, retirer délicatement le cadre : la feuille humide reste sur la planche/feutre. Recouvrir d'une seconde planche/feutre et presser fortement (monter dessus, ou utiliser une presse/poids lourds) \SI{10}{min}.
\item \textbf{Séchage} : décoller la feuille (elle tient toute seule) et la faire sécher à plat à l'ombre (évite le gondolage) ou au soleil (plus rapide) pendant \SI{24}{h} à \SI{48}{h}.
\end{enumerate}

\textbf{Observations} : après séchage, on obtient une feuille solide, légèrement rugueuse, dont l'épaisseur dépend de la quantité de pulpe déposée (grammage). L'ajout de colorants dans la pulpe (coloration en masse) ou de pétales en surface lors du formage produit un papier-cadeau décoratif.

\textbf{Interprétation} : le trempage et le mixage rompent les liaisons hydrogène inter-fibres (désagrégation) ; le passage sur tamis répartit les fibres en nappe aléatoire (formage) ; le pressage rapproche les fibres et chasse l'eau ; le séchage permet la reformation définitive des liaisons hydrogène qui soudent les fibres en une feuille cohérente. On a reproduit, à l'échelle du laboratoire, le procédé industriel de la machine à papier (machine de Fourdrinier).
\end{experience}

\begin{popfigure}
\begin{center}
\begin{tikzpicture}[font=\small, scale=0.9]
% Cadre extérieur
\draw[very thick, popdark, fill=popgoldL] (0,0) rectangle (8,5);
% Toile (tamis)
\draw[fill=popblueL, draw=popline, line width=0.5pt] (0.5,0.5) rectangle (7.5,4.5);
% Mailles du tamis (représentation simplifiée)
\foreach \x in {1,1.5,...,7} {\draw[popline, very thin] (\x,0.5) -- (\x,4.5);}
\foreach \y in {1,1.5,...,4} {\draw[popline, very thin] (0.5,\y) -- (7.5,\y);}
% Fibres déposées (aspect feutre)
\foreach \i in {1,...,60} {
  \pgfmathsetmacro{\xi}{0.7+rand*6.6}
  \pgfmathsetmacro{\yi}{0.7+rand*3.6}
  \pgfmathsetmacro{\ang}{rand*180}
  \pgfmathsetmacro{\len}{0.15+rand*0.15}
  \draw[poporange, thick, line cap=round] (\xi,\yi) -- ++(\ang:\len);
}
% Flèches d'égouttage
\draw[-{Stealth[length=2mm]}, popblue, thick] (4, -0.3) -- (4, 0.2) node[midway, right] {Eau};
\draw[-{Stealth[length=2mm]}, popblue, thick] (6, -0.3) -- (6, 0.2);
% Légendes
\draw[-{Stealth[length=2.5mm]}, popdark] (8.3, 4.7) node[right, align=left, text width=3.5cm] {\textbf{1} Cadre de bois rigide\\ (ex: tasseaux \SI{3x3}{cm})} -- (7.8, 4.7);
\draw[-{Stealth[length=2.5mm]}, popblue] (8.3, 2.5) node[right, align=left, text width=3.5cm] {\textbf{2} Toile fine tendue (tamis)\\ moustiquaire nylon ou inox} -- (7.6, 2.5);
\draw[-{Stealth[length=2.5mm]}, poporange] (8.3, 1.0) node[right, align=left, text width=3.5cm] {\textbf{3} Nappe de fibres de cellulose\\ déposée par égouttage} -- (7.0, 1.2);
\draw[-{Stealth[length=2.5mm]}, popgreen] (-0.3, 2.5) node[left, align=right, text width=3.5cm] {\textbf{4} L'eau traverse les mailles\\ (égouttage gravitaire)} -- (0.5, 2.5);
\draw[-{Stealth[length=2.5mm]}, poppurple] (-0.3, 0.5) node[left, align=right, text width=3.5cm] {\textbf{5} Feuille humide couchée\\ sur feutre/planche pour pressage} -- (0.5, 0.5);
\end{tikzpicture}
\end{center}
\legende{Schéma de principe du cadre à tamis (moule à papier) pour fabrication artisanale. \textbf{1} Cadre bois ; \textbf{2} Toile (tamis) ; \textbf{3} Fibres retenues formant la feuille ; \textbf{4} Égouttage de l'eau ; \textbf{5} Couchage (transfert de la feuille humide sur support absorbant) avant pressage.}
\end{popfigure}

\courssub{Les produits obtenus : adapter la pâte au produit}

Le même procédé de base donne des produits différents selon le \textbf{degré de désagrégation de la pâte} (finesse), l'\textbf{épaisseur de la feuille} (grammage) et les \textbf{traitements complémentaires} (désencrage, blanchiment, coloration, gaufrage).

\begin{center}
\begin{tabularx}{\linewidth}{|l|X|X|}
\hline
\rowcolor{popblueL} \textbf{Produit visé} & \textbf{Caractéristiques de la pâte et de la feuille} & \textbf{Traitements et façonnage spécifiques} \\
\hline
Serviettes jetables & Pâte fine, bien désagrégée. Feuilles minces (\SI{\approx 15-20}{g/m^2}), résistantes à la traction, peu absorbantes. & Désencrage poussé, blanchiment (peroxyde) pour blancheur. Pressage fort (lissage). Découpe en carrés (\SI{30x30}{cm} ou \SI{40x40}{cm}), pliage 1/4 ou 1/8. \\
\hline
Papier hygiénique & Pâte \textbf{très finement désagrégée} (brassage prolongé, raffinage). Feuilles \textbf{très fines} (\SI{\approx 13-17}{g/m^2}), souples, à haute capacité d'absorption (volume). & Désencrage et blanchiment soignés (blancheur \textgreater 80\% ISO). Séchage sur cylindre Yankee (industriel) pour le « crêpage » (micro-plis donnant le volume). Enroulage en rouleaux perforés (mandrin carton). \\
\hline
Papier-cadeau & Pâte plus épaisse, désagrégation modérée (fibres plus longues = meilleure tenue au pli). Grammage \SI{\approx 60-80}{g/m^2}. & Désencrage facultatif (souvent simple lavage). \textbf{Coloration en masse} (pigments dans la pulpe) ou impression/enduction surface. Gaufrage possible (motifs en relief). Découpe en feuilles (\SI{50x70}{cm}) ou rouleaux. \\
\hline
\end{tabularx}
\end{center}

Ainsi, c'est en jouant sur la finesse de la pulpe (durée/intensité du brassage/raffinage), le grammage (quantité de pulpe déposée) et les finitions que l'artisan ou l'industriel passe d'un produit à l'autre : le papier hygiénique exige les fibres les plus finement séparées et un crêpage pour le moelleux, tandis que le papier-cadeau tolère une pâte plus grossière et valorise l'aspect texturé et la couleur.

\courssub{Intérêt du recyclage du papier : quelques ordres de grandeur}

\begin{exemplebox}[Bilan environnemental du recyclage d'une tonne de papier]
Recycler \SI{1}{tonne} de papier et carton permet d'économiser, par rapport à la fabrication de pâte vierge (procédé kraft) :
\begin{itemize}
\item environ \textbf{17 arbres} (bois d'œuvre ou bois d'industrie, \SI{\approx 3,5}{m^3} de bois) ;
\item près de \SI{26000}{L} d'eau (soit \SI{26}{m^3}) — la pâte vierge consomme \SI{\approx 60}{m^3/t} contre \SI{\approx 10-15}{m^3/t} pour le recyclé ;
\item environ \SI{4000}{kWh} d'énergie (soit \SI{14,4}{GJ}) — économie de \SI{\approx 50}{\%} d'énergie primaire ;
\item \SI{\approx 3}{m^3} d'espace en décharge (volume des déchets évités) ;
\item une réduction significative des rejets polluants (COV, SOx, NOx, chlorures organiques liés au blanchiment kraft).
\end{itemize}
À l'échelle du Cameroun, où la consommation de papier scolaire, de bureau et d'emballage croît de \SI{\approx 5}{\%} par an, le recyclage constitue un levier majeur de préservation des forêts du bassin du Congo et de réduction de la pollution urbaine (décharges à ciel ouvert comme Nkolfoulou à Yaoundé ou PK10 à Douala).
\end{exemplebox}

\begin{savaistu}[Combien de fois peut-on recycler le papier ?]
Le papier peut être recyclé environ \textbf{5 à 7 fois}. À chaque cycle, le brassage mécanique (raffinage) et les traitements chimiques raccourcissent et endommagent les fibres de cellulose (coupures, cornification) ; lorsque la longueur moyenne des fibres devient trop faible (\textless \SI{0,5}{mm}), elles ne s'enchevêtrent plus assez pour former une feuille solide aux propriétés mécaniques acceptables. C'est pourquoi l'industrie papetière ajoute toujours une part de \textbf{fibres vierges} (pâte kraft) dans la pâte recyclée (souvent \SI{20}{\%} à \SI{50}{\%} selon la qualité visée), et pourquoi les papiers recyclés plusieurs fois finissent en produits de basse qualité (carton ondulé, boîtes d'œufs, isolant cellulose) avant de quitter le cycle matière pour être valorisés énergétiquement.
\end{savaistu}

\begin{exoresolu}[Exploiter des données quantitatives sur le recyclage du papier]
\textbf{Énoncé.} Une coopérative de jeunes de Bafoussam collecte et recycle \SI{2,5}{tonnes} de papiers et cartons mélangés par mois. On admet que le recyclage d'une tonne de papier économise 17 arbres et \SI{26000}{L} d'eau. On admet aussi qu'un arbre de plantation de papeterie (eucalyptus, pin) fournit en moyenne \SI{60}{kg} de papier sec après transformation.
\begin{enumerate}
\item Calcule le nombre d'arbres et le volume d'eau économisés par cette coopérative en un mois.
\item Quelle masse de papier neuf cette coopérative « remplace-t-elle » chaque mois par son activité ? Commenter le résultat.
\item Si la coopérative vend le papier hygiénique produit à \SI{500}{FCFA} le rouleau de \SI{200}{g}, quel est le chiffre d'affaires mensuel potentiel si tout le papier collecté est transformé en papier hygiénique (rendement transformation \SI{85}{\%}) ?
\end{enumerate}
\tcblower
\textbf{Solution détaillée.}
\begin{enumerate}
\item \textbf{Arbres économisés} : ${2,5} \times 17 = 42,5$, soit environ \textbf{42 à 43 arbres par mois} préservés.
\textbf{Eau économisée} : ${2,5} \times \SI{26000}{L} = \SI{65000}{L} = \SI{65}{m^3}$ d'eau par mois.
\item \textbf{Masse de papier neuf remplacée} : $42,5 \text{ arbres} \times \SI{60}{kg/arbre} = \SI{2550}{kg} \approx \SI{2,55}{tonnes}$.
\textbf{Commentaire} : On retrouve logiquement la masse de papier collectée (\SI{2,5}{t}). Cela confirme que le recyclage \textbf{substitue} quasi parfaitement des fibres récupérées aux fibres vierges : 1 tonne de papier recyclé évite la coupe et la transformation des arbres qui auraient fourni 1 tonne de papier neuf. En plus, on évite l'enfouissement de \SI{2,5}{t} de déchets.
\item \textbf{Masse de papier hygiénique produit} : \SI{2,5}{t} $\times$ \num{0,85} = \SI{2,125}{t} = \SI{2125}{kg}.
\textbf{Nombre de rouleaux} : $\frac{\SI{2125000}{g}}{\SI{200}{g/rouleau}} = \SI{10625}{rouleaux}$.
\textbf{Chiffre d'affaires} : $\SI{10625}{rouleaux} \times \SI{500}{FCFA} = \SI{5312500}{FCFA} \approx \SI{5,3}{millions FCFA}$ par mois.
\textit{Note : ce calcul ne tient pas compte des coûts (énergie, eau, main-d'œuvre, produits chimiques, amortissement, collecte), mais illustre le potentiel économique réel de la filière.}
\end{enumerate}
\end{exoresolu}

\begin{aretenir}[L'essentiel sur le recyclage du papier — Bac Terminale C/TI]
\begin{itemize}
\item Le recyclage du papier repose sur la \textbf{récupération des fibres de cellulose} des papiers usagés pour fabriquer de nouveaux produits (serviettes, papier hygiénique, papier-cadeau, carton).
\item La technique comporte \textbf{5 étapes ordonnées et incontournables} :
      \begin{enumerate}
      \item \textbf{Collecte et tri} (élimination corps étrangers, papiers souillés/plastifiés) ;
      \item \textbf{Mise en pâte} (trempage, brassage mécanique $\rightarrow$ pulpe fibreuse) ;
      \item \textbf{Désencrage} (flotation industrielle + lavage, tensioactifs) \textbf{et blanchiment éventuel} (peroxyde) selon blancheur visée ;
      \item \textbf{Formage sur tamis, pressage, séchage} (reformations liaisons hydrogène $\rightarrow$ feuille solide) ;
      \item \textbf{Façonnage} (découpe, pliage, gaufrage, enroulement, coloration).
      \end{enumerate}
\item Le produit final dépend de la \textbf{finesse de la pâte} (degré de désagrégation/raffinage), du \textbf{grammage} et des \textbf{traitements de surface} : papier hygiénique = pâte très fine + crêpage ; serviette = pâte fine + lissage ; papier-cadeau = pâte plus grossière + coloration/décor.
\item \textbf{Limitations} : papiers plastifiés, gras, thermiques, souillés toxiques = non recyclables. Fibres recyclables \textbf{5 à 7 fois} maximum (raccourcissement fibres).
\item \textbf{Avantages majeurs} : 1 t recyclée = 17 arbres + \SI{26}{m^3} eau + \SI{4000}{kWh} énergie économisés ; réduction déchets en décharge ; création d'emplois verts (collecte, tri, transformation artisanale/industrielle).
\end{itemize}
\end{aretenir}

\coursec{Le recyclage des déchets plastiques en bouteilles recyclées}

Après le papier, attaquons-nous au déchet le plus visible de nos villes : le plastique. À Douala, à Yaoundé ou à Bafoussam, les bouteilles vides d'eau minérale et de boissons gazeuses jonchent les caniveaux, bouchent les drains et provoquent des inondations en saison des pluies. Pourtant, ces bouteilles ne sont pas des déchets sans valeur : elles sont constituées d'un matériau remarquable, le \cle{PET}, qui peut être fondu et remodelé presque à volonté. Comprendre pourquoi certaines matières plastiques sont recyclables et d'autres non, puis décrire la chaîne technique qui transforme une bouteille usagée en bouteille neuve, tel est l'objectif de cette section.

\courssub{Nature des plastiques : pourquoi certains sont recyclables}

\textbf{Observation et questionnement.} Approche une bouteille d'eau vide d'une flamme (sans la toucher) : elle se ramollit, se déforme, puis durcit de nouveau en refroidissant. Recommence l'opération : elle se ramollit encore. Fais maintenant la même expérience avec le manche d'une vieille marmite en bakélite ou avec un interrupteur électrique cassé : il noircit, brûle et se carbonise, mais ne fond jamais. Comment expliquer cette différence fondamentale de comportement ?

\textbf{Interprétation.} Les plastiques sont des \cle{polymères}, c'est-à-dire de très longues molécules (macromolécules) obtenues par assemblage répété de petites unités appelées monomères, elles-mêmes issues du raffinage du \cle{pétrole}. On distingue deux grandes familles :

\begin{itemize}
\item les \cle{thermoplastiques} : leurs longues chaînes moléculaires sont simplement enchevêtrées, sans liaisons chimiques fortes entre elles. La chaleur augmente l'agitation moléculaire, les chaînes glissent les unes sur les autres : le matériau se \textbf{ramollit} et peut être \textbf{refondu et remodelé} de nombreuses fois. Ils sont donc \textbf{recyclables} ;
\item les \cle{thermodurcissables} : leurs chaînes sont reliées entre elles par des liaisons covalentes formant un réseau tridimensionnel rigide. Chauffé, ce réseau ne peut pas se défaire : le matériau se \textbf{décompose} sans fondre. Il est \textbf{irréversiblement durci}, donc non recyclable par fusion.
\end{itemize}

\begin{definition}[Thermoplastique et PET]
Un \cle{thermoplastique} est un polymère qui se ramollit sous l'effet de la chaleur et durcit en refroidissant, de façon réversible : il peut être fondu et remodelé plusieurs fois, ce qui le rend recyclable. Le \cle{PET} (polyéthylène téréphtalate) est un thermoplastique obtenu par polycondensation d'acide téréphtalique et d'éthylène glycol, deux dérivés du pétrole. Transparent, léger, solide et imperméable aux gaz, il est le matériau des bouteilles d'eau et de boissons gazeuses. Il porte le code de recyclage \textbf{1} (symbole triangulaire gravé sous les bouteilles).
\end{definition}

\begin{attention}[Tous les plastiques ne sont pas recyclables]
Erreur fréquente au Baccalauréat : écrire que « les déchets plastiques sont fondus puis recyclés » sans distinction. C'est faux. Seuls les \textbf{thermoplastiques} (PET, PE, PP, PVC\ldots) se re-fondent ; les \textbf{thermodurcissables} (bakélite, résines époxy, polyuréthanes durcis) ne se ramollissent jamais : ils brûlent ou se carbonisent. Le recyclage par fusion n'est donc possible que pour les thermoplastiques, et encore faut-il les \textbf{trier par type}, car des polymères différents fondus ensemble donnent un matériau hétérogène et fragile.
\end{attention}

\courssub{La chaîne technique de recyclage du PET en bouteilles}

\textbf{Situation-problème.} Une entreprise de la zone industrielle de Douala reçoit chaque jour plusieurs tonnes de bouteilles PET collectées dans les quartiers. Comment, concrètement, une bouteille sale, écrasée, avec son bouchon et son étiquette, devient-elle une bouteille neuve, propre et apte au contact alimentaire ? La réponse tient en cinq étapes successives.

\begin{methode}[Expliquer la technique de recyclage des plastiques en bouteilles en 5 étapes]
Pour rédiger une explication complète et bien ordonnée, suis toujours la chronologie de la chaîne :
\begin{enumerate}
\item \textbf{Collecte et tri} : rassembler les plastiques et séparer le PET (code 1) des autres polymères et des corps étrangers ;
\item \textbf{Lavage et séchage} : éliminer impuretés, colle et étiquettes, puis sécher ;
\item \textbf{Broyage} : réduire les bouteilles en petites paillettes appelées \emph{flakes} ;
\item \textbf{Fusion et extrusion} : fondre les paillettes et les transformer en filaments puis en granulés ;
\item \textbf{Moulage par soufflage} : chauffer les granulés et souffler la matière fondue dans un moule pour former la nouvelle bouteille.
\end{enumerate}
Conclure par le \textbf{contrôle qualité}, obligatoire pour le contact alimentaire.
\end{methode}

\textbf{Étape 1 : collecte et tri.} Les bouteilles usagées sont ramassées par des collecteurs (souvent des associations ou des ramasseurs informels), acheminées vers le centre de tri, puis séparées : le PET transparent des bouteilles est isolé des autres plastiques (PE des bidons, PVC, polystyrène), des bouchons (en PE, plus dense ou moins dense selon les cas, séparés par flottaison dans l'eau) et des déchets non plastiques. Le tri peut être manuel sur tapis roulant ou automatisé (tri optique par infrarouge). C'est l'étape \textbf{décisive} : la pureté du lot conditionne toute la qualité finale.

\textbf{Étape 2 : lavage et séchage.} Les bouteilles triées sont lavées à l'eau chaude (environ \SI{60}{\celsius} à \SI{80}{\celsius}) additionnée de détergent ou de soude diluée, afin d'éliminer les résidus de boisson, les graisses, la colle et les étiquettes. Elles sont ensuite essorées et séchées à l'air chaud. Un plastique mal lavé donnerait des granulés souillés, impropres à tout usage alimentaire.

\textbf{Étape 3 : broyage en paillettes.} Les bouteilles propres passent dans un broyeur à lames rotatives qui les découpe en petits fragments de quelques millimètres, appelés \cle{paillettes} ou \emph{flakes}. Une seconde séparation par flottaison peut avoir lieu ici : les paillettes de PET (densité voisine de \num{1,38}) coulent dans l'eau tandis que les fragments de bouchons en PE (densité voisine de \num{0,95}) flottent et sont évacués.

\textbf{Étape 4 : fusion et extrusion.} Les paillettes sont introduites dans une \cle{extrudeuse} : une vis sans fin tournant dans un cylindre chauffé (vers \SI{260}{\celsius} à \SI{280}{\celsius} pour le PET) les transporte, les comprime et les fond. La masse fondue est filtrée pour retenir les dernières particules solides, puis forcée à travers une \cle{filière} (sortie calibrée) d'où elle sort en filaments continus. Ces filaments sont refroidis dans un bain d'eau puis découpés en petits cylindres : les \cle{granulés} de PET recyclé, matière première de l'étape suivante.

\textbf{Étape 5 : moulage par soufflage.} Les granulés sont de nouveau fondus et injectés dans un moule pour former une \cle{préforme} (sorte d'éprouvette épaisse munie du goulot définitif). La préforme est réchauffée puis placée dans un moule en forme de bouteille ; de l'air comprimé y est insufflé : la matière ramollie épouse les parois du moule. Après refroidissement, on obtient la \textbf{bouteille recyclée}.

\textbf{Contrôle qualité.} Les bouteilles destinées au contact alimentaire doivent satisfaire à des \textbf{normes strictes} : absence de contaminants chimiques résiduels, de micro-organismes et de particules, résistance mécanique et étanchéité vérifiées. Le PET recyclé de qualité insuffisante est orienté vers des usages non alimentaires (fibres textiles, sanglage, rembourrage). Ce point est un savoir admis : aucune démonstration n'est exigible.

\begin{popfigure}
\begin{center}
\begin{tikzpicture}[
  font=\small,
  etape/.style={rectangle, rounded corners=3pt, draw=popblue, fill=popblueL, thick, minimum width=4.6cm, minimum height=1.05cm, align=center},
  detail/.style={font=\scriptsize\itshape, text=popmuted, align=center},
  fleche/.style={-{Stealth[length=3mm]}, very thick, poporange}
]
\node[etape, align=center] (e1) at (0,0){\textbf{1. Collecte et tri}\\séparation du PET (code 1)};
\node[etape, align=center] (e2) at (0,-1.9){\textbf{2. Lavage et séchage}\\eau chaude + détergent, élimination des étiquettes};
\node[etape, align=center] (e3) at (0,-3.8){\textbf{3. Broyage}\\réduction en paillettes (\emph{flakes})};
\node[etape, align=center] (e4) at (0,-5.7){\textbf{4. Fusion et extrusion}\\vers \SI{270}{\celsius}, filaments refroidis puis granulés};
\node[etape, align=center] (e5) at (0,-7.6){\textbf{5. Moulage par soufflage}\\préforme soufflée dans un moule $\to$ bouteille neuve};
\draw[fleche] (e1) -- (e2);
\draw[fleche] (e2) -- (e3);
\draw[fleche] (e3) -- (e4);
\draw[fleche] (e4) -- (e5);
\node[detail] at (5.1,-1.9) {pureté du lot};
\node[detail] at (5.1,-3.8) {fragments de quelques mm};
\node[detail] at (5.1,-5.7) {filtration des résidus};
\node[detail, align=center] at (5.1,-7.6){contrôle qualité (normes\\de contact alimentaire)};
\draw[-{Stealth[length=2mm]}, popmuted, dashed] (3.4,-1.9) -- (2.35,-1.9);
\draw[-{Stealth[length=2mm]}, popmuted, dashed] (3.4,-3.8) -- (2.35,-3.8);
\draw[-{Stealth[length=2mm]}, popmuted, dashed] (3.4,-5.7) -- (2.35,-5.7);
\draw[-{Stealth[length=2mm]}, popmuted, dashed] (3.4,-7.6) -- (2.35,-7.6);
\end{tikzpicture}
\end{center}
\legende{Organigramme des cinq étapes du recyclage des bouteilles PET en bouteilles recyclées. Chaque étape prépare la suivante : un tri ou un lavage insuffisant compromet la qualité du produit final.}
\end{popfigure}

\courssub{Le cœur du procédé : l'extrudeuse}

L'extrudeuse est la machine clé de la plasturgie : c'est elle qui réalise le passage de la paillette solide au granulé homogène. Son fonctionnement repose sur trois éléments : la \textbf{vis sans fin} qui transporte et comprime la matière, les \textbf{colliers chauffants} qui la fondent progressivement, et la \textbf{filière} qui lui donne sa forme de sortie.

\begin{popfigure}
\begin{center}
\begin{tikzpicture}[font=\small, scale=0.95]
% Cylindre de l'extrudeuse
\draw[thick, popdark, fill=popcream] (0,-0.9) rectangle (8.4,0.9);
% Trémie
\draw[thick, popdark, fill=popblueL] (0.7,0.9) -- (0.3,2.1) -- (1.9,2.1) -- (1.5,0.9) -- cycle;
\node[font=\scriptsize] at (1.1,2.45) {Trémie};
\node[font=\scriptsize, text=popmuted] at (1.1,1.55) {paillettes};
\draw[-{Stealth[length=2mm]}, popblue, thick] (1.1,1.3) -- (1.1,0.95);
% Vis sans fin
\draw[very thick, poporange] (0.3,0) -- (8.1,0);
\foreach \x in {0.5,1.1,1.7,2.3,2.9,3.5,4.1,4.7,5.3,5.9,6.5,7.1,7.7}{
  \draw[thick, poporange] (\x,-0.62) .. controls (\x+0.25,-0.1) and (\x+0.25,0.1) .. (\x,0.62);
}
\node[font=\scriptsize, text=poporange] at (2.2,-1.25) {Vis sans fin (transport + compression)};
% Zones de chauffe
\draw[thick, poppink, fill=poppinkL] (2.6,-0.9) rectangle (4.4,0.9);
\draw[thick, poppink, fill=poppinkL] (4.7,-0.9) rectangle (6.5,0.9);
\node[font=\scriptsize, text=poppink] at (3.5,1.25) {Zone de chauffe 1};
\node[font=\scriptsize, text=poppink] at (5.6,1.25) {Zone de chauffe 2};
\node[font=\scriptsize, text=poppink] at (5.6,-1.25) {$\approx$ \SI{270}{\celsius} : fusion du PET};
% Filière
\draw[thick, popdark, fill=popgreenL] (8.4,-0.5) -- (9.1,-0.18) -- (9.1,0.18) -- (8.4,0.5) -- cycle;
\node[font=\scriptsize] at (8.75,0.85) {Filière};
% Filaments
\draw[very thick, popgreen] (9.1,0.08) -- (10.6,0.08);
\draw[very thick, popgreen] (9.1,-0.08) -- (10.6,-0.08);
\node[font=\scriptsize, text=popgreen] at (10.0,0.5) {Filaments};
\draw[-{Stealth[length=2mm]}, popgreen, thick] (10.6,0) -- (11.3,0);
\node[font=\scriptsize, text=popgreen] at (11.9,0) {granulés};
% Sens de rotation
\draw[-{Stealth[length=2mm]}, poporange, thick] (0.1,1.3) arc (150:30:0.5);
\node[font=\scriptsize, text=poporange] at (0.75,1.75) {rotation};
\end{tikzpicture}
\end{center}
\legende{Schéma simplifié d'une extrudeuse. Les paillettes de PET introduites par la trémie sont saisies par la vis sans fin en rotation, fondues dans les zones de chauffe (environ \SI{270}{\celsius}), puis forcées à travers la filière d'où elles sortent en filaments continus, refroidis puis découpés en granulés.}
\end{popfigure}

\courssub{Exemples chiffrés et applications}

\begin{exemplebox}[Du déchet au vêtement : un exemple chiffré]
Environ \textbf{25 à 30 bouteilles} en PET recyclées fournissent assez de matière pour fabriquer \textbf{un vêtement polaire} : les paillettes fondues sont alors extrudées en fibres textiles fines au lieu d'être soufflées en bouteilles. De même, une tonne de PET recyclé évite l'extraction d'environ \SI{800}{kg} de pétrole brut et réduit les émissions de gaz à effet de serre par rapport à la production de plastique vierge. Recycler, c'est donc économiser une ressource fossile non renouvelable.
\end{exemplebox}

\begin{savaistu}[Le plastique, une richesse qui dort dans nos rues]
Au Cameroun, plusieurs entreprises et coopératives collectent les bouteilles plastiques usagées dans les grandes villes : une partie est \textbf{exportée} (vers l'Asie notamment) sous forme de paillettes ou de balles compressées, l'autre est \textbf{transformée localement} en granulés, en fibres de rembourrage ou en objets moulés. Cette filière fait vivre des milliers de ramasseurs et trieurs, souvent des jeunes et des femmes. À l'échelle d'un quartier, une simple opération de collecte organisée peut retirer plusieurs tonnes de plastique des caniveaux avant la saison des pluies, réduisant ainsi les inondations et les gîtes larvaires des moustiques vecteurs du paludisme. Recycler le plastique, c'est aussi protéger la santé publique.
\end{savaistu}

\begin{exoresolu}[Expliquer le recyclage du PET (type Bac)]
\textbf{Énoncé.} Un élève de Terminale C affirme à ses camarades : « Toutes les bouteilles vides de notre quartier peuvent redevenir des bouteilles neuves. » (a) Cette affirmation est-elle exacte pour toutes les bouteilles ? Justifier. (b) Décrire, en cinq étapes ordonnées, la technique qui permet de transformer les bouteilles en PET en bouteilles recyclées. (c) Pourquoi un contrôle qualité est-il indispensable pour les bouteilles destinées à contenir des boissons ?
\tcblower
\textbf{Solution détaillée.}
(a) L'affirmation n'est exacte que pour les bouteilles en matière \textbf{thermoplastique}, comme le PET (code 1) : celles-ci se ramollissent à la chaleur et peuvent être refondues et remodelées. Les récipients en plastique \textbf{thermodurcissable}, dont les chaînes moléculaires sont liées en réseau rigide, ne fondent pas : chauffés, ils se décomposent. Ils ne peuvent donc pas être recyclés par fusion.

(b) La technique comporte cinq étapes :
\begin{enumerate}
\item \textbf{Collecte et tri} : les bouteilles sont ramassées et le PET est séparé des autres plastiques, des bouchons et des corps étrangers ;
\item \textbf{Lavage et séchage} : lavage à l'eau chaude additionnée de détergent pour éliminer impuretés, colle et étiquettes, puis séchage ;
\item \textbf{Broyage} : les bouteilles propres sont réduites en paillettes (\emph{flakes}) ;
\item \textbf{Fusion et extrusion} : les paillettes sont fondues vers \SI{270}{\celsius} dans une extrudeuse, filtrées, puis transformées en filaments refroidis et découpés en granulés ;
\item \textbf{Moulage par soufflage} : les granulés fondus forment des préformes qui, réchauffées, sont soufflées à l'air comprimé dans des moules pour donner les nouvelles bouteilles.
\end{enumerate}

(c) Le contrôle qualité est indispensable car la bouteille sera au \textbf{contact d'un aliment} (eau, jus) : il faut garantir l'absence de contaminants chimiques résiduels, de micro-organismes et de particules, ainsi que la résistance mécanique et l'étanchéité du récipient, conformément à des normes strictes de sécurité sanitaire.
\end{exoresolu}

\begin{exercice}[À toi de jouer]
Dans un centre de tri de Yaoundé, on reçoit un lot mêlant bouteilles d'eau en PET, bidons d'huile en PE, vieux interrupteurs en bakélite et barquettes en polystyrène.
\begin{enumerate}
\item Quels éléments de ce lot peuvent être recyclés par fusion ? Justifier.
\item Pourquoi est-il nécessaire de séparer le PET des autres plastiques avant la fusion ?
\item Remettre dans l'ordre les étapes suivantes : extrusion ; broyage ; soufflage ; tri ; lavage.
\end{enumerate}
\end{exercice}

\begin{corrige}[Exercice : À toi de jouer]
\begin{enumerate}
\item Les bouteilles en PET, les bidons en PE et les barquettes en polystyrène sont des \textbf{thermoplastiques} : ils se ramollissent à la chaleur et sont recyclables par fusion. Les interrupteurs en bakélite sont \textbf{thermodurcissables} : leurs chaînes moléculaires forment un réseau rigide qui ne fond pas ; ils se décomposent à la chaleur et ne sont pas recyclables par ce procédé.
\item Parce que des polymères de natures différentes, fondus ensemble, sont \textbf{incompatibles} : ils donnent un matériau hétérogène, fragile et sans valeur commerciale. Le tri garantit la pureté du lot et donc la qualité des granulés.
\item Ordre correct : \textbf{tri} $\rightarrow$ \textbf{lavage} $\rightarrow$ \textbf{broyage} $\rightarrow$ \textbf{extrusion} $\rightarrow$ \textbf{soufflage}.
\end{enumerate}
\end{corrige}

\begin{aretenir}[L'essentiel de la section]
\begin{itemize}
\item Les plastiques sont des \cle{polymères} issus du pétrole ; seuls les \cle{thermoplastiques} (dont le PET des bouteilles) se re-fondent et sont recyclables ; les thermodurcissables ne se ramollissent jamais.
\item Le recyclage du PET en bouteilles suit cinq étapes : \textbf{tri} $\rightarrow$ \textbf{lavage-séchage} $\rightarrow$ \textbf{broyage en paillettes} $\rightarrow$ \textbf{fusion-extrusion en granulés} $\rightarrow$ \textbf{moulage par soufflage}.
\item L'\cle{extrudeuse} (vis sans fin + zones de chauffe + filière) transforme les paillettes en filaments puis en granulés.
\item Les bouteilles recyclées à usage alimentaire exigent un \textbf{contrôle qualité} strict (savoir admis).
\end{itemize}
\end{aretenir}

Les bouteilles ne sont pas le seul destin possible des déchets plastiques : fondus et mélangés à du sable, ils peuvent aussi devenir des matériaux de construction. C'est l'objet de la section suivante, consacrée à la transformation des déchets plastiques en pavés et aux avantages généraux de la valorisation des déchets.

\coursec{La transformation des déchets plastiques en pavés et les avantages de la valorisation}

Dans la section précédente, nous avons vu comment les déchets plastiques rigides (bouteilles en PET) peuvent être lavés, broyés et refondus pour donner naissance à de nouvelles bouteilles recyclées. Mais que faire des plastiques souillés, des sachets fins, des emballages mélangés qui ne peuvent pas être revalorisés en bouteilles ? Faut-il les abandonner dans les rues de Douala ou de Yaoundé, où ils bouchent les caniveaux et provoquent des inondations ? Non. Une innovation simple et ingénieuse permet de les transformer en \cle{pavés} de construction, plus résistants et moins chers que les pavés classiques en ciment. C'est cette technique, puis les avantages généraux de la valorisation des déchets, que nous allons étudier maintenant.

\courssub{Principe de la fabrication des pavés plastiques}

\textbf{Observation de départ.} Dans une cour d'école de Bafoussam, des pavés gris foncé, posés depuis plusieurs années, ne présentent ni fissure ni trace d'usure, alors que le pavage en ciment d'une cour voisine s'effrite. Or ces pavés ont été fabriqués à partir de... sachets plastiques ramassés dans les rues ! Comment un déchet peut-il devenir un matériau de construction performant ?

\textbf{Idée directrice.} Dans un pavé classique, le \cle{ciment} joue le rôle de \emph{liant} : il enrobe les grains de sable et durcit en prenant l'eau, solidarisant l'ensemble. Or le plastique fondu possède exactement les mêmes propriétés : liquide à chaud, il enrobe les grains de sable, puis durcit en refroidissant. On peut donc \emph{remplacer le ciment par du plastique fondu} : le plastique devient le liant, le sable reste le squelette granulaire (le \emph{granulat}).

\begin{definition}[Pavé plastique]
Un \cle{pavé plastique} (ou pavé écologique) est un \textbf{matériau de construction obtenu par mélange de plastique fondu et de sable}, coulé dans un moule puis refroidi. Dans ce matériau composite, le plastique solidifié joue le rôle de \cle{liant} (à la place du ciment) et le sable celui de \cle{granulat}.
\end{definition}

\begin{propriete}[Propriétés des pavés plastiques]
Les pavés plastique-sable présentent trois propriétés remarquables, vérifiables expérimentalement :
\begin{itemize}
\item une \textbf{grande résistance à la compression}, souvent supérieure à celle des pavés en ciment ;
\item une \textbf{imperméabilité} totale : le plastique étant hydrophobe, le pavé n'absorbe pas l'eau et ne se dégrade pas sous l'effet des pluies ;
\item une \textbf{grande durabilité} : insensible à l'humidité, aux champignons et aux termites, le pavé ne s'effrite pas avec le temps.
\end{itemize}
\end{propriete}

\courssub{Les cinq étapes de fabrication des pavés plastiques}

\begin{methode}[Décrire la fabrication des pavés plastiques en 5 étapes]
Pour décrire correctement cette technique (exigible au Bac), retiens la logique suivante : \textbf{préparer le plastique, le fondre, lui donner un squelette de sable, lui donner une forme, le laisser durcir}.
\begin{enumerate}
\item \textbf{Collecte, tri et lavage} des déchets plastiques (sachets, bouteilles, emballages) ;
\item \textbf{Découpe et fusion} en morceaux dans une cuve chauffée à température contrôlée ;
\item \textbf{Mélange} du plastique fondu avec du sable propre et sec ;
\item \textbf{Coulage} de la pâte dans des moules et \textbf{compactage} ;
\item \textbf{Refroidissement, démoulage et séchage} des pavés.
\end{enumerate}
\end{methode}

Détaillons chaque étape.

\textbf{Étape 1 : collecte, tri et lavage.} Les déchets plastiques (sachets, bouteilles vides, emballages) sont ramassés dans les rues, les marchés et les poubelles. On les \emph{trie} pour éliminer les corps étrangers (métaux, papiers, matières organiques), puis on les \emph{lave} à l'eau savonneuse et on les laisse sécher. Un plastique propre garantit une fusion homogène et un pavé sans défaut.

\textbf{Étape 2 : découpe et fusion.} Les plastiques propres sont découpés en petits morceaux (ce qui accélère la fusion), puis introduits dans une \emph{cuve chauffée}. On porte la température à environ \SI{170}{\degreeCelsius} à \SI{250}{\degreeCelsius}, température suffisante pour fondre les plastiques courants (PE, PP) mais \emph{contrôlée} pour éviter leur combustion. On obtient une pâte visqueuse homogène.

\textbf{Étape 3 : mélange avec le sable.} On ajoute progressivement du \emph{sable propre, sec et tamisé} au plastique fondu, en malaxant vigoureusement. Une proportion couramment utilisée est d'environ \textbf{1 volume de plastique pour 3 à 4 volumes de sable} (soit environ 20 à 25 \% de plastique en masse). Le plastique fondu enrobe chaque grain de sable : c'est lui qui, en refroidissant, cimentera l'ensemble.

\textbf{Étape 4 : coulage et compactage.} La pâte chaude est versée dans des \emph{moules} métalliques de forme parallélépipédique, puis \emph{compactée} (à la presse ou au pilon) pour chasser les bulles d'air et densifier le matériau. Un pavé bien compacté est un pavé résistant.

\textbf{Étape 5 : refroidissement, démoulage et séchage.} On laisse refroidir les moules (parfois en les plongeant dans l'eau), puis on \emph{démoule} les pavés et on les laisse sécher et durcir à l'air libre pendant quelques jours. Les pavés sont alors prêts à être posés.

\begin{popfigure}
\begin{center}
\begin{tikzpicture}[
  font=\small,
  etape/.style={draw=popblue, fill=popblueL, rounded corners=3pt, text width=3.6cm, align=center, minimum height=1.2cm, thick},
  fleche/.style={-{Stealth[length=3mm]}, very thick, poporange}
]
\node[etape, align=center] (e1) at (0,0){\textbf{1. Collecte, tri et lavage}\\ sachets, bouteilles};
\node[etape, align=center] (e2) at (5.0,0){\textbf{2. Découpe et fusion}\\ cuve chauffée ($\approx$ 170--250 \textdegree C)};
\node[etape, align=center] (e3) at (10.0,0){\textbf{3. Mélange avec le sable}\\ 1 vol. plastique / 3--4 vol. sable};
\node[etape, align=center] (e4) at (5.0,-3.0){\textbf{4. Coulage en moules}\\ et compactage};
\node[etape, fill=popgreenL, draw=popgreen, align=center] (e5) at (10.0,-3.0){\textbf{5. Refroidissement,}\\ démoulage et séchage};
\draw[fleche] (e1) -- (e2);
\draw[fleche] (e2) -- (e3);
\draw[fleche] (e3.south) |- ++(0,-0.6) -| (e4.north);
\draw[fleche] (e4) -- (e5);
\end{tikzpicture}
\end{center}
\legende{Organigramme des cinq étapes de fabrication des pavés plastiques. Le plastique fondu remplace le ciment comme liant ; le sable constitue le granulat.}
\end{popfigure}

\begin{attention}[Fusion du plastique : température contrôlée et espace ventilé]
Deux erreurs sont fréquentes, dans la pratique comme dans les copies :
\begin{itemize}
\item \textbf{Brûler le plastique au lieu de le fondre.} Si la température n'est pas contrôlée, le plastique se décompose en dégageant des \textbf{fumées toxiques} (dont des dioxines). La fusion doit donc se faire à température maîtrisée, jamais à flamme vive incontrôlée.
\item \textbf{Travailler en espace clos.} Même à bonne température, le plastique fondu dégage des vapeurs irritantes : l'opération doit se dérouler \textbf{en espace ventilé}, avec des équipements de protection (gants, masques).
\end{itemize}
À l'écrit, ne dis jamais qu'on « brûle » les déchets plastiques : on les \emph{fond}. La combustion des plastiques à l'air libre est une source de pollution, pas une valorisation.
\end{attention}

\courssub{Vérification expérimentale : essai de compression}

\textbf{Problème scientifique.} Affirmer que les pavés plastiques sont « plus résistants » ne suffit pas : en science, toute propriété annoncée doit être \emph{vérifiée par l'expérience}. Comment comparer objectivement la résistance d'un pavé plastique-sable et d'un pavé classique en ciment ?

\begin{experience}[Essai de compression comparatif]
\textbf{Matériel :} pavés plastique-sable et pavés en ciment de mêmes dimensions ; presse hydraulique (machine d'essai de compression) ; dispositif de mesure de la force de rupture.

\textbf{Protocole :} on place chaque pavé entre les plateaux de la presse et on applique une force croissante jusqu'à la rupture. On note la \emph{contrainte de rupture} (force maximale rapportée à la surface), exprimée en mégapascals (MPa). On répète l'essai sur plusieurs pavés de chaque type et on calcule la moyenne.

\textbf{Résultats (données d'un atelier pilote) :}
\begin{center}
\begin{tabular}{|l|c|c|}
\hline
\rowcolor{popblueL} \textbf{Type de pavé} & \textbf{Résistance moyenne (MPa)} & \textbf{Absorption d'eau} \\
\hline
Pavé classique en ciment & 18 & forte \\
\hline
Pavé plastique-sable & 26 & quasi nulle \\
\hline
\end{tabular}
\end{center}

\textbf{Interprétation :} le pavé plastique-sable rompt sous une contrainte moyenne de \SI{26}{MPa}, contre \SI{18}{MPa} pour le pavé en ciment : il est donc environ 1,4 fois plus résistant à la compression. De plus, sa quasi-imperméabilité le protège de l'érosion par les pluies.

\textbf{Conclusion :} les propriétés annoncées (résistance, imperméabilité, durabilité) sont confirmées expérimentalement ; le plastique fondu est un liant au moins aussi performant que le ciment pour ce type d'ouvrage.
\end{experience}

\begin{popfigure}
\begin{center}
\begin{tikzpicture}
\begin{axis}[
  ybar, bar width=1.6cm,
  width=0.75\linewidth, height=6.5cm,
  ymin=0, ymax=32,
  ylabel={Résistance à la compression (MPa)},
  symbolic x coords={Pavé ciment, Pavé plastique-sable},
  xtick=data,
  nodes near coords,
  every node near coord/.append style={font=\small\bfseries},
  ymajorgrids=true, grid style={popline!40},
  axis line style={popdark},
]
\addplot[fill=popmuted!60, draw=popdark] coordinates {(Pavé ciment,18)};
\addplot[fill=popgreen!70, draw=popdark] coordinates {(Pavé plastique-sable,26)};
\end{axis}
\end{tikzpicture}
\end{center}
\legende{Comparaison de la résistance à la compression d'un pavé classique en ciment (\SI{18}{MPa}) et d'un pavé plastique-sable (\SI{26}{MPa}). Le pavé plastique est environ 1,4 fois plus résistant (données expérimentales fictives à titre d'illustration).}
\end{popfigure}

\begin{exemplebox}[Un pavé plus résistant et moins cher]
Un atelier de valorisation produit un pavé plastique-sable pour un coût de revient d'environ 75 FCFA (sable, collecte et chauffage compris), contre environ 125 FCFA pour un pavé en ciment de mêmes dimensions, le ciment étant un matériau coûteux et souvent importé. Avec une résistance de \SI{26}{MPa} contre \SI{18}{MPa}, le pavé plastique est donc à la fois \textbf{plus résistant} (environ 40 \% de plus) et \textbf{moins cher} (environ 40 \% de moins) : un double avantage économique et technique.
\end{exemplebox}

\begin{exoresolu}[Calculer une quantité de matière première]
\textbf{Énoncé.} Une coopérative de jeunes de Buea fabrique des pavés contenant 25 \% de plastique en masse. Chaque pavé a une masse de \SI{4,0}{kg}. La coopérative reçoit une commande de 2 000 pavés pour paver la cour d'un lycée.
\begin{enumerate}
\item Quelle masse de déchets plastiques la coopérative doit-elle collecter ?
\item Un sachet plastique vide a une masse moyenne de \SI{5}{g}. À combien de sachets cela correspond-il ?
\end{enumerate}
\tcblower
\textbf{Solution détaillée.}
\begin{enumerate}
\item Masse totale des pavés : $2\,000 \times 4{,}0 = 8\,000$ kg. Masse de plastique nécessaire (25 \%) :
\[ m_{\text{plastique}} = 0{,}25 \times 8\,000 = 2\,000 \ \text{kg} = 2 \ \text{tonnes}. \]
\item Nombre de sachets correspondant :
\[ N = \frac{2\,000\,000 \ \text{g}}{5 \ \text{g}} = 400\,000 \ \text{sachets}. \]
\end{enumerate}
\textbf{Conclusion :} cette seule commande permet de retirer de l'environnement \textbf{2 tonnes de plastique}, soit environ \textbf{400 000 sachets} qui n'iront ni boucher les caniveaux, ni polluer les sols. La valorisation transforme un problème d'assainissement en ressource économique.
\end{exoresolu}

\begin{savaistu}[Des cours d'écoles pavées avec des déchets]
En Afrique, la filière « pavés plastiques » est portée par de jeunes entrepreneurs. Au \textbf{Kenya}, l'ingénieure Nzambi Matee a fondé à Nairobi une entreprise qui transforme chaque jour des centaines de kilogrammes de déchets plastiques en pavés deux à trois fois plus résistants que le béton, utilisés pour paver des cours d'écoles et des trottoirs. Au \textbf{Cameroun}, des initiatives similaires voient le jour à Douala, Yaoundé et Bafoussam : des jeunes collectent sachets et bouteilles, fabriquent des pavés et des carreaux, et créent ainsi des emplois tout en assainissant leurs quartiers. La biotechnologie environnementale n'est pas un luxe des pays riches : c'est une opportunité de développement local.
\end{savaistu}

\courssub{Les avantages de la valorisation des déchets}

La fabrication des pavés plastiques illustre, à elle seule, presque tous les avantages de la valorisation des déchets. Dressons-en le bilan général, valable aussi pour le recyclage du papier, des bouteilles et la production de biocarburants étudiés dans les sections précédentes.

\begin{aretenir}[Avantages de la valorisation des déchets]
\begin{itemize}
\item \textbf{Réduction de la pollution} : moins de déchets abandonnés dans la nature, les sols, les rivières et l'océan ; moins de plastiques brûlés à l'air libre (fumées toxiques).
\item \textbf{Assainissement des villes} : la collecte des déchets débouche les caniveaux, ce qui réduit les inondations et limite les gîtes larvaires des moustiques, donc le paludisme.
\item \textbf{Création d'emplois et de revenus} : collecte, tri, transformation et vente génèrent des activités économiques, notamment pour les jeunes et les femmes.
\item \textbf{Économie de matières premières} : recycler le papier préserve les forêts ; recycler le plastique économise le pétrole ; le sable et les déchets remplacent le ciment importé.
\item \textbf{Production d'énergie} : biocarburants, biogaz et combustion contrôlée des résidus fournissent une énergie renouvelable.
\item \textbf{Lutte contre le réchauffement climatique} : le recyclage évite des émissions de gaz à effet de serre (méthane des décharges, CO$_2$ de la fabrication du ciment et du papier neuf) et préserve les forêts, puits de carbone.
\end{itemize}
\end{aretenir}

\begin{attention}[Ne pas confondre les formulations]
À l'écrit, distingue bien : « \emph{enfouissement} » et « \emph{incinération} » (élimination sans valorisation, sources de pollution) d'une part, « \emph{recyclage} » et « \emph{valorisation} » (transformation du déchet en nouvelle ressource ou en énergie) d'autre part. Seule la valorisation procure les avantages listés ci-dessus. Autre piège : le pavé plastique n'est pas « biodégradable » — c'est justement sa durabilité qui fait sa qualité de matériau de construction ; sa valeur écologique vient du fait qu'il \emph{capt durablement} des déchets qui auraient pollué l'environnement.
\end{attention}

\begin{exercice}[Pour t'entraîner]
Une commune souhaite paver 500 m$^2$ de trottoirs. On hésite entre des pavés en ciment (125 FCFA l'unité, résistance \SI{18}{MPa}) et des pavés plastique-sable (75 FCFA l'unité, résistance \SI{26}{MPa}). Il faut 40 pavés par mètre carré.
\begin{enumerate}
\item Calcule le coût total de chaque option et l'économie réalisée avec les pavés plastiques.
\item Sachant qu'un pavé contient \SI{1}{kg} de plastique, quelle masse de déchets plastiques sera valorisée ?
\item Cite trois avantages de ce choix pour la commune.
\end{enumerate}
\end{exercice}

\begin{corrige}[Exercice : pavage des trottoirs]
\begin{enumerate}
\item Nombre de pavés : $500 \times 40 = 20\,000$ pavés. Coût en ciment : $20\,000 \times 125 = 2\,500\,000$ FCFA. Coût en pavés plastiques : $20\,000 \times 75 = 1\,500\,000$ FCFA. Économie : $2\,500\,000 - 1\,500\,000 = 1\,000\,000$ FCFA, soit 40 \% du budget.
\item Masse de plastique valorisée : $20\,000 \times 1 = 20\,000$ kg, soit \textbf{20 tonnes} de déchets plastiques retirées de l'environnement.
\item Avantages : réduction de la pollution et assainissement de la ville (caniveaux débouchés, moins d'inondations) ; économie financière et de matières premières (moins de ciment importé) ; création d'emplois locaux (collecte, fabrication, pose) ; contribution à la lutte contre le réchauffement climatique.
\end{enumerate}
\end{corrige}

Ainsi, le déchet plastique, hier symbole de la pollution urbaine, devient une ressource : un liant pour construire nos villes, un revenu pour nos jeunes, un geste pour la planète. Il reste une condition essentielle au succès de toutes ces techniques : l'adhésion de la population. C'est pourquoi la dernière section sera consacrée à la \emph{sensibilisation de la communauté au recyclage}.

\coursec{Sensibilisation de la communauté au recyclage}

Nous venons de voir, dans la section précédente, comment les déchets plastiques peuvent être transformés en pavés et, plus largement, quels sont les avantages sanitaires, économiques et écologiques de la valorisation des déchets. Mais une question demeure : à quoi servent les meilleures techniques de recyclage si personne ne trie ses déchets, si personne n'apporte ses plastiques au point de collecte, si les ménages continuent de brûler leurs ordures à l'air libre ? La réponse est simple : à rien. La technique sans l'adhésion de la population reste lettre morte. C'est pourquoi la dernière étape de notre leçon, loin d'être un simple supplément moral, est une véritable \cle{compétence} : savoir \cle{informer la communauté} sur la nécessité de transformer et de recycler les déchets. C'est un savoir-faire exigible au Baccalauréat, et surtout une mission citoyenne qui te concerne directement, toi, élève de Terminale.

\courssub{Pourquoi faut-il informer la communauté ?}

\textbf{Observation et mise en évidence du problème.}\par

Prenons une situation que tu connais bien : un quartier de Yaoundé ou de Douala après une forte pluie. Les caniveaux sont obstrués par des sachets plastiques, l'eau stagne, les moustiques pullulent, et les services de santé enregistrent une recrudescence des cas de paludisme et de choléra. Ce n'est pas un hasard : c'est la conséquence directe de déchets \emph{non valorisés} et mal gérés.

\begin{experience}[Document : les conséquences sanitaires des déchets non valorisés]
\textbf{Document.} On compare deux quartiers d'une même ville camerounaise : le quartier A, où une association organise le tri et la collecte des déchets, et le quartier B, où les déchets sont abandonnés dans la nature et brûlés à l'air libre. On relève sur une année les données suivantes (valeurs indicatives) :

\begin{center}
\begin{tabularx}{\linewidth}{|l|X|X|}
\hline
\rowcolor{popblueL} \textbf{Indicateur} & \textbf{Quartier A (tri + collecte)} & \textbf{Quartier B (abandon + brûlage)} \\
\hline
Cas de paludisme recensés (pour \num{1000} habitants) & 120 & 260 \\
\hline
Cas d'infections respiratoires (pour \num{1000} habitants) & 45 & 150 \\
\hline
Inondations de rues après pluie & rares & fréquentes \\
\hline
\end{tabularx}
\end{center}

\textbf{Interprétation.} Les déchets abandonnés retiennent l'eau de pluie (gîtes larvaires pour les moustiques, vecteurs du paludisme), obstruent les caniveaux (inondations, eaux stagnantes favorisant le choléra), et leur combustion à l'air libre dégage des fumées toxiques (dioxines, particules fines) responsables d'infections respiratoires.

\textbf{Conclusion.} Les déchets non valorisés polluent les \cle{sols}, les \cle{eaux} et l'\cle{air}, et constituent un réservoir d'agents pathogènes. Informer la communauté n'est donc pas un luxe : c'est une mesure de \cle{santé publique}.
\end{experience}

Retenons l'essentiel de cette démarche :

\begin{itemize}
\item Les déchets plastiques jetés dans la nature mettent des siècles à se dégrader (un sachet plastique : environ 100 à 400 ans ; une bouteille en PET : environ 450 ans) ; ils s'accumulent dans les sols et les cours d'eau.
\item Les eaux stagnantes piégées dans les détritus sont des gîtes larvaires pour \emph{Anopheles} (paludisme) et \emph{Aedes} (dengue, fièvre jaune).
\item La combustion des plastiques à l'air libre libère des gaz toxiques (dioxines, furanes, monoxyde de carbone \ce{CO}) dangereux pour les poumons.
\item Les déchets organiques en décomposition anarchique attirent les rongeurs et les mouches, vecteurs de maladies diarrhéiques (choléra, typhoïde).
\end{itemize}

\begin{aretenir}[La nécessité de la sensibilisation]
La valorisation des déchets ne peut réussir que si la population \emph{comprend} les dangers des déchets non valorisés et \emph{adopte} les bons comportements. Sensibiliser, c'est transformer un savoir scientifique en changement de comportement collectif. C'est le maillon qui relie la biotechnologie à la société.
\end{aretenir}

\courssub{Les cibles de la sensibilisation}

Une campagne efficace ne s'adresse pas « à tout le monde » de la même façon. Elle identifie des \cle{cibles} précises et adapte son message à chacune :

\begin{center}
\begin{tabularx}{\linewidth}{|l|X|X|}
\hline
\rowcolor{popblueL} \textbf{Cible} & \textbf{Pourquoi elle est prioritaire} & \textbf{Message adapté} \\
\hline
Les ménages & Ils produisent la majorité des déchets domestiques ; c'est à la maison que commence le tri & « Trie tes déchets à la source : organiques, plastiques, papiers » \\
\hline
Les écoles & Les élèves apprennent vite et transmettent à leurs parents ; ils sont les citoyens de demain & Clubs environnement, concours de recyclage, ateliers pratiques \\
\hline
Les marchés & Très forte production de déchets organiques et de sachets plastiques & « Vends propre, jette juste » ; compostage des déchets de légumes \\
\hline
Les autorités locales (mairies, chefs de quartier) & Elles décident des points de collecte, des budgets et des règlements & Plaidoyer pour l'installation de bacs de tri et de points de collecte \\
\hline
\end{tabularx}
\end{center}

\courssub{Les outils de communication}

Chaque outil a ses forces. Le bon communicateur les combine :

\begin{itemize}
\item \textbf{Les affiches} : placées aux carrefours, aux marchés, dans les écoles et les centres de santé, elles délivrent un message visuel simple, compréhensible même par les personnes peu scolarisées.
\item \textbf{Les spots radio en langues locales} (ewondo, douala, fulfuldé, bamiléké, pidgin\ldots) : la radio reste le média le plus écouté au Cameroun ; un spot de 30 secondes dans la langue du quartier touche les personnes âgées et les milieux ruraux.
\item \textbf{Le théâtre forum} : une courte pièce jouée dans la rue, montrant une famille qui brûle ses plastiques et tombe malade, puis qui adopte le tri ; le public est invité à intervenir et à proposer des solutions. Très efficace car il crée l'émotion et le débat.
\item \textbf{Les réseaux sociaux} (WhatsApp, Facebook, TikTok) : idéaux pour mobiliser les jeunes, diffuser des vidéos de démonstration (fabrication d'un pavé, d'un papier recyclé) et organiser des défis (« challenge zéro plastique »).
\item \textbf{Le porte-à-porte} : le plus lent mais le plus convaincant ; il permet de répondre aux questions, de montrer concrètement comment trier et d'indiquer le point de collecte le plus proche.
\end{itemize}

\courssub{Les messages clés et la règle des 3R}

Une sensibilisation efficace repose sur des messages \emph{simples, répétés et actionnables}. Les quatre messages clés à faire passer sont :

\begin{enumerate}
\item \textbf{Trier ses déchets à la source} : séparer, dès la maison, les déchets organiques (compostables), les papiers et cartons, les plastiques et les métaux.
\item \textbf{Réutiliser} : une bouteille plastique peut servir d'arrosoir, de pot de fleurs, de mangeoire pour poules ; un vieux cahier peut servir de brouillon.
\item \textbf{Recycler} : apporter les plastiques et papiers triés au point de collecte ou à l'atelier de transformation (bouteilles recyclées, pavés, papier hygiénique, papier-cadeau).
\item \textbf{Ne jamais brûler les plastiques à l'air libre} : les fumées dégagées sont toxiques pour toute la famille et tout le voisinage.
\end{enumerate}

Ces messages se résument dans un outil mémotechnique universel, très utile à connaître pour le Bac :

\begin{definition}[La règle des 3R]
La \cle{règle des 3R} est le principe de gestion durable des déchets fondé sur trois actions hiérarchisées :
\begin{itemize}
\item \textbf{Réduire} : diminuer la quantité de déchets produits à la source (acheter moins d'emballages, refuser les sachets inutiles) ;
\item \textbf{Réutiliser} : employer un objet plusieurs fois pour le même usage ou un nouvel usage avant de le jeter ;
\item \textbf{Recycler} : transformer un déchet en nouvelle matière première ou en nouveau produit (papier recyclé, bouteilles recyclées, pavés).
\end{itemize}
L'ordre compte : il vaut toujours mieux réduire que réutiliser, et réutiliser que recycler, car chaque étape en amont économise l'énergie et les matières premières des étapes suivantes.
\end{definition}

\begin{popfigure}
\begin{center}
\begin{tikzpicture}[scale=1.0]
% Trois noeuds en triangle
\node[circle, draw=popgreen, fill=popgreenL, very thick, minimum size=2.6cm, align=center, font=\bfseries] (R1) at (90:3) {R\'eduire\\ \footnotesize moins de\\ \footnotesize d\'echets};
\node[circle, draw=poporange, fill=poporangeL, very thick, minimum size=2.6cm, align=center, font=\bfseries] (R2) at (210:3) {R\'eutiliser\\ \footnotesize un objet\\ \footnotesize plusieurs fois};
\node[circle, draw=popblue, fill=popblueL, very thick, minimum size=2.6cm, align=center, font=\bfseries] (R3) at (330:3) {Recycler\\ \footnotesize transformer\\ \footnotesize en nouveau produit};
% Flèches de la boucle
\draw[-{Stealth[length=4mm]}, line width=2pt, popdark] (R1.south west) to[bend right=25] (R2.north east);
\draw[-{Stealth[length=4mm]}, line width=2pt, popdark] (R2.east) to[bend right=25] (R3.west);
\draw[-{Stealth[length=4mm]}, line width=2pt, popdark] (R3.north west) to[bend right=25] (R1.south east);
% Centre
\node[align=center, font=\itshape\small, text=popmuted] at (0,0) {Gestion durable\\ des d\'echets};
\end{tikzpicture}
\end{center}
\legende{Le cycle des 3R. Les trois actions forment une boucle vertueuse : réduire la production de déchets, réutiliser les objets, puis recycler ce qui ne peut plus être réutilisé. La flèche de « Recycler » vers « Réduire » rappelle que la matière recyclée remplace les matières premières neuves, ce qui réduit l'extraction des ressources.}
\end{popfigure}

\begin{savaistu}[Le tri à la source, l'étape la plus déterminante]
Savais-tu que le \cle{tri à la source} (c'est-à-dire le tri effectué par chacun, à la maison, au moment où l'on jette) est l'étape la plus déterminante pour la réussite de tout recyclage ? Un plastique souillé de restes alimentaires ou mélangé à des déchets organiques devient très coûteux, voire impossible à recycler : il faut le laver, le trier à nouveau, et sa qualité chute. En revanche, un déchet propre et bien trié dès le départ peut être valorisé directement, à moindre coût et avec un excellent rendement. C'est pourquoi toutes les campagnes de sensibilisation du monde insistent sur le même geste : \emph{trier chez soi}. Le succès d'une filière de recyclage se joue donc dans ta cuisine, pas seulement dans l'usine !
\end{savaistu}

\courssub{Construire un message de sensibilisation efficace}

Savoir sensibiliser est une compétence qui s'apprend. Voici la démarche complète, exigible comme savoir-faire.

\begin{methode}[Construire un message de sensibilisation en 4 étapes]
\begin{enumerate}
\item \textbf{Identifier la cible} : à qui parles-tu ? (ménages, élèves, commerçants du marché, autorités). Le langage, le support et l'exemple doivent être adaptés à cette cible.
\item \textbf{Nommer le problème} : formuler clairement le danger constaté (ex. : « les sachets brûlés dégagent des fumées qui provoquent la toux et l'asthme chez nos enfants »).
\item \textbf{Proposer une solution concrète} : indiquer un geste précis et faisable (ex. : « dépose tes plastiques dans le bac vert devant le marché, tous les samedis »).
\item \textbf{Lancer un appel à l'action} : terminer par une invitation claire, courte et motivante (ex. : « Ensemble, un quartier propre dès ce week-end ! »).
\end{enumerate}
Un bon message = une cible + un problème + une solution + un appel à l'action.
\end{methode}

\begin{attention}[Erreur fréquente : sensibiliser sans proposer de solution concrète]
L'erreur la plus fréquente — et la plus sanctionnée dans une copie — est de rédiger un message qui se contente de dénoncer : « Ne jetez pas vos déchets n'importe où ! » ou « La pollution est mauvaise ! ». Une sensibilisation \emph{sans solution concrète proposée} (point de collecte précis, atelier de recyclage, bac de tri, jour de ramassage) reste inefficace : la population, même convaincue, ne sait pas quoi faire de ses déchets et retombe dans ses habitudes. Au Bac, ton message de sensibilisation doit toujours indiquer \textbf{où}, \textbf{quand} et \textbf{comment} agir.
\end{attention}

\begin{popfigure}
\begin{center}
\begin{tikzpicture}
% Cadre de l'affiche
\draw[line width=2.5pt, popgreen, fill=popcream, rounded corners=8pt] (0,0) rectangle (12,9);
% Bandeau titre
\fill[popgreen, rounded corners=6pt] (0.4,7.2) rectangle (11.6,8.6);
\node[text=white, font=\bfseries\large, align=center] at (6,7.9) {MON QUARTIER PROPRE, MA SANT\'E D'ABORD !};
% Pictogramme 1 : poubelle de tri
\draw[very thick, popblue, fill=popblueL] (1.2,3.4) -- (1.5,5.6) -- (3.1,5.6) -- (3.4,3.4) -- cycle;
\draw[very thick, popblue] (1.0,5.6) -- (3.6,5.6);
\draw[very thick, popblue] (1.9,5.6) -- (1.9,5.95) -- (2.7,5.95) -- (2.7,5.6);
\node[font=\footnotesize\bfseries, text=popblue] at (2.3,4.4) {TRI};
\node[font=\scriptsize, align=center, text=popdark] at (2.3,2.9) {Je trie mes d\'echets\\ \`a la maison};
% Pictogramme 2 : flèche recyclage (triangle de flèches)
\draw[-{Stealth[length=3mm]}, line width=2.2pt, popgreen] (5.2,4.0) -- (6.6,5.4);
\draw[-{Stealth[length=3mm]}, line width=2.2pt, popgreen] (6.9,5.1) -- (6.9,3.3);
\draw[-{Stealth[length=3mm]}, line width=2.2pt, popgreen] (6.4,3.2) -- (5.0,3.9);
\node[font=\footnotesize\bfseries, text=popgreen] at (6.05,4.35) {3R};
\node[font=\scriptsize, align=center, text=popdark] at (6.05,2.9) {Je r\'eduis, je r\'eutilise,\\ je recycle};
% Pictogramme 3 : interdiction de brûler
\draw[very thick, poppink] (9.4,4.4) circle (1.15);
\draw[line width=2.5pt, poppink] (8.6,5.2) -- (10.2,3.6);
\draw[thick, poporange, fill=poporangeL] (9.4,3.9) .. controls (9.0,4.3) and (9.3,4.5) .. (9.4,4.9) .. controls (9.5,4.5) and (9.8,4.3) .. (9.4,3.9);
\node[font=\scriptsize, align=center, text=popdark] at (9.4,2.9) {Je ne br\^ule jamais\\ les plastiques};
% Bandeau appel à l'action
\fill[poporange, rounded corners=6pt] (0.4,0.4) rectangle (11.6,2.1);
\node[text=white, font=\bfseries, align=center] at (6,1.55) {D\'epose tes plastiques et papiers au point de collecte};
\node[text=white, font=\bfseries\large, align=center] at (6,0.85) {Place du March\'e --- tous les samedis de 8 h \`a 12 h};
\end{tikzpicture}
\end{center}
\legende{Maquette d'affiche de sensibilisation. Elle respecte la méthode en 4 étapes : cible (habitants du quartier), problème rappelé par les pictogrammes (déchets mal gérés, brûlage des plastiques), solution concrète (tri, 3R, point de collecte) et appel à l'action daté et localisé (bandeau orange). Les pictogrammes permettent de toucher aussi les personnes ne lisant pas le français.}
\end{popfigure}

\courssub{Le rôle de l'élève : ambassadeur du recyclage}

Toi, élève de Terminale C ou TI, tu possèdes quelque chose que beaucoup de membres de ta communauté n'ont pas : une formation scientifique qui te permet de \emph{comprendre} et d'\emph{expliquer} pourquoi et comment recycler. Cela fait de toi un \cle{ambassadeur du recyclage} :

\begin{itemize}
\item \textbf{Dans ta famille} : montre l'exemple en triant les déchets de la maison, explique à tes parents et à tes cadets pourquoi il ne faut pas brûler les sachets, propose de réutiliser les bouteilles et les boîtes.
\item \textbf{Dans ton école} : crée ou rejoins un club environnement, organise une journée « zéro déchet », un concours de la meilleure affiche ou un atelier de fabrication de papier recyclé.
\item \textbf{Dans ton quartier} : participe aux campagnes de porte-à-porte et aux journées de nettoyage, informe les voisins de l'existence des points de collecte, dialogue avec le chef de quartier pour l'installation de bacs de tri.
\end{itemize}

Le savoir-être attendu ici est la \cle{responsabilité citoyenne} : la biotechnologie n'atteint son but — un environnement sain — que si chacun devient acteur du changement.

\begin{exoresolu}[Rédiger un message de sensibilisation]
\textbf{Énoncé.} Ton lycée se situe près d'un marché où les commerçants brûlent chaque soir les sachets plastiques invendus ou usagés. Le club environnement te charge de rédiger un court message de sensibilisation (5 à 8 lignes) destiné aux commerçants, qui sera diffusé en spot radio dans la langue locale. Rédige ce message en respectant la méthode en 4 étapes.

\tcblower
\textbf{Solution détaillée.}

\emph{Étape 1 — Cible :} les commerçants du marché (adultes, langage simple et respectueux, radio en langue locale).

\emph{Étape 2 — Problème :} la fumée des plastiques brûlés provoque toux, asthme et maladies respiratoires chez les commerçants eux-mêmes, leurs enfants et leurs clients.

\emph{Étape 3 — Solution concrète :} au lieu de brûler, rassembler les sachets propres dans un sac et les déposer au bac de collecte installé à l'entrée du marché ; ils seront transformés en pavés et en bouteilles recyclées.

\emph{Étape 4 — Appel à l'action :} invitation datée et motivante.

\textbf{Message rédigé :}

« Chers commerçants, chaque soir, la fumée des sachets brûlés entre dans nos poumons et rend malades nos enfants. Ne brûlons plus nos plastiques ! Rassemblons-les dans un sac et déposons-les dans le grand bac vert à l'entrée du marché : ils deviendront des pavés pour nos rues. Dès ce samedi, apporte ton sac de plastiques au bac vert. Un marché propre, c'est des clients en bonne santé ! »

\textbf{Vérification :} cible précise, problème nommé, solution concrète (bac vert, entrée du marché), appel à l'action daté (« dès ce samedi »). Le message évite le piège de la dénonciation sans solution.
\end{exoresolu}

\begin{exercice}[À toi de jouer]
Ta commune organise une journée « Quartier propre ». Rédige le texte d'une affiche destinée aux ménages de ton quartier, comportant : un slogan, les trois pictogrammes à dessiner (décrits en une phrase chacun), les quatre messages clés de la leçon, et un appel à l'action précisant le lieu et la date du point de collecte. Indique ensuite, en trois lignes, pourquoi ton affiche respecte la méthode de construction d'un message efficace.
\end{exercice}

\begin{corrige}[Exercice]
Éléments attendus : (1) un slogan court et positif (ex. : « Trie aujourd'hui, respire demain ! ») ; (2) trois pictogrammes : une poubelle de tri (je trie à la maison), le symbole des 3R en boucle (je réduis, je réutilise, je recycle), une flamme barrée (je ne brûle pas les plastiques) ; (3) les quatre messages clés : trier à la source, réutiliser, recycler, ne pas brûler les plastiques à l'air libre ; (4) un appel à l'action avec lieu et date du point de collecte. Justification : l'affiche identifie sa cible (les ménages), rappelle le problème par les pictogrammes, propose une solution concrète (point de collecte localisé et daté) et se termine par un appel à l'action : elle suit donc les 4 étapes de la méthode.
\end{corrige}

\begin{aretenir}[L'essentiel de la section]
\begin{itemize}
\item Les déchets non valorisés polluent les sols, les eaux et l'air, et propagent des maladies (paludisme, choléra, infections respiratoires) : informer la communauté est une mesure de santé publique.
\item Les cibles prioritaires sont les ménages, les écoles, les marchés et les autorités locales ; les outils sont les affiches, les spots radio en langues locales, le théâtre forum, les réseaux sociaux et le porte-à-porte.
\item Les messages clés : trier à la source, réutiliser, recycler, ne jamais brûler les plastiques à l'air libre ; ils se résument par la règle des \cle{3R} : Réduire, Réutiliser, Recycler.
\item Un message efficace suit 4 étapes : cible, problème, solution concrète, appel à l'action ; sans solution concrète, la sensibilisation reste inefficace.
\item L'élève est un ambassadeur du recyclage dans sa famille, son école et son quartier.
\end{itemize}
\end{aretenir}

\newpage

\coursec{Activité d'intégration}

\coursec{Énergies renouvelables et valorisation des déchets}

\courssub{Objectifs de l'activité}
\par
Cette activité d'intégration vise à mobiliser l'ensemble des savoirs et savoir-faire de la leçon sur les énergies renouvelables et la valorisation des déchets. Elle te permettra de :
\begin{itemize}
    \item Analyser des données réelles de gestion des déchets dans un contexte urbain camerounais (Douala).
    \item Interpréter des résultats expérimentaux (production de biogaz, fermentation alcoolique) pour dimensionner des filières de valorisation.
    \item Comparer les avantages techniques, économiques et écologiques de différentes techniques de recyclage (papier, plastique).
    \item Concevoir une proposition argumentée de filière de valorisation locale adaptée à la composition des déchets.
    \item Communiquer sous forme d'une affiche de sensibilisation à destination de la communauté.
\end{itemize}

\courssub{Situation : Projet « Maképé Zéro Déchet »}
\par
Le quartier \textbf{Maképé}, situé dans le 3\ieme\ arrondissement de Douala, compte environ \num{10 000} habitants. La gestion des déchets y est un défi majeur : les ordures sont souvent déversées dans les caniveaux ou brûlées à ciel ouvert, provoquant des inondations récurrentes et une pollution de l'air. La mairie, en partenariat avec une ONG locale, lance le projet « Maképé Zéro Déchet » et mandate ton groupe d'experts juniors (élèves de Terminale C/TI) pour proposer une filière de valorisation locale.
\par
Voici le dossier documentaire mis à ta disposition.

\medskip
\noindent \textbf{Document 1 : Caractérisation des déchets ménagers du quartier Maképé (moyenne annuelle)}
\par
Une campagne de tri manuel sur \num{500} kg d'échantillons représentatifs a donné la composition massique suivante :

\begin{center}
\begin{tabularx}{\linewidth}{|l|X|c|}\hline
\rowcolor{popblueL}
\textbf{Catégorie} & \textbf{Description principale} & \textbf{Pourcentage massique (\%)} \\ \hline
Papiers / Cartons & Journaux, emballages, cahiers usagés & 20 \\ \hline
Plastiques & Bouteilles PET, sachets PE, bouchons, films & 25 \\ \hline
Matières organiques & Épluchures, restes de repas, déchets de marché, herbes & 45 \\ \hline
Autres & Verre, métaux, textiles, inertes, couches, déchets dangereux & 10 \\ \hline
\end{tabularx}
\end{center}
\par
La production moyenne de déchets est estimée à \SI{0,8}{kg/habitant/jour}.

\medskip
\noindent \textbf{Document 2 : Cinétique de production de biogaz par méthanisation des déchets organiques}
\par
Un digesteur pilote de \SI{10}{m^3} (type dôme fixe, température mésophile \SI{35}{\celsius}) a été alimenté avec \SI{200}{kg} de matières organiques broyées (taux de matière sèche volatile \textit{MSV} = \SI{80}{\%} de la matière fraîche). Le volume cumulé de biogaz produit (corrigé STP : \SI{0}{\celsius}, \SI{1}{atm}) a été mesuré quotidiennement pendant \num{30} jours.

\begin{popfigure}
\begin{tikzpicture}
\begin{axis}[
    width=\linewidth, height=8cm,
    xlabel={Temps (jours)}, ylabel={Volume cumulé de biogaz (m$^3$)},
    domain=0:30, samples=200,
    xmin=0, xmax=31, ymin=0, ymax=65,
    xtick={0,5,10,15,20,25,30}, ytick={0,10,20,30,40,50,60},
    grid=major, grid style={popline}, legend pos=north west,
    axis line style={popdark}, tick style={popdark},
    label style={font=\small}, tick label style={font=\small}
]
% Sigmoid curve: Vmax = 60 m3, k = 0.3, t0 = 10
\addplot[popcurve, thick] {60/(1+exp(-0.3*(x-10)))};
\addlegendentry{Production cumulée (modèle logistique)}
% Points expérimentaux simulés
\addplot[only marks, mark=*, mark size=2, poporange] coordinates {
(0,0) (2,2.5) (5,12) (8,30) (10,38) (12,45) (15,52) (18,56) (22,58.5) (26,59.5) (30,60)
};
\addlegendentry{Données expérimentales}
\end{axis}
\end{tikzpicture}
\legende{Évolution de la production cumulée de biogaz (en m$^3$ aux conditions STP) en fonction du temps de rétention hydraulique dans un digesteur pilote de 10 m$^3$ alimenté en déchets organiques de cuisine et de marché.}
\end{popfigure}

\medskip
\noindent \textbf{Document 3 : Expérience de fermentation du jus de canne à sucre par \textit{Saccharomyces cerevisiae} (production de bioéthanol)}
\par
\emph{Protocole :} \SI{1,0}{L} de jus de canne stérilisé (titre sucré \SI{18}{\%} massique en saccharose, densité \SI{1,07}{g/mL}) est inoculé avec \SI{10}{g} de levure sèche active. Le fermenteur (volume \SI{1,5}{L}, muni d'un piège à eau) est maintenu à \SI{30}{\celsius} sous agitation lente. Le volume de \ce{CO2} dégagé est mesuré par déplacement d'eau toutes les \num{12} heures pendant \num{5} jours.
\par
\emph{Résultats :}

\begin{center}
\begin{tabularx}{\linewidth}{|c|X|X|X|}\hline
\rowcolor{popgreenL}
\textbf{Temps (h)} & \textbf{Volume \ce{CO2} cumulé (mL)} & \textbf{Masse \ce{CO2} (g)} & \textbf{Observations} \\ \hline
0 & 0 & 0,00 & Inoculation, milieu trouble \\ \hline
12 & 120 & 0,21 & Mousse abondante, dégagement gazeux vif \\ \hline
24 & 380 & 0,67 & Activité maximale \\ \hline
36 & 580 & 1,02 & Ralentissement visible \\ \hline
48 & 720 & 1,27 & Quasi-arrêt bulles \\ \hline
60 & 760 & 1,34 & Fin de fermentation, dépôt levure \\ \hline
72 & 760 & 1,34 & Stable \\ \hline
\end{tabularx}
\end{center}
\par
Rappel : Équation bilan de la fermentation alcoolique : \ce{C12H22O11 + H2O -> 4 C2H5OH + 4 CO2}. Masse molaire : \ce{C12H22O11} = \SI{342}{g/mol} ; \ce{CO2} = \SI{44}{g/mol} ; \ce{C2H5OH} = \SI{46}{g/mol}. Densité éthanol = \SI{0,79}{g/mL}.

\medskip
\noindent \textbf{Document 4 : Comparatif technique et économique : Pavé ciment vs Pavé plastique-sable}
\par
Une startup locale « \textit{EcoPave Douala} » propose des pavés fabriqués à partir de déchets plastiques (PE, PP, PS) fondus et mélangés à du sable de rivière (ratio masse plastique : sable = 1 : 3). Les caractéristiques comparées sont :

\begin{center}
\begin{tabularx}{\linewidth}{|l|X|X|}\hline
\rowcolor{poppurpleL}
\textbf{Paramètre} & \textbf{Pavé ciment classique (vibré)} & \textbf{Pavé plastique-sable (EcoPave)} \\ \hline
Composition & Ciment Portland, sable, gravier, eau & Plastiques fondus (liant), sable (charge) \\ \hline
Résistance à la compression (28 j) & \SI{25}{MPa} (Classe B25) & \SI{30}{MPa} \\ \hline
Absorption d'eau & \SI{6}{\%} & \SI{0,5}{\%} (quasi imperméable) \\ \hline
Poids unitaire (20x10x6 cm) & \SI{2,2}{kg} & \SI{1,8}{kg} \\ \hline
Coût de production unitaire & \SI{500}{FCFA} & \SI{350}{FCFA} \\ \hline
Prix de vente unitaire (marché) & \SI{650}{FCFA} & \SI{500}{FCFA} \\ \hline
Durée de vie estimée & \num{20} ans & \num{25} ans (résistance UV additivée) \\ \hline
Impact environnemental & Émission \ce{CO2} cimentière (~\SI{0,8}{kg CO2/kg ciment}) & Valorisation \SI{0,45}{kg} plastique/pavé, pas de cuisson \\ \hline
\end{tabularx}
\end{center}

\courssub{Consignes}
\par
Travaillant en groupe de \num{4} à \num{5} élèves, tu devras rendre un rapport écrit (format A4, \num{4} pages max) et préparer une affiche A3 pour la restitution orale (\num{10} min + \num{5} min questions). Les questions ci-dessous guident ton travail :

\begin{enumerate}
    \item \textbf{Quantification du gisement :} À partir du Document 1, calcule la masse annuelle (en tonnes) de chaque catégorie de déchets produite par le quartier Maképé. Quelle catégorie représente le gisement le plus important ?
    \item \textbf{Analyse de la filière Biogaz (Doc 2) :}
    \begin{enumerate}
        \item Identifie les trois phases classiques de la méthanisation sur la courbe et associe-les aux groupes microbiens dominants (hydrolyseurs, acétogènes, méthanogènes).
        \item Détermine le rendement massique de biogaz (en \SI{}{m^3/kg MSV}) atteint au bout de \num{30} jours.
        \item Estime la production annuelle potentielle de biogaz (en \SI{}{m^3/an}) si \num{80}{\%} des déchets organiques du quartier sont collectés et méthanisés. Sachant que le biogaz contient \num{60}{\%} de \ce{CH4} (PCI \ce{CH4} = \SI{9,97}{kWh/m^3}), calcule l'énergie électrique récupérable annuellement (rendement cogénération \num{35}{\%}).
    \end{enumerate}
    \item \textbf{Analyse de la filière Bioéthanol (Doc 3) :}
    \begin{enumerate}
        \item Vérifie par le calcul stœchiométrique si la masse de \ce{CO2} mesurée à \num{72} h correspond à la consommation totale du saccharose initial.
        \item Calcule le volume d'éthanol pur (en L) théoriquement produit et le titre alcoolique volumique (\% vol) du vin de canne obtenu.
        \item Pourquoi cette filière est-elle moins adaptée aux déchets du quartier que la méthanisation ? (Indice : nature du substrat, prétraitement).
    \end{enumerate}
    \item \textbf{Analyse de la filière Pavés plastiques (Doc 4) :}
    \begin{enumerate}
        \item Calcule la masse de déchets plastiques nécessaire pour produire \num{10 000} pavés EcoPave.
        \item Compare le coût total d'un trottoir de \SI{100}{m^2} (pose comprise estimée à \SI{1500}{FCFA/m^2} pour les deux types) entre pavés ciment et pavés plastique-sable.
        \item Liste trois avantages techniques et deux avantages environnementaux majeurs du pavé plastique-sable.
    \end{enumerate}
    \item \textbf{Synthèse et Choix :} Propose \textbf{deux} filières de valorisation prioritaires pour Maképé. Justifie ton choix par une analyse croisée : gisement disponible, maturité technique, rentabilité économique, impact social (emplois) et environnemental. Décris les \cle{étapes techniques} principales de chaque filière choisie (schéma flux simplifié attendu).
    \item \textbf{Communication :} Conçois l'affiche de sensibilisation « Maképé Zéro Déchet : Mes déchets, ma richesse ». Elle doit comporter : un slogan accrocheur, le tri à la source (3 bacs minimum), les deux filières choisies avec 1 chiffre clé chacune, un appel à l'action pour les habitants (dépôt volontaire, horaire collecte).
\end{enumerate}

\courssub{Ressources à mobiliser}
\par
Pour résoudre cette activité, tu devras t'appuyer sur les notions suivantes (rappel synthétique) :
\begin{itemize}
    \item \textbf{Biocarburants solides/liquides/gazeux :} Différence entre bioéthanol (fermentation sucres/amidon), biodiesel (transestérification huiles) et biogaz (méthanisation anaérobie). Rendements typiques : biogaz \SIrange{0,3}{0,6}{m^3/kg MSV} ; bioéthanol \SIrange{0,4}{0,5}{L/kg} sucre.
    \item \textbf{Méthanisation :} 4 étapes (hydrolyse, acidogenèse, acétogenèse, méthanogenèse). Conditions : anaérobiose stricte, pH \num{6,8}-\num{7,2}, T mésophile (\SI{35}{\celsius}) ou thermophile (\SI{55}{\celsius}). Biogaz = \ce{CH4} (\num{50}-\num{70}\%) + \ce{CO2} + traces \ce{H2S}, \ce{H2O}.
    \item \textbf{Fermentation alcoolique :} \ce{C6H12O6 -> 2 C2H5OH + 2 CO2} (glucose) ou \ce{C12H22O11 + H2O -> 4 C2H5OH + 4 CO2} (saccharose). Rendement théorique : \SI{51,1}{\%} massique sucre $\to$ éthanol. Inhibition par l'éthanol > \num{12}-\num{15}\% vol.
    \item \textbf{Recyclage papier :} Tri $\to$ Broyage/Pulpage $\to$ Désencrage (flottation) $\to$ Épuration $\to$ Raffinage $\to$ Formage/Séchage $\to$ Produits (hygiène, emballage). Perte de fibres $\to$ nombre de cycles limité (\num{5}-\num{7}).
    \item \textbf{Recyclage plastique :} Tri (code resin) $\to$ Broyage $\to$ Lavage $\to$ Extrusion/Granulats $\to$ Soufflage/Injection (bouteilles) OU Extrusion/Compression (pavés, profilés). Thermoplastiques (PET, PE, PP, PS) recyclables ; Thermodurcissables non recyclables par fusion.
    \item \textbf{Formules utiles :}
    \begin{align*}
        \text{Masse annuelle} &= \text{Population} \times \text{Production journalière} \times 365 \times \% \text{catégorie} \\
        \text{Rendement biogaz} &= \frac{V_{\text{biogaz final}}}{m_{\text{MSV introduit}}} \\
        \text{Énergie électrique} &= V_{\ce{CH4}} \times \text{PCI}_{\ce{CH4}} \times \eta_{\text{cogénération}} \\
        n(\ce{CO2}) &= \frac{m(\ce{CO2})}{M(\ce{CO2})} \quad ; \quad n(\text{sucre}) = \frac{n(\ce{CO2})}{4} \quad (\text{saccharose})
    \end{align*}
\end{itemize}

\courssub{Démarche et corrigé commenté}
\begin{exoresolu}[Corrigé détaillé de l'activité d'intégration « Maképé Zéro Déchet »]
\par
\noindent \textbf{1. Quantification du gisement annuel (Doc 1)}
\par
Population $P = \num{10 000}$ hab. Production $p = \SI{0,8}{kg/hab/j}$.
Masse totale annuelle $M_{\text{tot}} = P \times p \times 365 = 10\,000 \times 0,8 \times 365 = \SI{2\,920\,000}{kg} = \mathbf{\SI{2\,920}{tonnes/an}}$.
\begin{itemize}
    \item Papiers : $20\% \times 2\,920 = \mathbf{\SI{584}{t/an}}$
    \item Plastiques : $25\% \times 2\,920 = \mathbf{\SI{730}{t/an}}$
    \item Organiques : $45\% \times 2\,920 = \mathbf{\SI{1\,314}{t/an}}$
    \item Autres : $10\% \times 2\,920 = \mathbf{\SI{292}{t/an}}$
\end{itemize}
\par
\emph{Conclusion :} La fraction \textbf{organique} est le gisement principal (\SI{1\,314}{t/an}), suivie des plastiques (\SI{730}{t/an}).

\medskip
\noindent \textbf{2. Filière Biogaz (Doc 2)}
\par
\textbf{a) Phases de la courbe :}
\begin{itemize}
    \item \textbf{Jours 0-5 (Latence/Lente) :} \textbf{Hydrolyse} \& \textbf{Acidogenèse}. Bactéries hydrolytiques (enzymes extracellulaires) et fermentaires transforment polymères (protéines, lipides, glucides) en acides gras volatils (AGV), \ce{CO2}, \ce{H2}, \ce{NH3}. pH baisse légèrement.
    \item \textbf{Jours 5-15 (Exponentielle) :} \textbf{Acétogenèse} \& \textbf{Méthanogenèse} acétoclastique (\textit{Methanosaeta}, \textit{Methanosarcina}). Conversion AGV (acétate) $\to$ \ce{CH4} + \ce{CO2}. Production maximale de gaz.
    \item \textbf{Jours 15-30 (Stationnaire) :} Épuisement substrat accessible. Méthanogenèse hydrogénotrophe (\ce{4H2 + CO2 -> CH4 + 2H2O}) résiduelle. Courbe plateau.
\end{itemize}
\par
\textbf{b) Rendement massique :}
Volume final $V_f = \SI{60}{m^3}$ (STP).
Masse substrat frais $m_f = \SI{200}{kg}$. MSV = $80\% \times 200 = \SI{160}{kg}$.
Rendement $R = \frac{60}{160} = \mathbf{\SI{0,375}{m^3/kg\,MSV}}$. (Valeur réaliste pour déchets de cuisine/marché, riches en lipides/protéines).
\par
\textbf{c) Potentiel quartier :}
Organiques collectés = $80\% \times \SI{1\,314}{t} = \SI{1\,051,2}{t/an} = \SI{1\,051\,200}{kg/an}$.
MSV disponible (hypothèse même taux $80\%$) = $\SI{840\,960}{kg\,MSV/an}$.
Biogaz potentiel $V_{\text{bio}} = 840\,960 \times 0,375 = \mathbf{\SI{315\,360}{m^3/an}}$.
Volume \ce{CH4} ($60\%$) = $0,60 \times 315\,360 = \SI{189\,216}{m^3/an}$.
Énergie thermique \ce{CH4} = $189\,216 \times \SI{9,97}{kWh/m^3} = \SI{1\,886\,483}{kWh/an}$.
Énergie électrique (cogénération $\eta = 35\%$) = $1\,886\,483 \times 0,35 = \mathbf{\SI{660\,269}{kWh/an} \approx \SI{660}{MWh/an}}$.
\emph{Interprétation :} Cela couvre les besoins électriques d'environ \textbf{3 300 foyers camerounais} (conso moy. $\sim \SI{200}{kWh/an}$/foyer) ou alimente l'éclairage public du quartier.

\medskip
\noindent \textbf{3. Filière Bioéthanol (Doc 3)}
\par
\textbf{a) Vérification stœchiométrique :}
Masse \ce{CO2} finale $m_{\ce{CO2}} = \SI{1,34}{g}$.
$n_{\ce{CO2}} = \frac{1,34}{44} = \SI{0,03045}{mol}$.
Équation : 1 mol saccharose $\to$ 4 mol \ce{CO2}.
$n_{\text{saccharose consommé}} = \frac{0,03045}{4} = \SI{0,00761}{mol}$.
Masse saccharose consommé = $0,00761 \times 342 = \SI{2,60}{g}$.
Masse saccharose initial = $18\% \times \SI{1070}{g}$ (masse \SI{1}{L} jus) = \SI{192,6}{g}.
\emph{Seulement \SI{1,35}{\%} du sucre a été consommé.} La fermentation est \textbf{incomplète} (arrêt précoce : inhibition éthanol ? manque nutriments azotés ? température ?). Le volume \ce{CO2} mesuré ne correspond PAS à la consommation totale du substrat.
\par
\textbf{b) Production éthanol théorique (si fermentation complète) :}
$n_{\text{saccharose total}} = \frac{192,6}{342} = \SI{0,563}{mol}$.
$n_{\ce{C2H5OH}} \text{ max} = 4 \times 0,563 = \SI{2,252}{mol}$.
$m_{\ce{C2H5OH}} \text{ max} = 2,252 \times 46 = \SI{103,6}{g}$.
$V_{\ce{C2H5OH}} \text{ pur} = \frac{103,6}{0,79} = \mathbf{\SI{131,1}{mL}}$.
Volume total vin $\approx \SI{1000}{mL}$ (variation volume négligeable).
Titre alcoolique volumique (TAV) = $\frac{131,1}{1000} \times 100 = \mathbf{\SI{13,1}{\%\,vol}}$.
\par
\textbf{c) Inadaptation au quartier :}
Le bioéthanol nécessite un substrat \textbf{sucrier pur} (canne, betterave, mélasse) ou \textbf{amylacé} (maïs, manioc) avec hydrolyse enzymatique coûteuse. Les déchets organiques de Maképé sont \textbf{hétérogènes, lignocellulosiques, pauvres en sucres simples}, riches en eau et azote. La méthanisation est la seule voie biologique robuste pour ce gisement « déchets humides mélangés ».

\medskip
\noindent \textbf{4. Filière Pavés Plastique-Sable (Doc 4)}
\par
\textbf{a) Masse plastique pour 10 000 pavés :}
Masse plastique/pavé = \SI{0,45}{kg} (donné tableau).
Masse totale = $10\,000 \times 0,45 = \mathbf{\SI{4\,500}{kg} = \SI{4,5}{tonnes}}$.
\par
\textbf{b) Coût trottoir \SI{100}{m^2} :}
Surface pavé = $0,20 \times 0,10 = \SI{0,02}{m^2}$. Nombre pavés/$m^2$ = 50.
Total pavés = $100 \times 50 = \num{5\,000}$ pavés.
\begin{itemize}
    \item \textbf{Ciment :} Achat $5\,000 \times 650 = \SI{3\,250\,000}{FCFA}$. Pose $100 \times 1\,500 = \SI{150\,000}{FCFA}$. \textbf{Total = \SI{3\,400\,000}{FCFA}}.
    \item \textbf{Plastique-sable :} Achat $5\,000 \times 500 = \SI{2\,500\,000}{FCFA}$. Pose $100 \times 1\,500 = \SI{150\,000}{FCFA}$. \textbf{Total = \SI{2\,650\,000}{FCFA}}.
\end{itemize}
\textbf{Économie : \SI{750\,000}{FCFA} (\num{22}\% moins cher).}
\par
\textbf{c) Avantages pavé plastique-sable :}
\begin{itemize}
    \item \textbf{Techniques :} 1) Résistance compression supérieure (\SI{30}{MPa} vs \SI{25}{MPa}). 2) Imperméabilité quasi totale (\SI{0,5}{\%} absorption vs \SI{6}{\%}) $\to$ pas de gel/dégel, durabilité accrue. 3) Poids léger (\SI{1,8}{kg} vs \SI{2,2}{kg}) $\to$ transport/pose facilités.
    \item \textbf{Environnementaux :} 1) Valorisation directe \SI{4,5}{t} plastique pour \num{10 000} pavés (évite incinération/océan). 2) Zéro émission \ce{CO2} liée au clinker (ciment) ; pas de cuisson four (\SI{1450}{\celsius}).
\end{itemize}

\medskip
\noindent \textbf{5. Synthèse : Choix des deux filières prioritaires pour Maképé}
\par
\emph{Filière 1 : Méthanisation centralisée des déchets organiques (Biogaz $\to$ Électricité/Chaleur).}
\emph{Justification :} Gisement massive (\SI{1\,314}{t/an}), technologie maîtrisée (digesteurs dôme fixe type chinois/indien adaptés), produit énergie locale (éclairage public, cuisson cantine scolaire, charge téléphones), digestat = engrais organique pour agriculture urbaine. Crée emplois : collecte triée (\num{10} jeunes), opérateur digesteur (\num{2}), maintenance.
\par
\emph{Étapes techniques (Schéma flux) :}
\begin{center}
\begin{tikzpicture}[node distance=1.2cm and 1.5cm, >=Stealth, font=\small]
\node[draw, rectangle, rounded corners, fill=popgreenL, minimum width=2.5cm, align=center] (tri){Tri à la source\\(Bac vert)};
\node[draw, rectangle, rounded corners, fill=poporangeL, right=of tri, align=center] (collecte){Collecte sélective\\porte-à-porte};
\node[draw, rectangle, rounded corners, fill=popblueL, right=of collecte, align=center] (pretrait){Prétraitement :\\Broyage + Déshydratation\\+ Mélange eau};
\node[draw, rectangle, rounded corners, fill=poppurpleL, right=of pretrait, align=center] (dig){Digesteur anaérobie\\(Dôme fixe 50-100 m$^3$)};
\node[draw, rectangle, rounded corners, fill=poppinkL, below=of dig] (biogaz) {Biogaz (\ce{CH4}/\ce{CO2})};
\node[draw, rectangle, rounded corners, fill=popgoldL, right=of dig, align=center] (cogen){Cogénération\\Moteur Gaz/Générateur};
\node[draw, rectangle, rounded corners, fill=popgreenL, below=of pretrait, align=center] (digestat){Digestat liquide\\$\to$ Compostage $\to$ Engrais};
\draw[->, thick, popdark] (tri) -- (collecte);
\draw[->, thick, popdark] (collecte) -- (pretrait);
\draw[->, thick, popdark] (pretrait) -- (dig);
\draw[->, thick, poporange] (dig) -- node[above] {Biogaz} (cogen);
\draw[->, thick, popgreen] (dig) -- node[left] {Digestat} (digestat);
\draw[->, thick, popblue] (cogen) -- ++(1,0) node[right] {Électricité / Chaleur};
\end{tikzpicture}
\end{center}
\par
\emph{Filière 2 : Unité de fabrication de pavés plastique-sable (Plastiques $\to$ Pavés).}
\emph{Justification :} Gisement plastique significatif (\SI{730}{t/an}), technologie low-tech (extrudeuse/compresseur chauffant, moules métalliques), forte demande locale (voirie, cours écoles, parkings), rentabilité immédiate (marge \SI{150}{FCFA}/pavé), nettoie visuellement le quartier. Crée emplois : trieurs/lavage (\num{15} femmes), opérateurs machine (\num{3}), commercialisation.
\par
\emph{Étapes techniques :}
Tri sélectif (Bac jaune) $\to$ Collecte $\to$ Tri fin (séparation PET/PE/PP) $\to$ Broyage (flocons \SI{<10}{mm}) $\to$ Lavage/Séchage $\to$ \textbf{Extrusion/Fusion} (T \SIrange{180}{220}{\celsius}) + \textbf{Injection Sable} (ratio 1:3) $\to$ \textbf{Moulage/Compression} (moules chauffés) $\to$ Démoulage/Refroidissement $\to$ Contrôle qualité $\to$ Stockage/Vente.

\medskip
\noindent \textbf{6. Affiche de sensibilisation (Proposition de contenu)}
\par
\begin{center}
\fbox{\begin{minipage}{0.9\linewidth}
\textbf{\Large MAKÉPÉ ZÉRO DÉCHET : MES DÉCHETS, MA RICHESSE !}\\[0.5em]
\textbf{TRIEZ À LA SOURCE : 3 BACS, 3 GESTES SIMPLES}\\
\begin{tabular}{ccc}
\cellcolor{popgreenL}\textbf{BAC VERT} & \cellcolor{poporangeL}\textbf{BAC JAUNE} & \cellcolor{popgray}\textbf{BAC GRIS} \\
Organiques & Plastiques & Autres/Rejets \\
Épluchures, restes & Bouteilles, sachets, & Verre, couches, \\
de cuisine, herbes & bouchons, films & balayures\ldots
\end{tabular}\\[1em]
\textbf{VOS DÉCHETS DEVIENNENT RESSOURCES :}\\
\checkmark \textbf{1 314 tonnes/an d'organique $\to$ BIOGAZ $\to$ ÉCLAIRAGE PUBLIC + ENGRAIS} (Économie \SI{660}{MWh/an})\\
\checkmark \textbf{730 tonnes/an de plastique $\to$ PAVÉS ÉCOLOGIQUES} (Économie \SI{22}{\%} vs ciment, \SI{4,5}{t} plastique sauvées/10 000 pavés)\\[1em]
\textbf{AGISSEZ MAINTENANT !}\\
Dépôt volontaire : Carrefour Maképé, Lundi-Samedi 6h-18h.\\
Collecte porte-à-porte : Mardi \& Vendredi (Bac Vert), Mercredi \& Samedi (Bac Jaune).\\
\textit{Contact : Mairie 3\ieme\ Arrondissement / ONG EcoMaképé \#ZéroDéchetDouala}
\end{minipage}}
\end{center}
\end{exoresolu}

\courssub{Bilan de l'activité}
\par
Cette activité d'intégration a permis de mettre en évidence plusieurs points clés du programme :
\begin{itemize}
    \item \textbf{L'approche systémique :} La gestion des déchets ne se résout pas par une technique unique mais par une \cle{combinaison de filières} adaptées à la composition du gisement (organique $\to$ méthanisation ; plastique $\to$ recyclage matière ; papier $\to$ recyclage fibre).
    \item \textbf{La rigueur quantitative :} Le dimensionnement des installations (volume digesteur, nombre de pavés, production énergétique) repose sur des calculs de bilan de masse et d'énergie fondés sur des données expérimentales (rendement biogaz \SI{0,375}{m^3/kg MSV}, masse plastique/pavé \SI{0,45}{kg}).
    \item \textbf{L'analyse critique des procédés :} La comparaison Doc 3 vs Doc 2 illustre l'importance du \cle{choix du procédé en fonction du substrat} : la fermentation alcoolique exige des sucres fermentescibles purs, la méthanisation valorise la matière organique complexe humide.
    \item \textbf{L'évaluation multicritères :} Le choix du pavé plastique-sable ne repose pas seulement sur le prix (\SI{350}{FCFA} vs \SI{500}{FCFA}) mais sur ses performances techniques supérieures (imperméabilité, résistance) et son bilan carbone évité (pas de ciment).
    \item \textbf{La dimension citoyenne :} La réussite technique suppose un \cle{tri à la source} efficace par la population, d'où la nécessité d'une communication claire (affiche, code couleur bacs, calendrier collecte) pour transformer le comportement social.
\end{itemize}
\par
Tu as ainsi mobilisé les savoir-faire « Expliquer les techniques... », « Informer la communauté... » et « Analyser des documents... » dans une situation authentique ancrée dans la réalité camerounaise (Douala, déchets urbains, essaimage technologique local). C'est cette capacité à \textbf{transférer tes connaissances scolaires vers la résolution de problèmes concrets} que le Baccalauréat et ton future engagement citoyen évaluent.

\newpage

\coursec{Méthodes et exercices résolus}

\courssub{Méthode 1 : Expliquer une technique de production de biocarburant}

\begin{methode}[Rédiger l'explication d'une filière biocarburant]
Pour expliquer correctement une technique de production de biocarburant (bioéthanol, biodiesel, biogaz), adopte toujours la démarche en chaîne suivante :
\begin{enumerate}
\item \textbf{Identifier la matière première} (biomasse) : canne à sucre, maïs, manioc pour le bioéthanol ; palmier à huile, soja, coton pour le biodiesel ; déjections animales, déchets organiques pour le biogaz.
\item \textbf{Nommer le procédé biologique ou chimique} : fermentation alcoolique par des levures (\emph{Saccharomyces cerevisiae}) pour le bioéthanol ; transestérification (huile + alcool) pour le biodiesel ; digestion anaérobie (méthanisation) pour le biogaz.
\item \textbf{Écrire l'équation-bilan} si elle existe : \ce{C6H12O6 -> 2 C2H5OH + 2 CO2} pour la fermentation alcoolique.
\item \textbf{Décrire les étapes technologiques} dans l'ordre : broyage/extraction $\rightarrow$ hydrolyse éventuelle $\rightarrow$ fermentation ou réaction $\rightarrow$ distillation ou purification $\rightarrow$ conditionnement.
\item \textbf{Conclure par l'intérêt} : énergie renouvelable, bilan carbone quasi neutre (le \ce{CO2} dégagé a été capté par la plante), valorisation agricole locale.
\end{enumerate}
Réflexe : le mot \cle{fermentation} doit toujours être associé à un micro-organisme et à l'\cle{anaérobiose} (absence d'oxygène).
\end{methode}

\begin{exoresolu}[Production de bioéthanol à partir de la canne à sucre]
\textbf{Énoncé.} Une coopérative de la région de l'Extrême-Nord souhaite produire du bioéthanol à partir du jus de canne à sucre. (a) Décris les étapes de cette production. (b) Écris l'équation-bilan de la réaction clé et précise l'agent biologique responsable. (c) Calcule la masse d'éthanol théoriquement obtenue à partir de \SI{900}{kg} de glucose. On donne : $M(\ce{C6H12O6}) = \SI{180}{g.mol^{-1}}$ ; $M(\ce{C2H5OH}) = \SI{46}{g.mol^{-1}}$.
\tcblower
\textbf{Solution.}

(a) \textbf{Étapes de production.}
\begin{enumerate}
\item \textbf{Broyage et extraction} : les tiges de canne sont broyées et pressées pour obtenir un jus riche en saccharose.
\item \textbf{Hydrolyse} : le saccharose est hydrolysé en glucose et fructose (sucres fermentescibles).
\item \textbf{Fermentation alcoolique} : le jus est ensemencé avec la levure \emph{Saccharomyces cerevisiae} en cuve close, en anaérobiose, à environ \SI{30}{\celsius}.
\item \textbf{Distillation} : le moût fermenté (environ 8 à 12 \% d'éthanol) est distillé pour concentrer l'éthanol.
\item \textbf{Déshydratation et conditionnement} : on obtient l'éthanol carburant.
\end{enumerate}

(b) \textbf{Équation-bilan.} La fermentation alcoolique est réalisée par les levures en absence d'oxygène :
\[ \ce{C6H12O6 -> 2 C2H5OH + 2 CO2} \]
L'agent biologique est la levure \emph{Saccharomyces cerevisiae}.

(c) \textbf{Calcul de la masse d'éthanol.}
D'après l'équation, 1 mole de glucose donne 2 moles d'éthanol. Donc \SI{180}{g} de glucose donnent $2 \times 46 = \SI{92}{g}$ d'éthanol.
\begin{align*}
m(\text{éthanol}) &= \SI{900}{kg} \times \dfrac{92}{180} \\
&= \SI{900}{kg} \times \num{0,5111} \\
&= \SI{460}{kg}
\end{align*}
\textbf{Conclusion :} la fermentation de \SI{900}{kg} de glucose produit théoriquement \SI{460}{kg} d'éthanol (en pratique moins, car le rendement de fermentation n'est jamais de 100 \%).
\end{exoresolu}

\courssub{Méthode 2 : Expliquer la production d'énergie à partir du papier}

\begin{methode}[Expliquer la valorisation énergétique du papier]
\begin{enumerate}
\item \textbf{Préciser la nature du papier} : matériau cellulosique d'origine végétale, donc une biomasse combustible.
\item \textbf{Distinguer les deux voies} : l'\cle{incinération} avec récupération de chaleur (combustion directe dans un four : la chaleur produit de la vapeur qui peut actionner un alternateur) et la \cle{méthanisation} des déchets papiers mélangés aux biodéchets (production de biogaz, mélange de \ce{CH4} et \ce{CO2}, brûlé pour produire chaleur et électricité).
\item \textbf{Décrire la chaîne de conversion d'énergie} : énergie chimique (cellulose) $\rightarrow$ énergie thermique (combustion) $\rightarrow$ énergie mécanique (turbine) $\rightarrow$ énergie électrique (alternateur).
\item \textbf{Mentionner les conditions} : tri préalable (éliminer plastiques et métaux), séchage, broyage ; contrôle des fumées pour limiter la pollution.
\item \textbf{Conclure} : réduction du volume des déchets et production d'énergie renouvelable.
\end{enumerate}
\end{methode}

\begin{exoresolu}[Incinération des déchets papiers d'un lycée]
\textbf{Énoncé.} Un lycée de Douala produit \SI{500}{kg} de déchets papiers par mois. Ces déchets sont incinérés dans une chaufferie dont le pouvoir calorifique du papier sec est $P = \SI{15}{MJ.kg^{-1}}$. (a) Explique les étapes de la conversion de l'énergie contenue dans le papier en énergie électrique. (b) Calcule l'énergie thermique libérée par la combustion mensuelle de ces déchets. (c) Le rendement global de conversion en électricité est de 25 \% ; calcule l'énergie électrique produite en kilowattheures ($\SI{1}{kWh} = \SI{3,6}{MJ}$).
\tcblower
\textbf{Solution.}

(a) \textbf{Étapes de la conversion.}
\begin{enumerate}
\item Tri et séchage du papier pour éliminer les corps étrangers et l'humidité.
\item Combustion du papier dans le four : l'énergie chimique de la cellulose devient énergie thermique.
\item La chaleur transforme l'eau d'une chaudière en vapeur sous pression.
\item La vapeur fait tourner une turbine : énergie thermique $\rightarrow$ énergie mécanique.
\item La turbine entraîne un alternateur : énergie mécanique $\rightarrow$ énergie électrique.
\end{enumerate}

(b) \textbf{Énergie thermique libérée.}
\[ E_{\text{th}} = m \times P = \SI{500}{kg} \times \SI{15}{MJ.kg^{-1}} = \SI{7500}{MJ} \]

(c) \textbf{Énergie électrique produite.}
\begin{align*}
E_{\text{él}} &= 0{,}25 \times \SI{7500}{MJ} = \SI{1875}{MJ} \\
E_{\text{él}} &= \dfrac{1875}{\num{3,6}} \approx \SI{521}{kWh}
\end{align*}
\textbf{Conclusion :} l'incinération des \SI{500}{kg} de papier produit environ \SI{521}{kWh} d'électricité par mois, de quoi alimenter plusieurs salles de classe, tout en réduisant le volume des déchets du lycée.
\end{exoresolu}

\courssub{Méthode 3 : Expliquer le recyclage du papier (serviettes, papier hygiénique, papier-cadeau)}

\begin{methode}[Décrire la filière de recyclage du papier]
La trame de rédaction est identique quel que soit le produit fini ; seules la finition et le traitement diffèrent :
\begin{enumerate}
\item \textbf{Collecte et tri} : séparer le papier des autres déchets ; éliminer agrafes, plastiques, papiers souillés ou plastifiés.
\item \textbf{Mise en pâte} : le papier est découpé puis trempé dans l'eau avec des produits de désencrage ; le malaxage donne une \cle{pâte à papier} (suspension de fibres de cellulose).
\item \textbf{Épuration et blanchiment} : filtration, désencrage, blanchiment éventuel (indispensable pour le papier hygiénique, pour des raisons d'hygiène et d'aspect).
\item \textbf{Formage} : la pâte est étalée en couche mince sur une toile tamisante ; l'eau s'égoutte.
\item \textbf{Pressage et séchage} : les fibres feutrent en une feuille continue.
\item \textbf{Finition selon le produit} :
\begin{itemize}
\item serviettes jetables : découpe, pliage, gaufrage ;
\item papier hygiénique : feuille très fine, blanchie, enroulée en rouleaux ;
\item papier-cadeau : feuille plus épaisse, colorée ou imprimée.
\end{itemize}
\end{enumerate}
Réflexe Bac : citer toujours la \cle{cellulose} comme constituant recyclé, et préciser que les fibres raccourcissent à chaque recyclage (le papier ne se recycle qu'environ 5 à 7 fois).
\end{methode}

\begin{exoresolu}[Recyclage du papier en papier hygiénique]
\textbf{Énoncé.} Une association de jeunes de Bafoussam collecte les vieux cahiers et journaux pour fabriquer du papier hygiénique. (a) Décris les étapes de cette transformation en justifiant chacune. (b) Pourquoi le blanchiment est-il particulièrement important pour ce produit ? (c) Cite deux avantages de cette activité pour la communauté.
\tcblower
\textbf{Solution.}

(a) \textbf{Étapes et justifications.}
\begin{enumerate}
\item \textbf{Tri} : on écarte les papiers plastifiés, les agrafes et les papiers gras, qui dégraderaient la qualité de la pâte.
\item \textbf{Découpage et trempage} dans l'eau : l'eau ramollit les fibres de cellulose.
\item \textbf{Désencrage et malaxage} : les produits détachent l'encre des fibres ; le malaxeur transforme le tout en pâte homogène.
\item \textbf{Filtration} : on élimine les impuretés restantes (particules d'encre, sable).
\item \textbf{Blanchiment} : la pâte est éclaircie pour obtenir un papier blanc.
\item \textbf{Formage sur tamis} : la pâte étalée en fine couche laisse s'écouler l'eau ; les fibres s'enchevêtrent.
\item \textbf{Pressage et séchage} : la feuille est essorée puis séchée.
\item \textbf{Découpe et enroulement} en rouleaux de papier hygiénique.
\end{enumerate}

(b) \textbf{Importance du blanchiment :} le papier hygiénique entre en contact avec le corps ; il doit être propre, sans encre résiduelle (potentiellement toxique) et d'aspect blanc, gage d'hygiène pour les consommateurs.

(c) \textbf{Avantages :} réduction des déchets et des feux sauvages de papier dans les rues ; création d'emplois et de revenus pour les jeunes ; production d'un bien de première nécessité à faible coût ; préservation des arbres (moins de cellulose vierge).

\textbf{Conclusion :} le recyclage du papier en papier hygiénique suit la filière tri $\rightarrow$ pâte $\rightarrow$ épuration-blanchiment $\rightarrow$ formage $\rightarrow$ séchage $\rightarrow$ conditionnement, avec un double bénéfice sanitaire et économique.
\end{exoresolu}

\courssub{Méthode 4 : Expliquer le recyclage des déchets plastiques en bouteilles}

\begin{methode}[Décrire le recyclage des plastiques en bouteilles]
\begin{enumerate}
\item \textbf{Collecte et tri par type de résine} : les bouteilles sont en \cle{PET} (polyéthylène téréphtalate, code 1) ; il faut séparer le PET des autres plastiques (PEHD, PVC) car leurs températures de fusion et propriétés diffèrent.
\item \textbf{Nettoyage} : lavage pour éliminer étiquettes, colles et résidus ; retrait des bouchons (souvent en PEHD).
\item \textbf{Broyage} : les bouteilles sont réduites en paillettes (flakes).
\item \textbf{Lavage et séchage des paillettes} : séparation par densité, rinçage, séchage.
\item \textbf{Fusion et extrusion} : les paillettes sont fondues (vers \SI{250}{\celsius} pour le PET) puis transformées en granulés ou directement en préformes.
\item \textbf{Soufflage} : les préformes chauffées sont soufflées dans un moule pour donner de nouvelles bouteilles.
\item \textbf{Contrôle qualité et conditionnement.}
\end{enumerate}
Vocabulaire exigé : \cle{thermoplastique} (plastique qui se ramollit à la chaleur et peut être refondu), à opposer aux thermodurcissables, non recyclables par fusion.
\end{methode}

\begin{exoresolu}[Recyclage du PET à Yaoundé]
\textbf{Énoncé.} Une entreprise de Yaoundé collecte les bouteilles plastiques usagées pour fabriquer des bouteilles recyclées. (a) Pourquoi seuls les plastiques thermoplastiques comme le PET peuvent-ils être recyclés par cette technique ? (b) Décris les étapes du procédé. (c) La collecte mensuelle est de \SI{2}{tonnes} de bouteilles ; les pertes au tri et au lavage représentent 15 \% de la masse, et une bouteille recyclée pèse \SI{25}{g}. Combien de bouteilles l'entreprise peut-elle produire par mois ?
\tcblower
\textbf{Solution.}

(a) Les plastiques \cle{thermoplastiques} comme le PET se ramollissent et fondent lorsqu'on les chauffe, puis se solidifient en refroidissant : on peut donc les refondre et les remouler. Les plastiques thermodurcissables, eux, se décomposent à la chaleur sans fondre et ne peuvent pas être remoulés.

(b) \textbf{Étapes :} collecte et tri du PET $\rightarrow$ retrait des bouchons et étiquettes $\rightarrow$ lavage $\rightarrow$ broyage en paillettes $\rightarrow$ lavage et séchage des paillettes $\rightarrow$ fusion-extrusion en préformes $\rightarrow$ soufflage en bouteilles $\rightarrow$ contrôle qualité.

(c) \textbf{Calcul.}
Masse utile après pertes :
\[ m = \SI{2000}{kg} \times (1 - 0{,}15) = \SI{2000}{kg} \times 0{,}85 = \SI{1700}{kg} = \SI{1700000}{g} \]
Nombre de bouteilles :
\[ N = \dfrac{\num{1700000}}{\num{25}} = \num{68000} \text{ bouteilles} \]
\textbf{Conclusion :} l'entreprise peut produire environ \num{68000} bouteilles recyclées par mois, tout en retirant \SI{2}{tonnes} de plastique de l'environnement urbain.
\end{exoresolu}

\courssub{Méthode 5 : Expliquer la transformation des déchets plastiques en pavés}

\begin{methode}[Décrire la fabrication de pavés plastique-sable]
\begin{enumerate}
\item \textbf{Collecte et tri} : ramasser les sachets et emballages plastiques (polyéthylène), les nettoyer et les sécher.
\item \textbf{Préparation du sable} : tamiser le sable pour obtenir un granulat propre.
\item \textbf{Fusion du plastique} : chauffer le plastique dans une cuve jusqu'à obtention d'une masse pâteuse (le plastique fondu joue le rôle de \cle{liant}, comme le ciment dans le béton).
\item \textbf{Malaxage} : incorporer le sable chaud au plastique fondu (proportions usuelles : environ 30 \% de plastique pour 70 \% de sable).
\item \textbf{Moulage} : verser le mélange dans des moules à pavés et compacter.
\item \textbf{Refroidissement et démoulage} : le pavé durcit en refroidissant.
\end{enumerate}
Points à citer : pavés plus résistants et moins chers que les pavés en ciment, imperméables ; technique accessible aux PME locales ; elle valorise les sachets plastiques non recyclables en bouteilles.
\end{methode}

\begin{exoresolu}[Fabrication de pavés à Garoua]
\textbf{Énoncé.} Une PME de Garoua fabrique des pavés avec un mélange contenant 30 \% de plastique fondu et 70 \% de sable. (a) Explique le rôle du plastique fondu dans ce matériau et décris les étapes de fabrication. (b) Pour produire \SI{1,5}{tonne} de mélange par jour, quelles masses de plastique et de sable faut-il ? (c) Un sachet plastique pèse en moyenne \SI{5}{g} ; combien de sachets sont valorisés chaque jour ?
\tcblower
\textbf{Solution.}

(a) Le plastique fondu enrobe les grains de sable et, en refroidissant, les solidarise : il joue le rôle de \cle{liant}, à la place du ciment. Étapes : collecte et nettoyage des sachets $\rightarrow$ tamisage du sable $\rightarrow$ fusion du plastique $\rightarrow$ malaxage plastique-sable $\rightarrow$ moulage et compactage $\rightarrow$ refroidissement et démoulage.

(b) \textbf{Masses des constituants.}
\begin{align*}
m_{\text{plastique}} &= 0{,}30 \times \SI{1500}{kg} = \SI{450}{kg} \\
m_{\text{sable}} &= 0{,}70 \times \SI{1500}{kg} = \SI{1050}{kg}
\end{align*}
Vérification : $450 + 1050 = \SI{1500}{kg}$.

(c) \textbf{Nombre de sachets valorisés.}
\[ N = \dfrac{\SI{450}{kg}}{\SI{5}{g}} = \dfrac{\num{450000}}{\num{5}} = \num{90000} \text{ sachets} \]
\textbf{Conclusion :} chaque jour, la PME utilise \SI{450}{kg} de plastique (soit \num{90000} sachets) et \SI{1050}{kg} de sable pour produire \SI{1,5}{tonne} de pavés, transformant ainsi une pollution urbaine en matériau de construction durable.
\end{exoresolu}

\courssub{Méthode 6 : Informer et sensibiliser la communauté au recyclage}

\begin{methode}[Construire un message de sensibilisation au recyclage]
\begin{enumerate}
\item \textbf{Partir d'un constat local chiffré} : volume de déchets, durée de vie des plastiques dans la nature (plusieurs centaines d'années), conséquences (bouchage des caniveaux, inondations, paludisme favorisé par les eaux stagnantes).
\item \textbf{Expliquer simplement les solutions} : tri à la source, les 3 R (\cle{Réduire, Réutiliser, Recycler}), les débouchés économiques (pavés, bouteilles, papier recyclé, compost, biogaz).
\item \textbf{Adapter le support au public} : affiches, sketches, radio communautaire en langues locales, démonstrations pratiques dans les écoles et marchés.
\item \textbf{Proposer des actions concrètes et réalisables} : poubelles de tri, points de collecte, journées de nettoyage, coopératives de recyclage.
\item \textbf{Valoriser les bénéfices} : santé (moins de maladies), environnement (sols et eaux préservés), économie (emplois, revenus).
\end{enumerate}
Savoir-être : adopter un ton positif et non culpabilisant, montrer l'exemple, impliquer les chefs traditionnels et les leaders communautaires.
\end{methode}

\begin{exoresolu}[Campagne de sensibilisation dans un quartier de Douala]
\textbf{Énoncé.} Ton club d'environnement organise une campagne de sensibilisation au recyclage dans un quartier de Douala où les caniveaux sont bouchés par les sachets plastiques. (a) Rédige les trois arguments principaux de ton message. (b) Propose deux actions concrètes à mettre en place avec les habitants. (c) Explique le lien entre l'accumulation des déchets plastiques et la santé des populations.
\tcblower
\textbf{Solution.}

(a) \textbf{Arguments du message.}
\begin{enumerate}
\item \textbf{Argument sanitaire} : les sachets bouchent les caniveaux ; les eaux stagnantes deviennent des gîtes larvaires pour les moustiques, ce qui aggrave le paludisme, première cause de consultation dans nos centres de santé.
\item \textbf{Argument environnemental} : un sachet plastique met plusieurs centaines d'années à se dégrader ; il pollue les sols, les cours d'eau et finit dans l'océan.
\item \textbf{Argument économique} : ces déchets sont une richesse — ils peuvent être transformés en pavés, en bouteilles recyclées, en compost, et créer des emplois pour les jeunes du quartier.
\end{enumerate}

(b) \textbf{Actions concrètes.}
\begin{enumerate}
\item Installer des points de collecte des plastiques près du marché, avec un système d'achat au kilo pour encourager le ramassage.
\item Organiser une journée mensuelle de nettoyage (« clean-up day ») suivie d'une démonstration de fabrication de pavés plastique-sable.
\end{enumerate}

(c) \textbf{Lien déchets-santé :} les déchets plastiques retiennent l'eau de pluie, créant des eaux stagnantes où se reproduisent les moustiques \emph{Anopheles} (paludisme) et \emph{Aedes} (dengue) ; les déchets organiques mélangés attirent mouches et rongeurs, vecteurs de maladies diarrhéiques (choléra, typhoïde) ; le brûlage sauvage des plastiques dégage des fumées toxiques responsables d'affections respiratoires.

\textbf{Conclusion :} sensibiliser, c'est relier le geste de tri aux bénéfices concrets pour la santé, l'environnement et le portefeuille de chaque habitant.
\end{exoresolu}

\begin{attention}[Les 5 erreurs qui coûtent des points au Bac]
\begin{enumerate}
\item \textbf{Confondre fermentation et combustion.} La fermentation alcoolique (\ce{C6H12O6 -> 2 C2H5OH + 2 CO2}) se fait \emph{sans oxygène}, grâce aux levures ; la combustion, elle, consomme du dioxygène et produit \ce{CO2} et \ce{H2O}. Ne jamais écrire qu'un biocarburant est « fabriqué par combustion ».
\item \textbf{Oublier le tri par type de plastique.} On ne peut pas fondre ensemble PET, PVC et PEHD : le recyclage en bouteilles exige le tri préalable du PET. Omettre cette étape rend le procédé faux.
\item \textbf{Dire que le plastique des pavés est un « ciment ».} Le plastique fondu est un \emph{liant} qui enrobe le sable ; il ne subit aucune prise chimique comme le ciment. Employer le terme exact : liant.
\item \textbf{Négliger les unités et les conversions dans les calculs.} Convertir systématiquement en unités cohérentes (kg $\leftrightarrow$ g, MJ $\leftrightarrow$ kWh avec $\SI{1}{kWh} = \SI{3,6}{MJ}$) et appliquer les pourcentages de pertes \emph{avant} de calculer le nombre de produits finis.
\item \textbf{Rédiger une liste d'étapes sans justification.} Au Bac, « expliquer » signifie décrire chaque étape \emph{et} dire à quoi elle sert (le trempage ramollit les fibres, le blanchiment garantit l'hygiène, le soufflage donne la forme de la bouteille...). Une simple énumération non commentée fait perdre la moitié des points.
\end{enumerate}
\end{attention}

\newpage

\coursec{Exercices}

\courssub{Je vérifie mes connaissances}

\begin{exercice}[QCM : Biocarburants et biogaz]
Parmi les affirmations suivantes, lesquelles sont \textbf{correctes} ? Justifiez brièvement chaque choix.
\begin{enumerate}
    \item Le bioéthanol de première génération est produit exclusivement à partir de résidus lignocellulosiques (paille, bagasse).
    \item La méthanisation (production de biogaz) est un processus aérobie réalisé par des bactéries méthanogènes.
    \item Le biodiesel (EMHV) est obtenu par transestérification d'huiles végétales ou de graisses animales avec un alcool (souvent le méthanol) en présence d'un catalyseur.
    \item La production de bioéthanol par fermentation alcoolique du glucose par \emph{Saccharomyces cerevisiae} suit l'équation bilan : \ce{C6H12O6 -> 2 C2H5OH + 2 CO2}.
    \item Le biogaz est principalement composé de méthane (\ce{CH4}) et de dioxyde de carbone (\ce{CO2}) ; il peut être utilisé directement comme carburant pour véhicules sans épuration.
\end{enumerate}
\end{exercice}

\begin{exercice}[Vrai ou Faux justifié : Recyclage du papier et du plastique]
Indiquez si les affirmations suivantes sont vraies (V) ou fausses (F). Justifiez chaque réponse par une phrase précise.
\begin{enumerate}
    \item L'étape de désencrage du papier récupéré utilise des tensioactifs et des bulles d'air pour séparer les particules d'encre des fibres de cellulose par flottation.
    \item Le recyclage du papier en papier hygiénique nécessite un blanchiment au chlore gazeux (\ce{Cl2}) pour obtenir la blancheur requise.
    \item Le PET (polyéthylène téréphtalate) est un thermodurcissable, ce qui empêche son recyclage par fusion pour fabriquer de nouvelles bouteilles.
    \item La technique de fabrication de pavés en plastique-sable consiste à fondre le plastique (souvent PE, PP, PS) pour l'utiliser comme liant à la place du ciment, puis à mélanger avec du sable.
    \item La règle des 3R (Réduire, Réutiliser, Recycler) place le recyclage comme la priorité absolue avant la réduction à la source.
\end{enumerate}
\end{exercice}

\begin{exercice}[Définitions et calcul direct : Rendement énergétique]
\begin{enumerate}
    \item Définissez les termes suivants : \textbf{biocarburant de 1\up{re} génération}, \textbf{méthanisation}, \textbf{transestérification}, \textbf{valorisation matière}, \textbf{valorisation énergétique}.
    \item Une usine de méthanisation traite \SI{10}{\tonne} de déchets organiques par jour. Le rendement en biogaz est de \SI{100}{\cubic\meter\per\tonne} de déchets. Le biogaz contient 60 \% de \ce{CH4} (v/v). Sachant que le PCI (Pouvoir Calorifique Inférieur) du méthane est de \SI{35,8}{\mega\joule\per\cubic\meter} (CNTP), calculez l'énergie primaire (en MJ/jour) potentiellement récupérable sous forme de chaleur si le rendement de la chaudière est de 90 \%. 
\end{enumerate}
\end{exercice}

\begin{exercice}[Schéma-bilan : De la canne à sucre au bioéthanol]
Complétez le schéma-bilan ci-dessous en indiquant pour chaque étape : le nom du procédé, le(s) micro-organisme(s) ou enzyme(s) clé(s), le substrat principal et le(s) produit(s) principal(aux).
\begin{center}
\begin{tabularx}{\linewidth}{|c|X|X|X|X|}\hline
\textbf{Étape} & \textbf{Procédé / Opération} & \textbf{Agent biologique / Enzyme} & \textbf{Substrat entrant} & \textbf{Produit(s) sortant(s)} \\ \hline
1 & Broyage / Diffusion & -- & Tiges de canne &  \\ \hline
2 &  & Levures (\emph{S. cerevisiae}) & Jus sucré (saccharose) &  \\ \hline
3 & Distillation & -- & Vin de canne (\SIrange{8}{10}{\percent} vol.) &  \\ \hline
4 & Déshydratation (zéolithe) & -- &  & Bioéthanol carburant (> \SI{99,5}{\percent}) \\ \hline
\end{tabularx}
\end{center}
\end{exercice}

\courssub{Je m'entraîne}

\begin{exercice}[Analyse de courbe : Fermentation du jus de manioc]
On réalise une fermentation du jus de manioc (titré à \SI{150}{\gram\per\liter} de glucose) par \emph{Saccharomyces cerevisiae} en milieu anaérobie, à \SI{30}{\celsius}, pH 4,5. On suit la concentration en éthanol (\SI{}{\gram\per\liter}) et le volume de \ce{CO2} dégagé (\SI{}{\liter}) en fonction du temps (heures).
\begin{center}
\begin{tikzpicture}
\begin{axis}[
    width=\linewidth, height=8cm,
    xlabel={Temps (h)}, ylabel={Concentration / Volume},
    xmin=0, xmax=50, ymin=0, ymax=80,
    legend pos=north west,
    grid=major,
    domain=0:50, samples=200,
]
\addplot[popblue, thick] {70*(1-exp(-0.15*x))}; \addlegendentry{Éthanol (g/L)}
\addplot[poporange, thick, dashed] {50*(1-exp(-0.12*x))}; \addlegendentry{CO$_2$ (L)}
\end{axis}
\end{tikzpicture}
\end{center}
\begin{enumerate}
    \item Identifiez la phase de latence, la phase exponentielle et la phase stationnaire sur la courbe de production d'éthanol. Justifiez par l'évolution de la population de levures et de la concentration en substrat.
    \item Calculez le rendement volumique maximal de production d'éthanol (en \SI{}{\gram\per\liter\per\hour}) entre la 10\up{ème} et la 20\up{ème} heure (lecture graphique approximative acceptée).
    \item Le volume de \ce{CO2} mesuré à la fin de la fermentation (50 h) est de \SI{50}{\liter} pour \SI{100}{\liter} de moût. Vérifiez si ce résultat est cohérent avec la stœchiométrie de la fermentation alcoolique (masses molaires : \ce{C6H12O6}=180, \ce{C2H5OH}=46, \ce{CO2}=44 ; \ce{CO2} gaz aux CNTP : \SI{22,4}{\liter\per\mole}).
    \item Expliquez pourquoi la fermentation s'arrête avant la consommation totale du glucose initial.
\end{enumerate}
\end{exercice}

\begin{exercice}[Ordre et commentaire : Recyclage du papier en papier hygiénique]
Les étapes suivantes, présentées dans le désordre, décrivent la fabrication de papier hygiénique à partir de papiers récupérés (journaux, bureaux). Classez-les dans l'ordre chronologique logique (de 1 à 8) et ajoutez un commentaire technique d'une phrase pour \textbf{chacune} des étapes 2, 4, 6 et 8.
\begin{enumerate}
    \item Raffinage (battage) des fibres.
    \item Presse et séchage sur cylindre chauffé (Yankee dryer).
    \item Tri manuel / optique des papiers (retrait agrafes, plastiques).
    \item Désencrage par flottation (bulles d'air + tensioactifs).
    \item Dépotage / Pulpage dans l'hydropulpeur (eau + agitation).
    \item Blanchiment (peroxyde d'hydrogène \ce{H2O2} ou hydrosulfite).
    \item Formation de la feuille sur toile (machine à papier).
    \item Épuration / Nettoyage (cyclones, tamisage) pour retirer colles, charges, fines.
\end{enumerate}
\end{exercice}

\begin{exercice}[Tableau comparatif : Pavé plastique-sable vs Pavé ciment]
Le tableau ci-dessous présente des caractéristiques moyennes de deux types de pavés utilisés pour le revêtement de voiries secondaires à Douala.
\begin{center}
\begin{tabularx}{\linewidth}{|l|X|X|}\hline
\textbf{Propriété} & \textbf{Pavé Plastique-Sable (recyclé PE/PP)} & \textbf{Pavé Ciment (béton vibré)} \\ \hline
Résistance à la compression (MPa) & 25 -- 35 & 30 -- 50 \\ \hline
Coefficient d'absorption d'eau (\%) & \textless 0,5 & 4 -- 6 \\ \hline
Coût de production (FCFA / m²) & \num{3500} -- \num{4500} & \num{5000} -- \num{7000} \\ \hline
Poids spécifique (kg/m³) & \num{1800} -- \num{2000} & \num{2200} -- \num{2400} \\ \hline
Durée de vie estimée (ans) & 15 -- 20 & 20 -- 30 \\ \hline
Impact CO₂ (kg éq. CO₂ / m²) & \num{15} -- \num{20} & \num{60} -- \num{90} \\ \hline
\end{tabularx}
\end{center}
\begin{enumerate}
    \item Dégagez \textbf{trois avantages majeurs} du pavé plastique-sable par rapport au pavé ciment, en vous appuyant sur des valeurs du tableau.
    \item Décrivez la \textbf{technique de fabrication} du pavé plastique-sable en 5 étapes principales, depuis le déchet trié jusqu'au pavé fini. Précisez la température de fusion typique du mélange PE/PP.
    \item Pourquoi ne utilise-t-on pas du PET (bouteilles d'eau) comme liant principal pour ces pavés ? (Indice : température de fusion, viscosité).
\end{enumerate}
\end{exercice}

\begin{exercice}[Filières de valorisation : Composition des déchets de Yaoundé]
La composition moyenne en masse des ordures ménagères collectées à Yaoundé est la suivante :
\begin{itemize}
    \item Matières organiques (déchets de cuisine, marchés) : 55 \%
    \item Plastiques (sachets, bouteilles, emballages) : 15 \%
    \item Papiers / Cartons : 10 \%
    \item Verres / Métaux / Textiles / Autres : 20 \%
\end{itemize}
Proposez \textbf{trois filières de valorisation adaptées} à cette composition. Pour chaque filière, précisez :
\begin{enumerate}
    \item Le type de déchet ciblé (flux d'entrée).
    \item La technique de valorisation (nom + principe en une phrase).
    \item Le(s) produit(s) final(aux) valorisé(s).
    \item Un avantage environnemental et un avantage socio-économique local.
\end{enumerate}
\end{exercice}

\courssub{Je me prépare au Bac}

\begin{exercice}[Sujet type Bac : Production de bioéthanol à partir de résidus de banane plantain (Document d'analyse)]
\textbf{Document 1 :} Contexte camerounais. Le Cameroun produit environ \SI{4,5}{\million\tonne} de bananes plantains par an. Les résidus (épluchures, rejets, pseudotiges) représentent \SI{40}{\%} de la biomasse et sont souvent abandonnés ou brûlés. Une unité pilote de bioéthanol de 2\up{ème} génération est envisagée à Njombé (région du Littoral) pour valoriser ces résidus lignocellulosiques.

\textbf{Document 2 :} Principales étapes du procédé (ordre chronologique).
\begin{enumerate}
    \item Prétraitement : Broyage fin + traitement vapeur explosion (\SI{190}{\celsius}, 5 min) ou acide dilué (\ce{H2SO4} \SI{1}{\%}, \SI{121}{\celsius}, 30 min) pour rompre la structure lignine-hémicellulose-cellulose.
    \item Hydrolyse enzymatique : Ajout d'un cocktail enzymatique (cellulases, hémicellulases, $\beta$-glucosidases) à \SI{50}{\celsius}, pH 4,8 pendant 48--72 h. Conversion de la cellulose en glucose et de l'hémicellulose en xylose/arabinose.
    \item Fermentation C5/C6 : Inoculation d'une souche recombinante de \emph{Saccharomyces cerevisiae} (ou \emph{Zymomonas mobilis}) capable de fermenter à la fois le glucose (C6) et le xylose (C5) en éthanol. Anaérobiose, \SI{32}{\celsius}, 24--48 h.
    \item Distillation + Déshydratation : Séparation de l'éthanol (azeotrope à 95,6 \%), puis passage sur zéolithe 3A pour obtenir du bioéthanol carburant (\textgreater 99,5 \%).
    \item Valorisation des co-produits : Le résidu solide (lignine + enzymes non réagies) $\to$ chaudière biomasse (énergie process). Le vinasse (liquide) $\to$ fertirrigation ou méthanisation.
\end{enumerate}

\textbf{Document 3 :} Courbe de production d'éthanol lors de l'hydrolyse-fermentation simultanée (SHF) vs simultanée (SSF) à l'échelle pilote (\SI{100}{\liter}).
\begin{center}
\begin{tikzpicture}
\begin{axis}[
    width=\linewidth, height=7cm,
    xlabel={Temps (h)}, ylabel={Concentration éthanol (g/L)},
    xmin=0, xmax=72, ymin=0, ymax=55,
    legend pos=south east, grid=major,
]
\addplot[popblue, thick, mark=*, mark repeat=10, samples=10] coordinates {(0,0) (12,15) (24,32) (36,42) (48,48) (60,50) (72,51)};
\addlegendentry{SSF (Simultaneous Saccharification \& Fermentation)}
\addplot[poporange, thick, mark=square*, mark repeat=10, samples=10] coordinates {(0,0) (12,5) (24,20) (36,35) (48,44) (60,48) (72,49)};
\addlegendentry{SHF (Separate Hydrolysis \& Fermentation)}
\end{axis}
\end{tikzpicture}
\end{center}

\textbf{Questions :}
\begin{enumerate}
    \item \textbf{Classement :} Replacez les étapes du Document 2 dans l'ordre logique (1 à 5) si elles étaient mélangées. (Ici, elles sont déjà dans l'ordre ; confirmez la logique).
    \item \textbf{Prétraitement :} Expliquez pourquoi le prétraitement (étape 1) est \emph{indispensable} pour la biomasse lignocellulosique (épluchures de plantain) alors qu'il est inutile pour le jus de canne ou le grain de maïs. Citez les trois polymères cibles.
    \item \textbf{Hydrolyse enzymatique :} Écrivez l'équation bilan simplifiée de l'hydrolyse de la cellulose (polymère du glucose) par les cellulases. Pourquoi l'ajout de $\beta$-glucosidase est-il crucial pour le rendement ?
    \item \textbf{Fermentation C5/C6 :} Pourquoi une levure \emph{S. cerevisiae} sauvage ne peut-elle pas valoriser le xylose (C5) ? Quel avantage métabolique présente la souche recombinante utilisée ?
    \item \textbf{Analyse de courbe (Doc 3) :} Comparez les deux procédés (SSF vs SHF) au regard de : (a) la cinétique de production (vitesse initiale), (b) le rendement final (titre éthanol à 72 h), (c) l'inhibition enzymatique par le glucose. Concluez sur le procédé le plus performant pour cette usine.
    \item \textbf{Bilan énergétique/environnemental :} Calculez la production théorique maximale d'éthanol (en litres) à partir de \SI{1}{\tonne} d'épluchures de plantain (teneur en cellulose+hémicellulose = 60 \% MS ; rendement hydrolyse/fermentation global = 80 \% théorique ; densité éthanol = \SI{0,789}{\kilogram\per\liter}). Discutez brièvement de l'intérêt de la valorisation énergétique du résidu solide (lignine) pour l'autonomie de l'usine.
\end{enumerate}
\end{exercice}

\begin{exercice}[Sujet type Bac : Gestion intégrée des déchets plastiques à Douala - Étude de cas]
La ville de Douala génère environ \SI{800}{\tonne\per\day} de déchets plastiques (majoritairement sachets PE-LD, bouteilles PET, bouchons PP, pots PS). La mairie lance un appel à projets pour une unité de valorisation locale. Trois options techniques sont proposées par des start-ups camerounaises :
\begin{itemize}
    \item \textbf{Option A :} Recyclage mécanique « bouteille à bouteille » (r-PET) pour l'industrie agroalimentaire (brasseries, eaux minérales).
    \item \textbf{Option B :} Fabrication de pavés « éco-pavés » plastique-sable (PE/PP/PS) pour le pavage des voiries et cours d'écoles.
    \item \textbf{Option C :} Pyrolyse catalytique des plastiques mixtes non triés (PE, PP, PS) pour produire du carburant liquide (diesel/gazole de synthèse) et du gaz de process.
\end{itemize}

\textbf{Données techniques complémentaires :}
\begin{itemize}
    \item PET : $T_{\text{f}} \approx \SI{250}{\celsius}$, sensible à l'hydrolyse (eau) $\to$ baisse de viscosité (IV). Norme alimentaire exige IV \textgreater \SI{0,75}{\deci\liter\per\gram}.
    \item PE/PP/PS : $T_{\text{f}} \approx \SI{110}{\celsius}$--$\SI{160}{\celsius}$, hydrophobes, stables thermiquement jusqu'à \SI{300}{\celsius}.
    \item Pyrolyse : \SI{400}{\celsius}--\SI{500}{\celsius}, absence d'O₂, catalyseur (zéolithe). Rendement liquide \approx 70 \% masse, Gaz 20 \%, Char 10 \%.
    \item Coût investissement : A \textgreater B \textgreater C. Complexité technique : A \textgreater C \textgreater B.
\end{itemize}

\textbf{Questions :}
\begin{enumerate}
    \item \textbf{Triage amont :} Expliquez pourquoi un tri rigoureux (séparation PET / PE-PP-PS) est \emph{critique} pour la réussite de l'Option A, mais moins critique pour l'Option B. Que se passe-t-il si 5 \% de PVC contaminent le flux PET en extrusion ?
    \item \textbf{Option A (r-PET) :} Décrivez les 6 étapes principales du recyclage mécanique « bouteille à bouteille » (de la balle de bouteilles triées aux granulés r-PET certifiés contact alimentaire). Indiquez l'étape où l'on restaure la viscosité intrinsèque (IV) et le principe utilisé (SPC - Solid State Polycondensation).
    \item \textbf{Option B (Pavés) :} Détaillez le procédé de fabrication d'un pavé plastique-sable (dosage typique plastique/sable, température de mélange, moulage, refroidissement). Donnez \textbf{deux avantages techniques} du pavé plastique-sable sur le pavé ciment pour le climat tropical humide (Douala).
    \item \textbf{Option C (Pyrolyse) :} Écrivez l'équation bilan globale simplifiée de la pyrolyse du polyéthylène (\ce{-(CH2-CH2)-_n}) en mélange d'alcanes (\ce{C_nH_{2n+2}}) et alcènes (\ce{C_nH_{2n}}). Citez \textbf{un avantage majeur} (valorisation flux sales/mixtes) et \textbf{un risque majeur} (pollution, économie) de cette filière par rapport au recyclage mécanique.
    \item \textbf{Choix stratégique :} En tant qu'expert SVT consultant pour la mairie, recommandez \textbf{une stratégie combinée} (ex: 50\% Option A, 30\% Option B, 20\% Option C) justifiée par la composition du gisement, la maturité technique, la création d'emplois locaux et la circularité matière vs énergie. Argumentez en 10 lignes maximum.
\end{enumerate}
\end{exercice}

\begin{exercice}[Sujet type Bac : Communication pour le changement de comportement - Sensibilisation au tri]
Dans le cadre de la semaine nationale de l'environnement au Cameroun, vous êtes chargé(e) de concevoir un message de sensibilisation grand public pour promouvoir le tri à la source et le recyclage des déchets plastiques dans les grands marchés (ex: Marché Central de Yaoundé, Marché Sandaga à Douala).

\textbf{Contraintes du message :}
\begin{itemize}
    \item Format : Affiche \textbf{OU} Spot radio de 45 secondes (choisissez l'un).
    \item Cible : Commerçants, acheteurs, porteurs (jeunes, peu scolarisés, multilingues : français, pidgin, langues locales).
    \item Objectifs : (1) Expliquer \emph{pourquoi} trier (impact canalisations, inondations, santé, ressource). (2) Montrer \emph{comment} trier (gestes simples : 2 bacs/poubelles). (3) Valoriser l'acte (gain argent, propreté, emploi jeunes). (4) Intégrer la règle des 3R.
\end{itemize}

\textbf{Travail demandé :}
\begin{enumerate}
    \item \textbf{Choix du support :} Justifiez en 3 lignes le choix de l'affiche ou du spot radio pour ce contexte de marché camerounais.
    \item \textbf{Contenu du message :} Rédigez le \textbf{texte intégral} de votre message (slogan, corps du texte, appel à l'action, contacts utiles). Utilisez un langage clair, imagé, percutant, intégrant des termes locaux (ex: « plastique », « sachet », « bouteille », « caniveaux », « inondation », « argent », « jeunesse »).
    \item \textbf{Argumentaire scientifique :} En annexe de votre message, rédigez un paragraphe (5--6 lignes) à l'intention des animateurs/relais communautaires, expliquant \textbf{scientifiquement} pourquoi le sac en PE-LD (sachet « noir » ou « transparent ») mis en décharge ou brûlé à l'air libre est une perte de matière (carbone fossile) et une source de polluants (dioxines, furanes, microplastiques, lixiviats), alors qu'il est parfaitement recyclable en pavés ou granulés.
\end{enumerate}
\end{exercice}

\newpage

\coursec{Corrigés des exercices}

\begin{corrige}[Exercice 1 — QCM : Biocarburants et biogaz]
\textbf{Réponses correctes : 3 et 4.}

\begin{enumerate}
    \item \textbf{FAUX.} Le bioéthanol de \cle{première génération} est produit à partir de matières \textbf{sucrières} (canne à sucre, betterave, mélasse) ou \textbf{amylacées} (maïs, blé, manioc), c'est-à-dire des ressources alimentaires. Les résidus lignocellulosiques (paille, bagasse, épluchures de plantain) relèvent de la \textbf{deuxième génération}, qui exige un prétraitement et une hydrolyse enzymatique de la cellulose.
    \item \textbf{FAUX.} La méthanisation est un processus strictement \textbf{anaérobie} : les bactéries méthanogènes (archées méthanogènes, ex. \emph{Methanobacterium}) sont tuées par l'oxygène. Elle se déroule en digesteur étanche, en quatre phases (hydrolyse, acidogenèse, acétogenèse, méthanogenèse).
    \item \textbf{VRAI.} Le biodiesel (EMHV : esters méthyliques d'huiles végétales) est obtenu par \cle{transestérification} : un triglycéride réagit avec du méthanol en présence d'un catalyseur (soude \ce{NaOH} ou potasse \ce{KOH}) pour donner 3 molécules d'esters méthyliques + 1 molécule de glycérol (co-produit).
    \item \textbf{VRAI.} La fermentation alcoolique du glucose par \emph{Saccharomyces cerevisiae} suit l'équation de Gay-Lussac : \ce{C6H12O6 -> 2 C2H5OH + 2 CO2}, avec un rendement théorique massique de $2 \times 46 / 180 \approx 0,51$ g d'éthanol par gramme de glucose.
    \item \textbf{FAUX.} Le biogaz brut (50--70 \% \ce{CH4}, 30--50 \% \ce{CO2}, traces de \ce{H2S}) ne peut pas alimenter directement un véhicule : il faut une \textbf{épuration} (élimination du \ce{CO2}, du \ce{H2S} corrosif et de l'eau) pour obtenir du \textbf{biométhane} ($>$ 95 \% \ce{CH4}), injectable au réseau ou utilisable comme carburant (GNV).
\end{enumerate}

\textbf{Points de vigilance :} ne pas confondre générations de biocarburants (1G = alimentaire ; 2G = lignocellulose ; 3G = microalgues) ; « méthanogène » n'implique pas « aérobie » — c'est l'inverse.
\end{corrige}

\begin{corrige}[Exercice 2 — Vrai ou Faux : Recyclage du papier et du plastique]
\begin{enumerate}
    \item \textbf{VRAI.} Le \cle{désencrage} par flottation repose sur l'injection de bulles d'air dans la pâte diluée en présence de tensioactifs : les particules d'encre, rendues hydrophobes, s'accrochent aux bulles et remontent en mousse qui est écumée, tandis que les fibres de cellulose hydrophiles restent en suspension.
    \item \textbf{FAUX.} Le blanchiment au chlore gazeux \ce{Cl2} est \textbf{proscrit} car il génère des composés organochlorés toxiques (dioxines). On utilise aujourd'hui le peroxyde d'hydrogène \ce{H2O2}, l'ozone \ce{O3} ou l'hydrosulfite de sodium (procédés TCF/ECF), suffisants pour le papier hygiénique.
    \item \textbf{FAUX.} Le PET est un \textbf{thermoplastique} : il se ramollit à la chaleur ($T_f \approx \SI{250}{\celsius}$) et peut être refondu et remoulé, ce qui autorise le recyclage mécanique « bouteille à bouteille ». Les thermodurcissables (bakélite, époxy), eux, se dégradent sans fondre et ne sont pas recyclables par fusion.
    \item \textbf{VRAI.} Dans la technique des pavés plastique-sable, les plastiques fusibles (PE, PP, PS) sont fondus vers \SI{160}{\celsius}--\SI{180}{\celsius} et servent de \textbf{liant} en remplacement total du ciment ; le sable chaud est ensuite incorporé, puis le mélange est moulé et refroidi.
    \item \textbf{FAUX.} La hiérarchie des 3R place la \textbf{Réduction} à la source en priorité absolue, puis la \textbf{Réutilisation}, et le \textbf{Recyclage} en dernier recours avant l'élimination. Recycler ne dispense pas de consommer moins.
\end{enumerate}

\textbf{Points de vigilance :} bien distinguer thermoplastique (recyclable par fusion) et thermodurcissable (non refondable) ; l'ordre des 3R est une hiérarchie, pas une liste équivalente.
\end{corrige}

\begin{corrige}[Exercice 3 — Définitions et calcul : Rendement énergétique]
\textbf{1. Définitions}
\begin{itemize}
    \item \textbf{Biocarburant de 1\up{re} génération :} carburant liquide (bioéthanol, biodiesel) produit à partir de la fraction sucrée, amylacée ou oléagineuse de cultures \textbf{alimentaires} (canne, betterave, maïs, palme, soja).
    \item \textbf{Méthanisation :} dégradation biologique de la matière organique en \textbf{anaérobiose} par un consortium microbien, produisant un biogaz riche en méthane et un digestat fertilisant.
    \item \textbf{Transestérification :} réaction chimique entre un triglycéride (huile ou graisse) et un alcool léger (méthanol) en présence d'un catalyseur, donnant des esters d'acides gras (biodiesel) et du glycérol.
    \item \textbf{Valorisation matière :} recyclage d'un déchet en nouvelle \textbf{matière première} (papier recyclé, granulés r-PET, pavés).
    \item \textbf{Valorisation énergétique :} récupération de l'\textbf{énergie} contenue dans le déchet (incinération avec récupération de chaleur, méthanisation, pyrolyse).
\end{itemize}

\textbf{2. Calcul de l'énergie récupérable}
\begin{align*}
V_{\text{biogaz}} &= \SI{10}{\tonne\per\day} \times \SI{100}{\cubic\meter\per\tonne} = \SI{1000}{\cubic\meter\per\day} \\
V_{\ce{CH4}} &= 0{,}60 \times 1000 = \SI{600}{\cubic\meter\per\day} \\
E_{\text{brute}} &= 600 \times 35{,}8 = \SI{21480}{\mega\joule\per\day} \\
E_{\text{utile}} &= 0{,}90 \times 21480 = \boxed{\SI{19332}{\mega\joule\per\day} \approx \SI{1,93e4}{\mega\joule\per\day}}
\end{align*}
Soit environ \SI{19,3}{\giga\joule\per\day} de chaleur utile (équivalent au pouvoir calorifique d'environ \SI{460}{\kilogram} de fioul par jour, PCI fioul $\approx$ \SI{42}{\mega\joule\per\kilogram}).

\textbf{Points de vigilance :} appliquer le pourcentage de \ce{CH4} \emph{avant} le PCI, et le rendement de chaudière \emph{en dernier} ; ne pas oublier les unités à chaque ligne.
\end{corrige}

\begin{corrige}[Exercice 4 — Schéma-bilan : De la canne à sucre au bioéthanol]
\begin{center}
\begin{tabularx}{\linewidth}{|c|X|X|X|X|}\hline
\textbf{Étape} & \textbf{Procédé} & \textbf{Agent biologique} & \textbf{Substrat entrant} & \textbf{Produit(s) sortant(s)} \\ \hline
1 & Broyage / Diffusion & -- & Tiges de canne & \textbf{Jus sucré (saccharose) + bagasse} \\ \hline
2 & \textbf{Fermentation alcoolique} (anaérobie, \SIrange{30}{32}{\celsius}, pH 4,5) & Levures (\emph{S. cerevisiae}) & Jus sucré (saccharose) & \textbf{Vin de canne (éthanol 8--10 \% vol.) + \ce{CO2}} \\ \hline
3 & Distillation & -- & Vin de canne & \textbf{Éthanol hydraté (azéotrope $\approx$ 95--96 \%) + vinasse} \\ \hline
4 & Déshydratation (zéolithe 3A) & -- & \textbf{Éthanol hydraté ($\approx$ 96 \%)} & Bioéthanol carburant ($>$ \SI{99,5}{\percent}) \\ \hline
\end{tabularx}
\end{center}

\textbf{Justifications :}
\begin{itemize}
    \item Étape 1 : le broyage libère le jus ; la bagasse (fibre résiduelle) est un co-produit valorisé en combustible de chaudière.
    \item Étape 2 : la levure hydrolyse d'abord le saccharose en glucose + fructose (enzyme \textbf{invertase}), puis fermente les hexoses selon \ce{C6H12O6 -> 2 C2H5OH + 2 CO2}.
    \item Étape 3 : la distillation exploite la différence de volatilité éthanol/eau, mais bute sur l'\textbf{azéotrope} à 95,6 \% : impossible d'aller au-delà par simple distillation.
    \item Étape 4 : la zéolithe 3A (tamis moléculaire) adsorbe sélectivement les molécules d'eau, permettant d'atteindre le grade carburant anhydre ($>$ 99,5 \%).
\end{itemize}

\textbf{Points de vigilance :} citer l'azéotrope pour justifier l'étape 4 ; ne pas écrire que la distillation donne directement l'éthanol absolu.
\end{corrige}

\begin{corrige}[Exercice 5 — Analyse de courbe : Fermentation du jus de manioc]
\textbf{1. Phases de la courbe d'éthanol}
\begin{itemize}
    \item \textbf{Phase de latence (0 -- $\approx$ 5 h) :} production quasi nulle ; la population de levures s'adapte au milieu (synthèse d'enzymes, multiplication démarrant), le substrat est encore abondant.
    \item \textbf{Phase exponentielle ($\approx$ 5 -- 25 h) :} pente maximale ; les levures se multiplient activement et consomment massivement le glucose ; l'éthanol et le \ce{CO2} s'accumulent rapidement.
    \item \textbf{Phase stationnaire ($>$ 35 h) :} plateau vers \SI{70}{\gram\per\liter} ; le glucose devient limitant et l'éthanol accumulé inhibe les levures (toxicité membranaire) ; la population cesse de croître.
\end{itemize}

\textbf{2. Rendement volumique maximal entre 10 h et 20 h}

Lecture graphique : $[\text{éthanol}]_{10\,\text{h}} \approx \SI{54}{\gram\per\liter}$ ; $[\text{éthanol}]_{20\,\text{h}} \approx \SI{66}{\gram\per\liter}$.
\[ r = \frac{66 - 54}{20 - 10} = \frac{12}{10} \approx \boxed{\SI{1,2}{\gram\per\liter\per\hour}} \]

\textbf{3. Cohérence du volume de \ce{CO2}}

Vérifions à partir de l'éthanol réellement produit (plateau $\approx$ \SI{70}{\gram\per\liter} pour \SI{100}{\liter}) :
\begin{align*}
m_{\text{éthanol}} &= 70 \times 100 = \SI{7000}{\gram} \\
n_{\text{éthanol}} &= \frac{7000}{46} \approx \SI{152,2}{\mole} \\
n_{\ce{CO2}} &= n_{\text{éthanol}} = \SI{152,2}{\mole} \quad (\text{stœchiométrie 1:1}) \\
V_{\ce{CO2}} &= 152{,}2 \times 22{,}4 \approx \SI{3409}{\liter}
\end{align*}
Le volume mesuré (\SI{50}{\liter}) est \textbf{très inférieur} aux $\approx$ \SI{3400}{\liter} attendus : le résultat est \textbf{incohérent} avec la stœchiométrie. Explications possibles : fuite du dispositif de recueil, dissolution partielle du \ce{CO2} dans le moût (acidification), ou erreur de mesure. (À titre indicatif, la fermentation \emph{complète} des \SI{150}{\gram\per\liter} donnerait $150 \times 100 / 180 \times 2 \times 22{,}4 \approx \SI{3733}{\liter}$.)

\textbf{4. Arrêt avant épuisement du glucose}

L'éthanol accumulé ($>$ \SI{70}{\gram\per\liter}) exerce une \textbf{toxicité} sur les levures : il fluidifie et désorganise la membrane plasmique, dénature les enzymes glycolytiques et inhibe la croissance. C'est une \textbf{inhibition par le produit} : la fermentation s'auto-arrête même s'il reste du glucose, d'où la nécessité de souches tolérantes ou de procédés d'extraction continue de l'éthanol.

\textbf{Points de vigilance :} une vérification stœchiométrique doit partir d'une grandeur \emph{mesurée} (ici l'éthanol) ; conclure explicitement « cohérent / incohérent » avec une valeur de comparaison chiffrée.
\end{corrige}

\begin{corrige}[Exercice 6 — Ordre et commentaire : Recyclage du papier en papier hygiénique]
\textbf{Ordre chronologique correct :}

\begin{center}
\begin{tabularx}{\linewidth}{|c|X|}\hline
\textbf{Rang} & \textbf{Étape} \\ \hline
1 & Tri manuel / optique (c) \\ \hline
2 & Dépotage / Pulpage dans l'hydropulpeur (e) \\ \hline
3 & Épuration / Nettoyage : cyclones, tamisage (h) \\ \hline
4 & Désencrage par flottation (d) \\ \hline
5 & Blanchiment (\ce{H2O2} ou hydrosulfite) (f) \\ \hline
6 & Raffinage / battage des fibres (a) \\ \hline
7 & Formation de la feuille sur toile (g) \\ \hline
8 & Presse et séchage sur cylindre Yankee (b) \\ \hline
\end{tabularx}
\end{center}

\textbf{Commentaires techniques demandés :}
\begin{itemize}
    \item \textbf{Étape 2 (pulpage) :} l'hydropulpeur désagrège les papiers dans l'eau chaude ($\approx$ \SI{50}{\celsius}) sous agitation mécanique, séparant les fibres de cellulose sans les couper, pour obtenir une pâte diluée à 4--6 \% de matière sèche.
    \item \textbf{Étape 4 (désencrage) :} les bulles d'air injectées en présence de tensioactifs capturent les particules d'encre hydrophobes qui remontent en mousse écumée en surface, tandis que les fibres hydrophiles demeurent en suspension — c'est une séparation par différence de mouillabilité.
    \item \textbf{Étape 6 (raffinage) :} le battage mécanique fibrille la surface des fibres de cellulose, ce qui augmente les liaisons hydrogène inter-fibres et donc la résistance et la souplesse de la feuille future.
    \item \textbf{Étape 8 (séchage) :} la feuille humide est pressée puis crêpée sur le cylindre Yankee chauffé à vapeur ($\approx$ \SI{100}{\celsius}--\SI{130}{\celsius}) : le crêpage confère au papier hygiénique sa souplesse et son pouvoir absorbant caractéristiques.
\end{itemize}

\textbf{Points de vigilance :} le désencrage ne peut précéder le pulpage (il faut une pâte) ; le blanchiment se place après le désencrage pour ne pas gaspiller de réactif sur l'encre ; le raffinage précède obligatoirement la formation de la feuille.
\end{corrige}

\begin{corrige}[Exercice 7 — Tableau comparatif : Pavé plastique-sable vs Pavé ciment]
\textbf{1. Trois avantages majeurs du pavé plastique-sable}
\begin{itemize}
    \item \textbf{Imperméabilité :} absorption d'eau $<$ 0,5 \% contre 4--6 \% pour le béton : le pavé plastique résiste mieux à l'humidité, aux pluies tropicales et à la fissuration par cycles mouillage/séchage (atout majeur à Douala).
    \item \textbf{Coût :} \num{3500}--\num{4500} FCFA/m² contre \num{5000}--\num{7000} FCFA/m², soit une économie de l'ordre de 30--40 \%.
    \item \textbf{Bilan carbone :} 15--20 kg éq. \ce{CO2}/m² contre 60--90 kg éq. \ce{CO2}/m² (le ciment est très émetteur), soit une division par 3 à 4 de l'empreinte carbone. (On peut aussi citer la légèreté : \num{1800}--\num{2000} kg/m³ contre \num{2200}--\num{2400} kg/m³, facilitant transport et pose.)
\end{itemize}

\textbf{2. Technique de fabrication (5 étapes)}
\begin{enumerate}
    \item \textbf{Tri et lavage} des déchets plastiques (PE, PP, PS), élimination du PET et du PVC, puis séchage.
    \item \textbf{Broyage} en paillettes de quelques millimètres.
    \item \textbf{Fusion-malaxage} dans une cuve chauffée à \SI{160}{\celsius}--\SI{180}{\celsius} (plage de fusion du PE/PP) jusqu'à obtention d'une pâte homogène servant de liant.
    \item \textbf{Incorporation du sable} chaud et tamisé (dosage typique : 25--35 \% plastique pour 65--75 \% sable en masse) et malaxage jusqu'à enrobage complet des grains.
    \item \textbf{Moulage à chaud} dans des moules métalliques, compactage (presse), puis \textbf{démoulage et refroidissement} à l'air ou à l'eau.
\end{enumerate}

\textbf{3. Pourquoi pas de PET comme liant ?}

Le PET fond vers \SI{250}{\celsius}, température bien plus élevée que celle du PE/PP : il exige plus d'énergie et, surtout, son point de fusion est proche de sa température de dégradation ; il se \textbf{hydrolyse} facilement en présence d'humidité à chaud, ce qui fait chuter sa viscosité (chaîne polymère rompue) et donne un liant cassant et de mauvaise qualité. Sa viscosité à l'état fondu reste en outre élevée, ce qui nuit à l'enrobage du sable. Le PE/PP, fusibles dès \SIrange{110}{160}{\celsius} et hydrophobes, sont bien mieux adaptés.

\textbf{Points de vigilance :} appuyer chaque avantage sur une valeur chiffrée du tableau ; citer une plage de température de fusion cohérente avec le polymère choisi.
\end{corrige}

\begin{corrige}[Exercice 8 — Filières de valorisation : Déchets de Yaoundé]
\begin{center}
\begin{tabularx}{\linewidth}{|l|X|X|X|}\hline
\textbf{Filière} & \textbf{Flux ciblé} & \textbf{Technique et principe} & \textbf{Produits et avantages} \\ \hline
\textbf{1. Compostage / Méthanisation} & Matières organiques (55 \%) : déchets de cuisine et de marchés & \textbf{Compostage} aérobie (dégradation par micro-organismes en andains retournés) ou \textbf{méthanisation} anaérobie en digesteur & Compost (amendement agricole) et/ou biogaz (cuisson) + digestat. \emph{Env. :} détourne de la décharge le flux majeur, réduit \ce{CH4} sauvage et odeurs. \emph{Socio-éco. :} compost bon marché pour maraîchers périurbains, emplois de collecte/tri. \\ \hline
\textbf{2. Recyclage plastique} & Plastiques (15 \%) : sachets PE, bouteilles PET & \textbf{Recyclage mécanique} : tri, lavage, broyage, extrusion (granulés r-PET) ; ou \textbf{pavés plastique-sable} (fusion du PE/PP comme liant) & Granulés pour nouvelles bouteilles, pavés de voirie. \emph{Env. :} moins de sachets dans les caniveaux $\to$ moins d'inondations. \emph{Socio-éco. :} revenus des ramasseurs (« éboueurs informels » formalisés), pavés moins chers. \\ \hline
\textbf{3. Recyclage papier-carton} & Papiers/cartons (10 \%) & \textbf{Repulpage} : hydropulpeur, désencrage, formation de feuille & Papier hygiénique, serviettes, carton d'emballage, papier-cadeau. \emph{Env. :} économise fibres vierges et forêts. \emph{Socio-éco. :} micro-unités artisanales locales, fournitures scolaires à bas coût. \\ \hline
\end{tabularx}
\end{center}

\textbf{Justification globale :} la composition (55 \% d'organiques fermentescibles) impose de prioriser la filière organique ; le plastique, bien que minoritaire en masse, pose le problème environnemental le plus visible (caniveaux, inondations à Yaoundé) ; le papier est le flux le plus simple à valoriser artisanalement. Le verre et les métaux (dans les 20 \% restants) peuvent faire l'objet d'une quatrième filière de revente aux fondeurs/verreries.

\textbf{Points de vigilance :} chaque filière doit cibler un flux \emph{majoritaire ou problématique} et associer systématiquement un avantage environnemental ET un avantage socio-économique distincts.
\end{corrige}

\begin{corrige}[Exercice 9 — Sujet type Bac : Bioéthanol de 2\up{nde} génération à partir de résidus de plantain]
\textbf{Barème indicatif (20 points) :} Q1 : 1 pt ; Q2 : 3 pts ; Q3 : 3 pts ; Q4 : 3 pts ; Q5 : 5 pts ; Q6 : 5 pts.

\textbf{1. Classement des étapes}

L'ordre du Document 2 est correct et logique : Prétraitement $\to$ Hydrolyse enzymatique $\to$ Fermentation $\to$ Distillation/Déshydratation $\to$ Valorisation des co-produits. Logique : on ne peut hydrolyser qu'une biomasse rendue accessible (prétraitement) ; on ne peut fermenter que des sucres libérés (hydrolyse) ; on ne peut distiller qu'un moût fermenté ; la valorisation des résidus clôt le procédé.

\textbf{2. Nécessité du prétraitement}

Dans la biomasse lignocellulosique (épluchures, pseudotiges), la cellulose est enfermée dans une matrice rigide de \textbf{lignine} et d'\textbf{hémicellulose} qui la rend physiquement inaccessible aux enzymes ; sa structure cristalline aggrave cette récalcitrance. Le prétraitement (vapeur explosion ou acide dilué) rompt cette matrice, élimine une partie de la lignine et de l'hémicellulose et augmente la surface accessible. Les trois polymères cibles sont : \textbf{cellulose}, \textbf{hémicellulose}, \textbf{lignine}. Au contraire, le jus de canne contient des sucres \emph{déjà libres} (saccharose soluble) et l'amidon du maïs est facilement hydrolysable par simple liquéfaction/saccharification : aucun prétraitement destructeur n'est nécessaire.

\textbf{3. Hydrolyse enzymatique}

Équation bilan simplifiée :
\[ \ce{(C6H10O5)_n + n H2O ->[cellulases] n C6H12O6} \]
Les cellulases (endoglucanases, exoglucanases) dégradent la cellulose en \textbf{célobiose} (dimère de glucose), qui est un \textbf{inhibiteur} des cellulases elles-mêmes. La $\beta$-glucosidase hydrolyse le célobiose en deux glucose, levant cette inhibition par le produit : sans elle, le célobiose s'accumule et le rendement en glucose s'effondre. Elle est donc cruciale pour le rendement global.

\textbf{4. Fermentation C5/C6}

\emph{S. cerevisiae} sauvage ne possède pas la voie métabolique d'assimilation des pentoses : elle ne dispose ni du transporteur spécifique ni des enzymes initiales (xylose réductase / xylitol déshydrogénase, ou xylose isomérase) pour convertir le xylose en xylulose-5-phosphate entrant dans la voie des pentoses phosphates. Le xylose issu de l'hémicellulose (jusqu'à 30--40 \% des sucres de la biomasse) resterait donc non valorisé. La souche \textbf{recombinante} a reçu les gènes de cette voie : elle co-fermente glucose (C6) et xylose (C5), ce qui augmente fortement le rendement en éthanol par tonne de biomasse.

\textbf{5. Analyse de la courbe (Doc 3)}
\begin{itemize}
    \item \textbf{(a) Cinétique :} la SSF démarre beaucoup plus vite (15 g/L à 12 h contre 5 g/L en SHF) : le glucose est consommé par la levure dès sa libération.
    \item \textbf{(b) Rendement final :} à 72 h, SSF $\approx$ \SI{51}{\gram\per\liter} contre $\approx$ \SI{49}{\gram\per\liter} en SHF : léger avantage SSF.
    \item \textbf{(c) Inhibition enzymatique :} en SHF, l'hydrolyse séparée accumule le glucose (et le célobiose), qui \textbf{inhibe les cellulases} (rétro-inhibition par le produit), ralentissant la saccharification. En SSF, la consommation simultanée du glucose maintient sa concentration basse et lève cette inhibition.
    \item \textbf{Conclusion :} la \textbf{SSF est le procédé le plus performant} pour l'usine de Njombé : plus rapide, titre final supérieur, un seul réacteur (moindre investissement) ; sa contrainte est le compromis de température (\SI{32}{\celsius}--\SI{37}{\celsius}, entre l'optimum enzymatique de \SI{50}{\celsius} et celui de la levure).
\end{itemize}

\textbf{6. Bilan énergétique et environnemental}
\begin{align*}
m_{\text{cellulose+hémicellulose}} &= 1000 \times 0{,}60 = \SI{600}{\kilogram} \\
m_{\text{éthanol théorique}} &= 600 \times \frac{92}{180} \approx \SI{306,7}{\kilogram} \\
m_{\text{éthanol réel}} &= 306{,}7 \times 0{,}80 \approx \SI{245,3}{\kilogram} \\
V_{\text{éthanol}} &= \frac{245{,}3}{0{,}789} \approx \boxed{\SI{311}{\liter} \text{ par tonne d'épluchures}}
\end{align*}
(Le facteur 92/180 correspond à \ce{C6H12O6 -> 2 C2H5OH + 2 CO2} ; on assimile l'hémicellulose à un rendement équivalent, ce qui est une approximation haute.)

\textbf{Valorisation de la lignine :} le résidu solide (lignine, PCI $\approx$ \SI{20}{\mega\joule\per\kilogram}) brûlé en chaudière biomasse fournit la \textbf{vapeur} nécessaire au prétraitement et à la distillation, postes très énergivores : l'usine devient \textbf{auto-suffisante en énergie thermique}, sans fossile, ce qui améliore le bilan \ce{CO2} et la rentabilité. À l'échelle du Cameroun (\num{4,5} millions de tonnes de plantains, 40 \% de résidus), le gisement théorique dépasse \SI{500000}{\tonne} d'éthanol par an si toute la filière était mobilisée.

\textbf{Points de vigilance :} exiger la citation des trois polymères (cellulose, hémicellulose, lignine) ; le calcul doit faire apparaître le facteur stœchiométrique 92/180 et le rendement de 80 \% ; la comparaison SSF/SHF doit mobiliser les valeurs de la courbe, pas des affirmations générales.
\end{corrige}

\begin{corrige}[Exercice 10 — Sujet type Bac : Gestion intégrée des déchets plastiques à Douala]
\textbf{Barème indicatif (20 points) :} Q1 : 4 pts ; Q2 : 5 pts ; Q3 : 4 pts ; Q4 : 3 pts ; Q5 : 4 pts.

\textbf{1. Triage amont}

L'Option A (r-PET contact alimentaire) exige une pureté quasi absolue : le PET est fondu à $\approx$ \SI{250}{\celsius}, température à laquelle le PE/PP fond bien plus tôt et forme des inclusions, et surtout le PVC se \textbf{dégrade} dès $\approx$ \SI{200}{\celsius}. Si 5 \% de PVC contaminent le flux, il se \textbf{décompose en extrusion en libérant de l'acide chlorhydrique \ce{HCl}} (gaz corrosif pour les machines) et des composés chlorés qui jaunissent et fragilisent le PET, détruisent ses propriétés mécaniques et rendent le lot impropre au contact alimentaire. Pour l'Option B (pavés), le procédé est tolérant : PE, PP et PS fondent ensemble entre \SI{110}{\celsius} et \SI{160}{\celsius} et servent tous de liant ; un tri grossier suffit (on écarte seulement PET et PVC).

\textbf{2. Option A : recyclage « bouteille à bouteille » (6 étapes)}
\begin{enumerate}
    \item \textbf{Tri} (optique NIR + manuel) : séparation du PET par couleur, élimination PVC, métaux, corps étrangers.
    \item \textbf{Broyage} en paillettes (flakes) avec séparation des bouchons (PP/PE) et étiquettes par densité (cuve de flottation : PET coule, PE/PP flottent).
    \item \textbf{Lavage} à chaud (soude, \SI{80}{\celsius}) et rinçage pour éliminer colles, résidus alimentaires.
    \item \textbf{Séchage} poussé (le PET est sensible à l'hydrolyse : l'eau résiduelle romprait les chaînes à l'extrusion).
    \item \textbf{Extrusion-filtration} : fusion à $\approx$ \SI{270}{\celsius}, filtration fine ($<$ \SI{50}{\micro\meter}), décontamination sous vide, granulation.
    \item \textbf{Polycondensation à l'état solide (SSP/SPC)} : c'est ici que l'on \textbf{restaure la viscosité intrinsèque (IV $>$ \SI{0,75}{\deci\liter\per\gram})} : les granulés sont chauffés ($\approx$ \SI{200}{\celsius}--\SI{220}{\celsius}) sous vide ou azote, sous la température de fusion ; les réactions de polycondensation se poursuivent en éliminant eau et éthylène glycol, ce qui rallonge les chaînes et élimine les contaminants volatils (décontamination $\to$ certification contact alimentaire).
\end{enumerate}

\textbf{3. Option B : pavés plastique-sable}

Procédé : tri/lavage/broyage du PE-PP-PS ; fusion-malaxage à \SI{160}{\celsius}--\SI{180}{\celsius} ; incorporation du sable chaud au dosage typique \textbf{30 \% plastique / 70 \% sable} en masse ; moulage à chaud sous presse ; refroidissement et démoulage. Avantages techniques en climat tropical humide (Douala) :
\begin{itemize}
    \item \textbf{Absorption d'eau quasi nulle ($<$ 0,5 \%)} : pas de fissuration par humidité, pas de développement de mousses, résistance aux pluies diluviennes, contrairement au béton poreux (4--6 \%).
    \item \textbf{Insensibilité à la carbonatation et à la corrosion} : pas d'armatures ni de liant hydraulique à hydrater ; le pavé ne s'effrite pas en ambiance saline/humide du littoral et supporte mieux les cycles mouillage-séchage.
\end{itemize}

\textbf{4. Option C : pyrolyse}

Équation bilan globale simplifiée (craquage thermique du polyéthylène en l'absence d'oxygène) :
\[ \ce{-(CH2-CH2)_{n}- ->[\text{\SIrange{400}{500}{\celsius}, catalyseur}] x\ C_{a}H_{2a+2} + y\ C_{b}H_{2b} + H2 + CH4 + char} \]
soit, schématiquement : $\ce{(-CH2-CH2-)_{n}} \to$ mélange d'alcanes $\ce{C_{n}H_{2n+2}}$ et d'alcènes $\ce{C_{n}H_{2n}}$ de masses molaires variées (fraction diesel/gazole).

\textbf{Avantage majeur :} elle valorise les flux \textbf{sales, mélangés et non recyclables mécaniquement} (sachets souillés, multicouches) que les options A et B refusent, avec $\approx$ 70 \% de liquide valorisable en carburant. \textbf{Risque majeur :} risque d'émissions toxiques (dioxines, furanes, surtout si du PVC passe au four) et de rejets de char contaminé si le contrôle des fumées est insuffisant ; économiquement, la filière est sensible au prix du pétrole et exige un contrôle technique pointu — et elle détruit la matière (perte de circularité matière).

\textbf{5. Choix stratégique (recommandation d'expert)}

Je recommande une stratégie combinée hiérarchisée : \textbf{40 \% Option B (pavés), 30 \% Option A (r-PET), 20 \% Option C (pyrolyse), 10 \% refus/enfouissement résiduel}. Justification : le gisement doualien est dominé par les sachets PE-LD et plastiques mixtes souillés, peu aptes au r-PET mais parfaits pour les pavés, filière à faible investissement, techniquement mature et très créatrice d'emplois locaux (collecte, tri, fabrication, pose) ; les pavés répondent en outre à un besoin communal direct (voiries, cours d'écoles) et substituent un produit importé en ciment. L'Option A, plus capitalistique et exigeante (tri fin, SSP, certification alimentaire), ne doit traiter que le flux de bouteilles PET propres et triées, en partenariat avec les brasseries et embouteilleurs locaux. L'Option C, en complément, absorbe les refus de tri et plastiques sales non recyclables, transformant une nuisance en carburant, mais reste minoritaire car elle sacrifie la matière au profit de l'énergie et présente les plus gros risques environnementaux. Cette combinaison respecte la hiérarchie des modes de traitement (matière d'abord, énergie ensuite), maximise la circularité, sécurise les débouchés et répartit les risques techniques et financiers.

\textbf{Points de vigilance :} Q1 : la contamination PVC $\to$ \ce{HCl} est la justification-clé attendue ; Q2 : l'étape de restauration de l'IV (SSP) et son principe doivent être nommés explicitement ; Q5 : toute stratégie cohérente est acceptée si elle s'appuie sur le gisement, la maturité technique, l'emploi et la hiérarchie matière/énergie.
\end{corrige}

\begin{corrige}[Exercice 11 — Sujet type Bac : Sensibilisation au tri dans les marchés]
\textbf{Barème indicatif (10 points) :} Q1 : 2 pts ; Q2 : 5 pts ; Q3 : 3 pts.

\textbf{1. Choix du support}

Je choisis le \textbf{spot radio de 45 secondes} : au Marché Central ou à Sandaga, la radio (y compris les radios communautaires et les téléphones portables) touche tout le monde, y compris les commerçants et porteurs peu scolarisés ; elle permet le multilinguisme (français, pidgin, ewondo, douala) et l'usage de voix, de musique et d'humour, plus percutants qu'un texte écrit pour cette cible.

\textbf{2. Texte intégral du spot radio (45 s)}

\begin{quote}
\textbf{Spot radio « Deux poubelles, zéro inondation »}

\emph{(Bruitages de marché, tam-tam, voix jeune et dynamique)}

— « Hé papa, hé maman ! Au marché, le sachet plastique que tu jettes dans le caniveau, c'est lui qui bloque l'eau quand la pluie tombe... et c'est ta boutique qui inonde !

— Mais écoute bien : le plastique, ce n'est pas un déchet, c'est de l'argent ! Les jeunes du quartier le ramassent, le recyclent en pavés et en nouvelles bouteilles.

— Alors fais le geste simple : \textbf{deux poubelles au marché !} Dans la poubelle \textbf{jaune} : sachets, bouteilles, plastiques. Dans la poubelle \textbf{verte} : restes de nourriture, feuilles, carton sale.

— Retiens la règle des 3R : \textbf{Réduis} tes sachets, \textbf{Réutilise} ton cabas, \textbf{Recycle} le reste !

— Deux poubelles, zéro inondation, plus d'argent pour la jeunesse ! Marché propre, quartier gagnant !

— Message de la mairie et du comité d'hygiène du marché. Info : point de collecte des plastiques à l'entrée principale, tous les samedis. »
\end{quote}

\textbf{3. Argumentaire scientifique (annexe pour les relais communautaires)}

Le sachet « noir » est en polyéthylène basse densité (\text{PE-LD}), un polymère de carbone et d'hydrogène issu du pétrole : le jeter en décharge ou le brûler à l'air libre, c'est gaspiller une \textbf{matière fossile} qui a coûté de l'énergie à produire. La combustion à basse température et mal oxygénée (feux sauvages) est incomplète : elle émet des \textbf{dioxines et furanes} (cancérogènes persistants), des particules fines et des hydrocarbures aromatiques polycycliques. En décharge, le sachet se fragmente en \textbf{microplastiques} qui contaminent sols, nappes phréatiques (via les lixiviats) et chaîne alimentaire, et il obstrue les caniveaux, favorisant inondations et gîtes larvaires de moustiques (paludisme). Or le \text{PE-LD} est un thermoplastique parfaitement \textbf{recyclable} : fondu vers \SI{110}{\celsius}--\SI{130}{\celsius}, il sert de liant dans les pavés plastique-sable ou est régénéré en granulés, substituant des matériaux neufs. Recycler, c'est donc transformer une nuisance sanitaire et environnementale en ressource économique locale.

\textbf{Points de vigilance :} le message doit contenir les 4 objectifs (pourquoi, comment, valorisation, 3R) de façon identifiable ; l'argumentaire scientifique doit citer au moins deux polluants nommés (dioxines/furanes, microplastiques) et la notion de perte de matière carbonée fossile.
\end{corrige}

\newpage

\coursec{Fiche bilan}

\courssub{Carte mentale de synthèse}

\begin{popfigure}
\begin{tikzpicture}[
    mindmap,
    grow cyclic,
    every node/.style={concept, execute at begin node=\hskip0pt},
    concept color=popblue,
    level 1/.append style={level distance=4.5cm, sibling angle=60, font=\bfseries\small},
    level 2/.append style={level distance=3cm, sibling angle=45, font=\footnotesize},
    biocarburant/.style={concept color=popgreen},
    papier/.style={concept color=poporange},
    recyclagepapier/.style={concept color=poppurple},
    recyclageplastique/.style={concept color=poppink},
    paves/.style={concept color=popgold},
    avantages/.style={concept color=popdark}
]
\node[root concept, text width=3.5cm, align=center, font=\bfseries\normalsize] {Énergies renouvelables \& Valorisation des déchets}
    child[biocarburant] { node[concept, text width=2.8cm, align=center] {Biocarburants\\ (Bioéthanol, Biodiesel, Biogaz)}
        child { node[concept, text width=2.2cm, align=center] {1ère gén : amidon/sucre (maïs, canne, manioc)} }
        child { node[concept, text width=2.2cm, align=center] {2ème gén : lignocellulose (résidus agricoles)} }
        child { node[concept, text width=2.2cm, align=center] {3ème gén : micro-algues} }
        child { node[concept, text width=2.2cm, align=center] {Méthanisation : déchets organiques $\to$ CH$_4$} }
    }
    child[papier] { node[concept, text width=2.5cm, align=center] {Énergie depuis le papier}
        child { node[concept, text width=2cm, align=center] {Combustion directe (briquettes)} }
        child { node[concept, text width=2cm, align=center] {Pyrolyse / Gazéification} }
    }
    child[recyclagepapier] { node[concept, text width=2.8cm, align=center] {Recyclage papier}
        child { node[concept, text width=2cm, align=center] {Serviettes jetables} }
        child { node[concept, text width=2cm, align=center] {Papier hygiénique} }
        child { node[concept, text width=2cm, align=center] {Papier-cadeau / carton} }
    }
    child[recyclageplastique] { node[concept, text width=2.8cm, align=center] {Recyclage plastique (bouteilles)}
        child { node[concept, text width=2cm, align=center] {Tri : PET, PEHD, PP} }
        child { node[concept, text width=2cm, align=center] {Broyage $\to$ Lavage $\to$ Flocons} }
        child { node[concept, text width=2cm, align=center] {Extrusion $\to$ Granulats} }
        child { node[concept, text width=2cm, align=center] {Soufflage : nouvelles bouteilles} }
    }
    child[paves] { node[concept, text width=2.8cm, align=center] {Plastique $\to$ Pavés}
        child { node[concept, text width=2cm, align=center] {Fusion mélange plastiques/sable} }
        child { node[concept, text width=2cm, align=center] {Moulage sous pression} }
        child { node[concept, text width=2cm, align=center] {Pavés écologiques, durables} }
    }
    child[avantages] { node[concept, text width=2.5cm, align=center] {Avantages \& Sensibilisation}
        child { node[concept, text width=2cm, align=center] {Réduction pollution (sol, eau, air)} }
        child { node[concept, text width=2cm, align=center] {Économie circulaire, emplois verts} }
        child { node[concept, text width=2cm, align=center] {Substitution énergies fossiles} }
        child { node[concept, text width=2cm, align=center] {Éducation citoyenne (tri à la source)} }
    };
\end{tikzpicture}
\legende{Carte mentale récapitulant les filières de valorisation énergétique et matière des déchets (contexte camerounais : manioc, déchets urbains Yaoundé/Douala).}
\end{popfigure}

\courssub{Bilan des savoirs et savoir-faire}

\begin{bilan}

\textbf{1. Déclinations clés}
\begin{itemize}
    \item \cle{Énergie renouvelable} : Énergie issue de sources que la nature renouvelle assez vite pour les considérer inépuisables à l'échelle humaine (soleil, vent, eau, biomasse, géothermie).
    \item \cle{Biocarburant} : Carburant produit à partir de biomasse (végétaux, déchets organiques). Principaux : \textbf{bioéthanol} (essence), \textbf{biodiesel} (gazole), \textbf{biogaz} (méthane).
    \item \cle{Valorisation des déchets} : Ensemble des opérations visant à donner une valeur (énergétique ou matière) à un déchet : recyclage, compostage, méthanisation, incinération avec récupération d'énergie, transformation en nouveaux matériaux.
    \item \cle{Économie circulaire} : Modèle de production/consommation : réduire, réutiliser, recycler, pour boucler le cycle de vie des produits.
\end{itemize}

\textbf{2. Filières de production des biocarburants (Savoir-faire : Expliquer les techniques)}

\begin{center}
\begin{tabularx}{\linewidth}{|l|X|X|}\hline
\rowcolor{popblueL}
\textbf{Filière} & \textbf{Matières premières (Ex. Cameroun)} & \textbf{Étapes clés / Réactions} \\ \hline
\textbf{Bioéthanol (1ère gén)} & Canne à sucre, maïs, manioc, patate douce, sorgho & 1. Broyage/extraction sucres/amidon. 2. \textbf{Hydrolyse} amidon $\xrightarrow{\text{amylases}}$ glucose. 3. \textbf{Fermentation alcoolique} : \ce{C6H12O6 ->[levures] 2C2H5OH + 2CO2}. 4. Distillation (éthanol $\approx$ 96\,\%), déshydratation (éthanol anhydre $>99,5\,\%$). \\ \hline
\textbf{Biodiesel (1ère gén)} & Huiles : palmiste, coton, jatropha, ricin, usagées & \textbf{Transestérification} : Huile (triglycérides) + Méthanol $\xrightarrow{\text{catalyseur (NaOH/KOH)}}$ Esters méthyliques (biodiesel) + Glycérine. \\ \hline
\textbf{Biogaz (Méthanisation)} & Déchets agro-industriels (bouses, lisiers, résidus brasseries, marchés), boues stations d'épuration & \textbf{Digestion anaérobie} (4 phases) : Hydrolyse $\to$ Acidogenèse $\to$ Acétogenèse $\to$ \textbf{Méthanogenèse} (archées méthanogènes) : \ce{CH3COOH -> CH4 + CO2} / \ce{CO2 + 4H2 -> CH4 + 2H2O}. Biogaz $\approx$ 50-70\,\% CH$_4$. Épurage $\to$ Biométhane. \\ \hline
\textbf{2ème / 3ème génération} & Résidus lignocellulosiques (bagasse, paille, sciure), micro-algues & Pré-traitement (vapeur, acide) $\to$ Hydrolyse enzymatique (cellulases) $\to$ Fermentation (bioéthanol) OU Pyrolyse/Gazéification $\to$ Syngaz $\to$ Carburants Fischer-Tropsch. \\ \hline
\end{tabularx}
\end{center}

\textbf{3. Valorisation énergétique du papier}
\begin{itemize}
    \item \textbf{Combustion directe} : Papiers/cartons non recyclables (souillés, composites) $\to$ Briquettes / granulés $\to$ Chaudières industrielles, fours à ciment (substitution charbon).
    \item \textbf{Pyrolyse / Gazéification} : Chauffage sans oxygène (400-800\,°C) $\to$ Huile de pyrolyse, gaz de synthèse (H$_2$, CO), biochar.
\end{itemize}

\textbf{4. Recyclage du papier : Serviettes, Hygiénique, Cadeau (Savoir-faire : Expliquer la technique)}
\begin{methode}[Procédé industriel de recyclage papier (désencrage)]
\begin{enumerate}
    \item \textbf{Tri \& Broyage} : Déchiquetage + eau $\to$ \textbf{Pâte à papier} (pulper / hydropulpeur). Séparation grosses impuretés (agrafes, plastiques) par criblage / purification centrifuge.
    \item \textbf{Désencrage (Flottation)} : Injection bulles d'air + tensioactifs $\to$ Particules d'encre hydrophobicques fixées sur bulles $\to$ Mousse en surface $\to$ Évacuation. (Répété 2-3 fois).
    \item \textbf{Blanchiment (si blanc requis)} : Peroxyde d'hydrogène (H$_2$O$_2$) ou Hydrosulfite de sodium (Na$_2$S$_2$O$_4$) $\to$ Éclaircissement fibres sans chlore (TCF).
    \item \textbf{Affinage \& Nettoyage fin} : Affineur (défibrage final), nettoyeurs centripètes (sable, colle).
    \item \textbf{Fabrication produit final} :
    \begin{itemize}
        \item \textbf{Serviettes / Hygiénique} : Pâte \emph{tissue} (grammage faible 15-30\,g/m$^2$), crêpage (cylindre Yankee chauffé $\approx$ 160\,°C) $\to$ Douceur, absorption. Bobonnage, découpe, emballage.
        \item \textbf{Papier-cadeau / Carton} : Pâte plus épaisse, calandrage (lissage), enduction éventuelle (brillance), impression, découpe.
    \end{itemize}
\end{enumerate}
\end{methode}

\textbf{5. Recyclage des déchets plastiques en bouteilles recyclées (r-PET)}
\begin{methode}[Filière bouteille $\to$ bouteille (Boucle fermée PET)]
\begin{enumerate}
    \item \textbf{Collecte \& Tri} : Tri manuel/optique (NIR) : séparation PET (code 1), PEHD (code 2), PP (code 5), PVC (interdit). Élimination bouchons (PEHD/PP), étiquettes.
    \item \textbf{Broyage} : Bouteilles $\to$ \textbf{Flocons} (flakes) 8-12\,mm.
    \item \textbf{Lavage intensif} : Eau chaude (80-90\,°C) + détergent + soude $\to$ Élimination colle, étiquettes, contaminants, résidus boisson. Rinçage, séchage (humidité $<$ 0,5\,\%).
    \item \textbf{Extrusion \& Granulation} : Flocons fondus (260-280\,°C) $\to$ Filtration (maille 40-80\,µm) $\to$ Granulats (pellets) r-PET. \textbf{SSP (Solid State Polycondensation)} : Rehausse de la viscosité intrinsèque (IV) pour tenue mécanique/barrière.
    \item \textbf{Injection-Soufflage (ISBM)} : Préfomes injectées $\to$ Réchauffage $\to$ Soufflage dans moule $\to$ Nouvelles bouteilles (incorporation 25-100\,\% r-PET selon usage alimentaire).
\end{enumerate}
\end{methode}

\textbf{6. Transformation des déchets plastiques en pavés (Savoir-faire : Expliquer la technique)}
\begin{methode}[Pavés plastiques (Valorisation matière « bas de gamme » / mélange)]
\begin{enumerate}
    \item \textbf{Tri \& Prép.} : Plastiques souples (sachets PEBD, PP, PS) souvent sales, non triables finement. Lavage sommaire, séchage, broyage grossier.
    \item \textbf{Mélange} : Plastiques fondus (liant) + \textbf{Sable} (charge minérale, $\approx$ 60-70\,\% masse) $\pm$ colorants. Ratio typique 1 vol. plastique / 3 vol. sable.
    \item \textbf{Extrusion / Moulage} : Four rotatif ou extrudeuse bi-vis $\to$ Pâte homogène $\to$ Injection/Moulage sous pression dans moules métalliques (formats pavés autobloquants, dalles).
    \item \textbf{Refroidissement \& Démoulage} : Refroidissement air/eau $\to$ Pavés solides, imputrescibles.
\end{enumerate}
\end{methode}

\textbf{7. Avantages de la valorisation des déchets (Savoir-faire : Informer la communauté)}
\begin{center}
\begin{tabularx}{\linewidth}{|l|X|}\hline
\rowcolor{popblueL}
\textbf{Dimension} & \textbf{Avantages majeurs} \\ \hline
\textbf{Environnementale} & Réduction volume décharges (ex. Nkolfoulou, PK10 à Yaoundé ; Dibamba à Douala) ; Lutte pollution sols/eaux (lixiviats) \& air (brûlage à l'air libre = dioxines, furanes) ; Préservation biodiversité forêts (moins pâte vierge) ; Réduction GES (substitution fossile, évitement CH$_4$ décharge). \\ \hline
\textbf{Économique} & Création emplois (collecte, tri, usines recyclage, PME pavés) ; Matières premières secondaires moins chères (r-PET, pâte recyclée) ; Indépendance énergétique (biogaz, bioéthanol local) ; Revenus pour trieurs informels (organisation coopératives). \\ \hline
\textbf{Sociale / Citoyenne} & Amélioration cadre de vie (rues propres, caniveaux débouchés $\to$ moins inondations/paludisme) ; Éducation tri à la source (bacs sélectifs) ; Valorisation savoir-faire local (artisans pavés). \\ \hline
\end{tabularx}
\end{center}

\end{bilan}

\courssub{Checklist avant le Bac}

\begin{aretenir}[Checklist Bac - Leçon 11 : Énergies renouvelables \& Valorisation déchets]
    $\square$ Je définis \textbf{énergie renouvelable}, \textbf{biocarburant}, \textbf{valorisation}, \textbf{économie circulaire}.
    $\square$ Je décris les \textbf{3 générations de biocarburants} et donne un exemple de matière première camerounaise pour chacune (manioc, bagasse, algues).
    $\square$ J'écris l'équation bilan de la \textbf{fermentation alcoolique} ($\ce{C6H12O6 -> 2C2H5OH + 2CO2}$) et le principe de la \textbf{transestérification} (biodiesel).
    $\square$ J'explique les \textbf{4 étapes de la méthanisation} (hydrolyse, acidogenèse, acétogenèse, méthanogenèse) et la composition du biogaz.
    $\square$ Je détaille le \textbf{procédé de désencrage par flottation} pour le recyclage papier (rôle bulles + tensioactif).
    $\square$ Je distingue les produits finis \textbf{papier tissue} (serviettes, hygiénique : crêpage Yankee) et \textbf{papier-cadeau/carton} (calandrage).
    $\square$ Je cite dans l'ordre les étapes du recyclage \textbf{bouteille PET $\to$ r-PET $\to$ nouvelle bouteille} (Tri, Broyage, Lavage, Extrusion/SSP, Soufflage).
    $\square$ J'explique la fabrication des \textbf{pavés plastiques} (mélange plastiques fondus + sable, moulage sous pression).
    $\square$ Je cite \textbf{3 avantages environnementaux}, \textbf{2 économiques} et \textbf{1 social} de la valorisation des déchets.
    $\square$ Je sais argumenter pour \textbf{sensibiliser ma communauté} : tri à la source (bacs distincts), interdiction brûlage déchets, soutien filières locales recyclage.
    $\square$ Je relie cette leçon aux enjeux \textbf{ODD 7 (Énergie propre), ODD 11 (Villes durables), ODD 12 (Consommation responsable), ODD 13 (Climat)}.
\end{aretenir}

\findecours

\end{document}
